%!TEX root = /Users/kbj/Library/CloudStorage/Dropbox/TransformationalGrammar/601master.tex
\setcounter{chapter}{2}
\chapter{Positioning Arguments}
\label{ch:positargs}
\setcounter{exx}{0}

\noindent
Let's turn now to the question of what fixes the order of complements. We'll focus here on the question of what determines the linear position of the internal arguments to verbs. Recall that the way in which internal arguments to verbs are ordered does not vary with the choice of verb. We will not want to encode the order of arguments on a verb by verb basis, then. Instead, the relevant factor in determining the order of objects seems to be the categorial type that the objects are. If we focus on DP, PP and CP, the pattern our initial phrase structure rules encoded as follows.
\begin{exe}
	\ex \label{ex3:comptemplate} DP PP\\ DP CP\\ PP PP\\ PP CP
	\ex \label{ex3:examcomptempl} \marginnote[0pt]{As a simplifying measure, we consider only finite CPs, that is: CPs whose IP contains tense inflection or a modal.}
	\begin{xlist}
	\ex[]{I told [\DP her] [\CP that she should join].}
	\ex[*]{I told [\CP that she should join] [\DP her].}
	\ex[]{Amy Rose explained [\PP to me] [\CP that I should join].}
	\ex[?*]{Amy Rose explained [\PP that I should join] [\PP to me].}
	\ex[]{Taka told [\DP it] [\PP to Della].}
	\ex[??]{Taka told [\PP to Della] [\DP it].}
	\end{xlist}
\end{exe}
I have not considered cases where a verb's internal arguments are both CPs. These cases are either very rare or non-existent. I've also not considered cases where a verb's internal arguments are both DPs. That scenario does arise, but it has some rather special properties. We will come back to consider those cases later.

In his dissertation,\nocite{Stowell:1981} Tim Stowell argued for an answer to this question that has been influential. He suggested that systems of ``Case assignment'' are responsible for ordering these terms. DPs differ from other kinds of phrases in Indo-European in being able to host Case morphology. In English this happens with pronouns only: the pronouns \emph{she} and \emph{her}, for instance, have the same meaning but have a different ``Case'' form. The particular Case borne by a pronoun is determined by its syntactic position. In languages that are richer in Cases than English is, we can see that the Case borne by a DP is determined by a term proximate to the DP. In German, for instance, a certain class of prepositions and verbs determine Accusative Case for their complements, while others determine Dative Case. It is also often the case in other Indo-European languages that DPs other than just pronouns can bear Case morphology---sometimes on the head of the DP, sometimes on the determiner, sometimes spread across modifiers of various sorts within the DPs. Let's imagine then that, in general, DPs must be related to Case assigners. We'll call this the Case Filter.
\begin{exe}
	\ex \label{ex3:.21} \textsc{Case Filter}\\ 
		A DP must be assigned Case if it is an argument.
\end{exe}
I've restricted the Case Filter to argument DPs because modifying DPs do not seem to be positioned in a way that suggests they are sensitive to Case assignment. In (\ref{ex3:Tuesday}), for instance, the DP \emph{Tuesday} is not in a position that an argument DP, like the pronoun \emph{it}, can be.
\begin{exe}
	\raggedright
	\ex \label{ex3:Tuesday} 
	\small
	\begin{forest}
	[\ibar{I}
		[\obar{I} [will] ]
		[VP [\ibar{V} 
			[\ibar{V}
				[\obar{V} [talk] ]
				[PP [to the waiter, roof]]
			 ]
			[DP [Tuesday, roof]]
		] ]
	]
	\end{forest}
	\ex
	\begin{forest}
	[IP
		[DP [\ibar{D} [\obar{D} [she]]]]
		[\ibar{I} 
			[\obar{I} [will] ]
			[VP [\ibar{V} 
				[\ibar{V}
					[\obar{V} [tell] ]
					[PP [to the waiter, roof] ]
				 ]
				[DP [\ibar{D} [\obar{D} [it]]] ]
			] ]
		 ]
	]
	\end{forest}
\end{exe}
Moreover, typically the Case morphology that modifying DPs bear in languages where Case morphology is found on them is fixed, and not sensitive to nearby Case assigners. The Case filter, then, requires that some parse in a sentence's derivation puts every argument DP which that sentence contains in a Case marked position.

As I noted before, Case marked positions are ones that are close to terms that are responsible for assigning (i.e, determining) the Case. So, all we have to do now is know what those terms are and what ``close'' means, and we'll be able to use the Case filter to rule out certain strings as sentences. The ``object'' Cases---so-called Accusative, Dative, Locative, and a host of others---are assigned by particular lexical items. In English there are just three cases: accusative, nominative and genitive. The form of pronouns in English depends on this Case information, as well as the number, person, and, in some cases, gender of the referent of the pronoun. 
\begin{exe}
\sn \label{ex3:casetable} 
\begin{tabular}{l|ccc}
  & \emph{Nominative} & \emph{Accusative} & \emph{Genitive}\\ \hline
  \nth{1}, sing & I & me & my\\
  \nth{1}, plur & we & us & our\\
  \nth{2}, sing & you & you & your\\
  \nth{2}, plur & you & you & your\\
  \nth{3}, sing, fem & she & her & her\\
  \nth{3}, sing, masc & he & him & his\\
  \nth{3}, sing, neut & it & it & its\\
  \nth{3}, plur & they & them & their\\
\end{tabular}
\end{exe}
Our immediate task, then, is to determine what terms are responsible for assigning these Cases. The genitive appears on DPs that are in the Specifier of DP in English, as in (\ref{ex3:genex}).
\begin{exe}
\ex \label{ex3:genex} 
\begin{tabular}{llll}
a. & \textbf{your} hat & \emph{compare:} & *\textbf{your} left.\\
b. & \textbf{his} marbles & \emph{compare:} & *I talked to \textbf{his}.\\
c. & \textbf{our} problems & \emph{compare:} & *She met \textbf{our}.\\
\end{tabular}
\end{exe}
We put off for a couple chapters investigating what the item is that assigns Genitive; there are facts about the shape of DPs that we have yet to discover that are relevant.

Nominative Case is found on the ``subjects'' of finite IPs, as in (\ref{ex3:nomex}).
\begin{exe}
\ex \label{ex3:nomex} 
\begin{tabular}{llll}
a. & \textbf{She} left. & \emph{compare:} & *I left \textbf{she}.\\
b. & \textbf{We} arrived before \textbf{he} left. & \emph{compare:} & *I talked to \textbf{he}.\\
c. & \textbf{They} thought \textbf{I} would stay. & \emph{compare:} & *\textbf{they} hat
\end{tabular}
\end{exe}
Finite IPs are distinguished from other phrases by the presence of a finite \obar{I}: tense/agree inflection or a modal. Moreover, the nominative Case-marked DP appears in Specifier of IP, at least in the examples we have examined so far. (We will encounter soon cases that don't meet this description.) An initial stab at the conditions on nominative Case assignment, then, is~(\ref{ex3:nomass1}).
\begin{exe}
\ex \label{ex3:nomass1} Finite \obar{I} assigns Nominative Case to its Specifier position.
\end{exe}

Accusative Case is found on DPs that immediately follow some verbs and prepositions, as in~(\ref{ex3:accass1}).
\begin{exe}
\ex \label{ex3:accass1} 
\begin{tabular}{llll}
a. & We met \textbf{him}. & \emph{compare:} & *\textbf{Him} left.\\
b. & They talked to \textbf{me} & \emph{compare:} & *They are happy \textbf{me}.\\
c. & She stood near \textbf{us}. & \emph{compare:} & *\textbf{us} hats.
\end{tabular}
\end{exe}
It appears that verbs and prepositions are the terms that assign accusative Case. Tim Stowell argued that one of the conditions on accusative Case assignment is adjacency. We will see soon that additional conditions are required too. An initial stab at the conditions on accusative Case assignment, then, is~(\ref{ex3:accassc}).
\begin{exe}
\ex \label{ex3:accassc} Accusative Case is assigned to the position adjacent to a verb or preposition.
\end{exe}

It perhaps seems sensible to say that pronouns are subject to the Case filter. Pronouns simply have no pronunciation that isn't sensitive to the syntactic position they are in. To say a \nth{3} person singular female pronoun, for instance, requires choosing the nominative form (\emph{she}) or the accusative form (\emph{her}) or the genitive form (again: \emph{her}), and those forms depend entirely on whether that pronoun is in one of the positions assigned Case. But all other DPs can be pronounced in English without morphology that depends on where that DP is in a sentence. So the evidence that all (argument) DPs are subject to the Case filter is less obvious. Let me sketch here an argument for taking all DPs to be subject to the Case filter, even those that do not have a morphology that signals that they are Case marked.

The argument is based on a fact about adjectives in English. Adjectives in English can have internal arguments, just like verbs can, but they are never DPs. We have contrasts like~(\ref{ex3:approveof}).
\begin{exe}
	\raggedright
	\ex \label{ex3:approveof} 
	\begin{xlist}
		\ex[*]{She might be fond this sentence.}
		\ex[]{She might be fond of this sentence.}
	\end{xlist}
\end{exe}
This was recorded in our phrase structure rules. No rule for APs put a DP inside an \ibar{A}. If we abandon putting information of this kind in phrase structure rules, then there must be some other source of the fact that DPs are never internal arguments of adjectives.

One possibility is that adjectives, for whatever reason, simply never c-select an internal argument that is a DP. Perhaps there is some force that prevents adjectives from having that kind of argument structure. I want to show you this isn't correct. We have evidence that adjectives \uline{can} c-select DPs. It's only that those DPs cannot appear inside the \ibar{A} they project. To see this, we first have to learn something about how some adjectives can express their arguments. Consider the adjective \emph{obvious}. It c-selects a CP argument, forming sentences like~(\ref{ex3:obvious1}).
\begin{exe}
	\raggedright
	\ex \label{ex3:obvious1} 
	It is obvious [\CP that pigs fly].
\end{exe}
Notice that \emph{obvious} has no external \mbox{$\theta$-role}. The Specifier of IP in (\ref{ex3:obvious1}) is occupied by an expletive pronoun. Interestingly, the CP argument of \emph{obvious} can also be found in Specifier of IP, as in~(\ref{ex3:obvious2}).
\begin{exe}
	\raggedright
	\ex \label{ex3:obvious2} 
	That pigs fly is obvious.
\end{exe}
We will see an idea about how to account for this flexibility of expression of the argument of \emph{obvious}, but for now we can simply observe that there are two positions that the argument of \emph{obvious} can be positioned. Unsurprisingly, \emph{obvious}, like all adjectives, cannot take a DP argument that shows up in its \ibar{A}. But it can take a DP argument that shows up in Specifier of IP. This is indicated by the contrast in~(\ref{ex3:obvious3}).\marginnote[0pt]{There is a way of pronouncing (\ref{ex3:obvious3}b) in which it is grammatical. On that pronunciation, \emph{it} is the argument of \emph{obvious}, and \emph{this fact} is a kind of after-thought that spells out what \emph{it} refers to. On that way of construing (\ref{ex3:obvious3}), \emph{this fact} isn't an argument of \emph{obvious} and \emph{it} is. Tha alternative way of getting (\ref{ex3:obvious3}) is not available for the sentence \emph{it is obvious it}, because \emph{it} is not the kind of DP that can express an identifying after-thought.}
\begin{exe}
	\raggedright
	\ex \label{ex3:obvious3} 
	\begin{xlist}
		\ex[]{This fact is obvious.}
		\ex[*]{It is obvious this fact.}
	\end{xlist}
\end{exe}
What we learn from (\ref{ex3:obvious3}) is that the DP argument of \emph{obvious} exists, but must be found in one of the positions that we have determined assigns Case: Specifier of finite IP. This would make sense if the thing that makes adjectives different from verbs and prepositions is their ability to assign Case, and if all DPs are subject to the Case filter.

If it is correct that all argument DPs are subject to the Case filter, then we have two forces that determine where argument DPs can be. They must be c-selected by some head, and they must be able to be in one of the positions that is assigned Case. Only when both of these requirements can be met, will a grammatical sentence have an argument~DP.

\citet{Stowell:1981} proposes to derive the ordering of complements using precisely this logic. If accusative Case is only assigned to positions adjacent to a verb or preposition, then of the many places a c-selected argument DP can be, only those that are adjacent to a verb or preposition will do. Because argument DPs are subject to the Case filter, when they get a \mbox{$\theta$-role} from a verb, they will necessarily be positioned adjacent to this verb. By contrast, PPs are not subject to the Case filter and they are free to appear anywhere in the VP that contains their \mbox{$\theta$-role} assigner. This will ensure that when a DP and PP share a VP, the DP will come closer to the verb than does the PP. In a language like English, in which the linearization parameters cause the verb to be first in its VP, this will mean that an object DP will precede an object PP. 

The Case filter, and the adjacency condition on Accusative Case assignment, together will ensure that DP internal arguments end up positioned before other internal arguments in a \ibar{V}. Let's next consider some of the other cases where a verb has two internal arguments. One of these is when a verb c-selects two PPs as arguments. So far as I have been able to tell, there is no ordering on these PPs. Either order is equally grammatical. There are few verbs with this argument structure, so it is not easy to investigate the generality of this case. One such verb is \emph{talk}, and (\ref{ex3:talk}) indicates that either order is grammatical.
\begin{exe}
	\raggedright
	\ex \label{ex3:talk} 
	\begin{xlist}
		\ex Anna talked about Angela to Matt.
		\ex Anna talked to Angela about Matt.
	\end{xlist}
\end{exe}

Finally, let's consider what happens when one of the two internal arguments is a finite CP. In these cases, the CP comes last.
\begin{exe}
	\raggedright
	\ex \label{ex3:cp1} 
	\begin{xlist}
		\ex[]{Meredith explained to her that pigs are excellent fliers.}
		\ex[??]{Meredith explained that pigs are excellent fliers to her.}
	\end{xlist}
	\ex \label{ex3:cp2}
	\begin{xlist}
		\ex[]{Meredith told her that pigs are excellent fliers.}
		\ex[*]{Meredith told that pigs are excellent fliers her.}
	\end{xlist}
\end{exe}
The contrast in (\ref{ex3:cp1}) is much weaker than the one in (\ref{ex3:cp2}). It might appear that this is because (\ref{ex3:cp2}b) is a Case filter violation, since the argument DP \emph{fliers} is not adjacent to the verb. But there is more going on in this example than meets the eye, and it will take us some time before we're ready to face it. We can simply note at this point that there must be something additional going on in (\ref{ex3:cp2}b) that is responsible for the fact that it is worse than~(\ref{ex3:cp1}b).

There is some evidence that what is relevant for the contrast in (\ref{ex3:cp1}) has to do with factors that are not those we've seen relevant for syntax. If, for instance, the relative lengths of the PP and CP are manipulated, the contrast in word-order changes. Many speakers report finding no, or little, contrast in~(\ref{ex3:cp3}).
\begin{exe}
	\raggedright
	\ex \label{ex3:cp3} 
	\begin{xlist}
		\ex Meredith explained to the tourists crowded along the guardrail that pigs fly.
		\ex Meredith explained that pigs fly to the tourists crowded along the guardrail.
	\end{xlist}
\end{exe}
One way of thinking about this is that the syntax allows either order but one is preferred over the other when it reduces an online processing load. On this view, the weak badness of (\ref{ex3:cp1}b) arises because the online processor is tasked with holding in its buffer all of the words that make up the smallest \ibar{V} before being able to properly place the PP that follows. If while transitioning through this \ibar{V}, the processor's buffer size gets tasked too heavily, or if the guesses the parser makes online lead it to decide that it has finished parsing that \ibar{V}, there will be a boggle upon encountered the PP at the end of the \ibar{V}. There is evidence that all of these ingredients play a role in the strength of the perceived goodness of these sentences. This, then, might be one place where speakers' reports about the goodness or badness of a sentence should not be modeled with a syntax. 

Interestingly, these same factors that we might credit to the parser's behavior show up in the alternations in~(\ref{ex3:cp4}).
\enlargethispage*{20pt}
\begin{exe}
	\raggedright
	\ex \label{ex3:cp4} 
	\begin{xlist}
		\ex \begin{xlist}
			\ex[]{Meredith said [\PP on Tuesday] [\CP that pigs are excellent flier].}
			\ex[??]{Meredith said [\CP that pigs are excellent fliers] [\PP on Tuesday].}
		\end{xlist}
		\ex
		\begin{xlist}
			\ex[]{Meredith said [\PP on every day I showed up] [\CP that pigs fly].}
			\ex[]{Meredith said [\CP that pigs fly] [\PP on every day I showed up].}
		\end{xlist}
	\end{xlist}
\end{exe}
The PP \emph{on Tuesday} is not an argument in these examples, and so the Projection Principle leads to the expectation that it should necessarily follow the argument~CP. In fact, however, the argument CP can follow the PP and, in some cases, is even preferred to be in that position. These examples look like counter-examples to the Projection Principle.

A very similar paradigm of contrasts arises in cases with DP internal arguments. We've seen that there is a preference for putting an internal argument that is a DP adjacent to the verb that it is an argument of. But this preference is lost when the relative length of the DP and the rest of the material in the \ibar{V} is manipulated. When the DP argument is a pronoun or is a one or two word phrase, the ordering preferences we have given an account of emerge:
\begin{exe}
	\raggedright
	\ex \label{ex3:DPfa} 
	\begin{xlist}
		\ex \begin{xlist}
			\ex[]{Mako explained it to her classmates.}
			\ex[*]{Mako explained to her classmates it.}
		\end{xlist}
		\ex
		\begin{xlist}
			\ex[]{Mako explained stuff to her classmates.}
			\ex[?*]{Mako explained to her classmates stuff.}
		\ex
		\end{xlist}
			\ex[]{Mako explained Endocentricity to her classmates.}
			\ex[?*]{Mako explained to her classmates Endocentricity.}
	\end{xlist}
\end{exe}
The contrasts in (\ref{ex3:DPfa}b) and (\ref{ex3:DPfa}c) aren't as strong as the contrast in (\ref{ex3:DPfa}a), and can be overcome if the DP is given prosodic stress. But if the relative length of the DP and PP is manipulated, even this isn't necessary to get the rogue word-orders.
\begin{exe}
	\raggedright
	\ex \label{ex3:DPfb} 
	\begin{xlist}
			\ex[]{Mako explained the problems kept in her files to them.}
			\ex[]{Mako explained to them the problems kept in her files.}
	\end{xlist}
\end{exe}
These should be Case filter violations. In fact, these same forces seem to control the relative order of argument DPs and \underline{non-argument} PPs. The PP \emph{on Tuesday} is a non-argument in (\ref{ex3:DPfc}), and yet its relative order with the argument DP depends on the same factors that we've seen controlling the relative orders of arguments.
\begin{exe}
	\raggedright
	\ex \label{ex3:DPfc} 
	\begin{xlist}
		\ex 
		\begin{xlist}
			\ex[]{Mako explained it on Tuesday.}
			\ex[*]{Mako explained on Tuesday it.}
		\end{xlist}
		\ex
		\begin{xlist}
			\ex[]{Mako explained stuff on Tuesday.}
			\ex[?*]{Mako explained on Tuesday stuff.}
		\end{xlist}
		\ex
		\begin{xlist}
			\ex[]{Mako explained the utility of Endocentricity on Tuesday.}
			\ex[]{Mako explained on Tuesday the utility of Endocentricity.}
		\end{xlist}
		\ex
		\begin{xlist}
			\ex[?]{Mako told the story I'd made up about pigs on Tuesday.}
			\ex[]{\label{ex3:pigs}Mako told on Tuesday the story that I'd made up about pigs.}
		\end{xlist}
	\end{xlist}
\end{exe}
These are violations of the Projection Principle.

All of these cases don't fit our present grammar, so we should diagnose what the problem with our grammar is and then fix it. One diagnosis is that the principles which are violated by these grammatical examples should be loosened; perhaps we should replace the Projection Principle with something looser and remove the Case filter, or the adjacency condition built into it which forces argument DPs to be next to the verb. This would permit all of these word orders, and we could let other forces---perhaps ones having to do with ease of parsing---choose among them. This would involve revising the conditions on Accusative Case assignment and the Projection Principle to~(\ref{ex3:revised}).
\begin{exe}
	\raggedright
	\ex \label{ex3:revised} 
	\begin{xlist}
		\ex Accusative Case is assigned by a verb or preposition to a position inside their highest \ibar{X} projection.
		\ex An internal $\theta$-role is assigned by \obar{X} to a position inside its highest~\ibar{X}.
	\end{xlist}
\end{exe}
This revision will handle the cases we've just seen, if it can be wedded to something that chooses correctly from the range of word order options the looser conditions in (\ref{ex3:revised}) permit.

This is not the direction usually taken, however, and I will show you why in the next section. We have to learn about another construction type for the problems attendant with (\ref{ex3:revised}) to become visible. The solution usually taken involves a more involved modification to the grammar, one that is a bit of overkill for the problems we are now facing. But, as we shall see, this modification handles a large range of additional cases; there is an independent need for the modification I am about to develop. The leading idea behind this modification to our grammar is that the DPs and CPs that we see in the positions that the Projection Principle and Case filter predict they should not be are \underline{also} in the positions that the Projection Principle and Case filter predict they are. Or, putting it differently, in these examples there is a violation of Contiguity. The DPs and CPs that are in the ``wrong'' positions are actually in the right positions syntactically, but they are spoken in a position different than where Contiguity says they should be. We will see many cases of this type. Our evidence about the phrase structure is at odds with the strings that those phrases produce. We are going to need a general model of such cases.

One popular way of capturing what's happening in such cases is to let a single sentence get mapped onto not just one parse, but a set of parses. This set of parses can differ with respect to where one, or more, phrases are positioned. The linearization algorithm, and so Contiguity, sees only one of those parses, but the syntax and semantics can see the others. The picture is~(\ref{ex3:der1s}).
\begin{exe}
	\raggedright
	\ex \label{ex3:der1s} 
	Sentence =_{\text{def.}} \{ P_1, P_2,$\ldots$,P_n\}, where
	\begin{xlist}
		\ex Each P_i is a structure produced by the \ibar{X} Skeleton, and
		\ex One P_i is subject to the linearization algorithm.
	\end{xlist}
\end{exe}

We will need some way of defining the set of parses that are assigned to a sentence. We want sentences to be given a syntactic representation that gets interpreted semantically in one way. Our goal, recall, is to characterize which substrings of a sentence have a compositional meaning. So we will want this set of parses to all be related to one parse which characterizes the way in which these semantic building blocks are assembled. To capture this, I'll add another criterion to the list that defines a sentence.
\begin{exe}
	\raggedright
	\ex \label{ex3:der2s} 
	Sentence =_{\text{def.}} \{ P_1, P_2,$\ldots$,P_n\}, where
	\begin{xlist}
		\ex Each P_i is a structure produced by the \ibar{X} Skeleton, and
		\ex One P_i is subject to the linearization algorithm.
		\ex One P_i is the result of lexical insertion and is semantically interpreted.
	\end{xlist}
\end{exe}
Part of semantic interpretation is assigning $\theta$ roles, so the parse that is semantically interpreted will also be the one where the linking rules and Theta Criterion are satisfied. Now that we have one designated parse that captures the structure of the semantically interpreted phrases, we can define the other parses in the set with relations to that designated parse. We'll say that the parses in the set are just those that could be produced by applying a rule to the parse that gets semantically interpreted. The parse that is semantically interpreted is sometimes called a ``D-structure,'' and the parse that is linearized is called an ``S-structure.''\marginnote[0pt]{``D-structure'' is meant to evoke ``deep structure,'' and ``S-structure'' is meant to evoke ``surface structure.'' The framers of this theory must have imagined that meanings bubbled up from the depths and became visible only once they were afloat.}
\begin{exe}
	\raggedright
	\ex \label{ex3:der22s} 
	Sentence =_{\text{def.}} \{ P_1, P_2,$\ldots$,P_n\}, where
	\begin{xlist}
		\ex Each P_i is a structure produced by the \ibar{X} Skeleton, and
		\ex One P_i is subject to the linearization algorithm. (S-structure), and
		\ex One P_i is the result of lexical insertion and is semantically interpreted (D-structure)j, and
		\ex Each P_i $\neq$ D-structure is formed from the D-structure by a rule.
	\end{xlist}
\end{exe}
We will see that the rules might not create structures that are producible by the \ibar{X} Skeleton. One hypothesis is that the only deviations from what the \ibar{X} Skeleton allows are those created by the rules. We can capture this by imposing the requirements of the \ibar{X} Skeleton on D-structure alone.
\begin{exe}
	\raggedright
	\ex \label{ex3:der3s} 
	Sentence =_{\text{def.}} \{ P_1, P_2,$\ldots$,P_n\}, where
	\begin{xlist}
		\ex P_1 is a structure produced by the \ibar{X} Skeleton, and lexical insertion, and it is semantically interpreted (D-structure), and
		\ex One P_i is subject to the linearization algorithm (S-structure), and
		\ex Each P_{i \neq 1} is formed from P_{i-1} by a rule.
	\end{xlist}
\end{exe}

What (\ref{ex3:der3s}) defines is a cloud of parses that are different from a D-structure in the way that one rule specifies. In the rest of this book, we will be investigating what the rules that form this set look like, and we'll discover that some rules can apply to the output of other rules. We will need to relax (\ref{ex3:der3s}c) so that it allows a rule to apply to parses other than the D-structure. The set of parses, then, will have structure: there will be chains of parses in it that are formed by serial application of rules. These chains are called ``derivations,'' and we will think of them as a series.
\begin{exe}
	\raggedright
	\ex \label{ex3:deriv} 
	$\mathcal{D}$ = (P_1, P_2\ldots,P_n) is a derivation iff:
	\begin{xlist}
		\ex P_1 is formed by the \ibar{X} Skeleton and Lexical Insertion, and it is semantically interpreted, and
		\ex Each P_{i+1} is formed from P_{i} by a rule, and
		\ex P_n is subject to the linearization algorithm.
	\end{xlist}
	\ex A sentence is the set of derivations which have the same D-structure and S-structure.
\end{exe}

With these additions to our model, we can now formulate a rule that will characterize the cases where argument DP and CPs are in positions that the Projection Principle and Case filter don't predict. The rule is one that \citet{Ross:1967} devised.\marginnote[0pt]{In fact, Ross devised two rules, one for DPs and one for CPs, which is called ``extraposition.'' It's not clear to me that they need to be distinguished, so I've collapsed them.}
\begin{exe}
	\raggedright
	\ex \label{ex3:hnps} Complex NP Shift\\
	Move a DP or CP so that it is right adjoined to the smallest IP containing it.
\end{exe}
Ross's work predates the move to giving nominal phrases a determiner as the head, and so he named his rule after NPs, one of the terms the rule can move. Reference to ``complex'' in the name is intended to reference the fact that this rule is sensitive to the size of the DP or CP being moved. We should define a few terms that the rule makes reference to. First ``smallest'' and ``right adjoined.''
\begin{exe}
	\raggedright
	\ex \label{ex3:smallest} 
	XP is smaller than YP, if XP dominates YP.
	\ex \label{ex3:rightadjoined}
	Adjoining XP to YP involves the following steps.
	\begin{xlist}
		\ex YP is replaced by another YP, YP$'$
		\ex YP$'$ has the old YP as one daughter and XP as the other daughter.
	\end{xlist}
\end{exe}
We say that XP right adjoins to YP if the linearization algorithm makes YP<XP.\marginnote[0pt]{The ``<'' symbol means precedes.} And we say that XP left adjoins to YP if the linearization algorithm makes XP<YP. Finally, we should define ``move.''
\begin{exe}
	\raggedright
	\ex \label{ex3:movedef} $\alpha$, a node in a phrase marker tree, \underline{moves} to $\beta$, a position in a phrase marker tree, means:
	\begin{xlist}
		\ex $\alpha$ is put in position $\beta$, and
		\ex $\alpha$ is removed from its original position.
	\end{xlist}
\end{exe}
We will see that there are a variety of rules which make use of move.
\pagebreak

Ross's rule allows (\ref{ex3:pigs}) to have the derivation in~(\ref{ex3:pigsd}).
\begin{exe}
	\raggedright
	\exr{ex3:pigs} Mako told on Tuesday the story that I'd made up about pigs.
	\ex \label{ex3:pigsd}
	\small
	\begin{forest}
	[IP
		[DP [Mako, roof]]
		[\ibar{I} 
			[\obar{I} [past] ]
			[VP [\ibar{V} 
				[\ibar{V} 
					[\obar{V} [tell] ]
					[DP [the story that\\ I'd made up\\ about pigs, roof] ]
				 ]
				[PP [on Tuesday, roof] ]
			] ]
		]
	]
	\end{forest}
	{\Huge ,}
	\small
	\begin{forest}
	[IP	
		[IP	
			[DP ]
			[\ibar{I} 
				[\obar{I} [past] ]
				[VP [\ibar{V} 
						[\ibar{V} [\obar{V} [tell]] ]
						[PP [on Tuesday, roof] ]
					 ]
				] ]
		 ]
		[DP [the story that\\I'll make up\\about pigs, roof] ]
	]
	\end{forest}
\end{exe}
In this simple derivation, the D-structure is the first parse. It corresponds to the parses our system has been able to produce up to now. It satisfies the Projection Principle, the Theta Criterion, and the Case filter. It is the parse that is semantically interpreted, so it represents the phrases whose meanings are compositional. The second parse is derived from this D-structure by applying Complex NP Shift, and it is the S-structure (in this case). It is this phrase structure that the linearization algorithm produces a string from. As a consequence, the internal argument DP is pronounced in a position different than where its meaning is computed. In this simple case, the sentence (\ref{ex3:pigs}) has only this one derivation in its set. This is the only series of parses that have the D-structure and S-structure shown in~(\ref{ex3:pigsd}).

Stowell uses Ross's Complex NP Shift operation to control the relative position of PP and CP arguments too. But he suggests that Case plays a role in making Complex NP Shift relatively favored. He speculates that finite CPs, like DPs, must receive Case but that unlike DPs they cannot sit in Case-marked positions at S-structure. He calls this the ``Case Resistance Principle.''
\begin{exe}
\ex \label{ex3:caseresist} Case Resistance Principle\\
A (finite) CP may not be in a Case marked position at S-structure.
\end{exe}
The Case Resistance Principle requires finite CPs to be moved by S-structure to some non-Case marked position. NP Shift is capable of doing this, and thus, in cases where a complement CP shows up string finally in a VP, it has satisfied Stowell's injunction against surfacing in Case marked positions by undergoing NP Shift. There is, surprisingly, a certain amount of evidence for this picture.

Consider, for example, situations where the finite clause is the ``subject'' of the clause. Jan Koster\cite{Koster:1978} argued that a subject finite CP isn't actually in the Nominative Case-marked, Specifier of IP position. One of Koster's ar\-gu\-ments makes reference to a process that is found in certain question-types in English. Normally, in these question contexts, it is possible for \obar{I} to be at the front of a sentence, as in (\ref{ex3:kos1}).\marginnote[0pt]{Koster's conclusions, and a modernization of some of his arguments, can be found in a more thorough form in \citet{Alrenga:2005}. See also \citet{Elliott:2017} for critical discussion.}
\begin{exe}
	\ex \label{ex3:kos1} Maria will put the book on the table. $\rightarrow$\\
						 Will Maria put the book on the table?
\end{exe}
How precisely this happens is the subject of the next chapter. What's relevant here is that this word order is blocked if it interacts with another process that moves something to the left edge of IP, as in~(\ref{ex3:top1}).
\begin{exe}
\ex \label{ex3:top1} Maria will put the book on the table. $\rightarrow$\\
					 On the table, Maria will put the book.
\end{exe}
In (\ref{ex3:top1}) we are seeing the action of another movement operation, sometimes called Topicalization, which we can formulate as~(\ref{ex3:top}).
\begin{exe}
	\raggedright
	\ex \label{ex3:top} Topicalization\\
	Move an XP so that it is left adjoined to the smallest IP containing it.
\end{exe}
The derivation for (\ref{ex3:top1}) is~(\ref{ex3:topd}).
\begin{fullwidth}
	\medskip
\begin{exe}
	\raggedright
	\ex \label{ex3:topd} 
	\small
	\begin{forest}
	[IP
		[DP [Maria, roof]]
		[\ibar{I} 
			[\obar{I} [will] ]
			[VP [\ibar{V} 
				[\ibar{V} 
					[\obar{V} [put] ]
					[DP [the book, roof] ]
				 ]
				[PP [on the table, roof] ]
			] ]
		]
	]
	\end{forest}
	{\Huge ,}
	\small
	\begin{forest}
	[IP
		[DP [the book, roof] ]
		[IP
			[DP [Mary, roof]]
			[\ibar{I} 
				[\obar{I} [will] ]
				[VP [\ibar{V} 
					[\obar{V} [put] ]
					[PP [on the table, roof] ]
				] ]
			 ]
		]
	]
	\end{forest}
\end{exe}
\medskip
\end{fullwidth}

Topicalization cannot apply to a sentence that has the word-order of a question, that is, a word order that puts the \obar{I} first. (\ref{ex3:kos2}) illustrates.
\begin{exe}
	\ex \label{ex3:kos2} 
	\begin{xlist}
	\ex[]{Maria will put the book on the table $\rightarrow$}
	\ex[]{On the table, Maria will put the book. $\rightarrow$}
	\ex[*]{Will on the table, Maria put the book?}
	\end{xlist}
\end{exe}

Now, interestingly, for many speakers of English the presence of a finite CP as a subject of a sentence is also incompatible with this \obar{I} first kind of sentence.
% \pagebreak
\begin{exe}
	\ex \label{ex3:kos3} 
	\begin{xlist}
	\ex[*]{Should that Maria has left be obvious?}
	\ex[*]{Does that Paula sleep late bother Andries?}
	\end{xlist}
\end{exe}
This would follow if finite CPs are driven from the nominative Case marked Specifier of IP, and adjoin to the left of IP in these cases. Stowell's method of forcing Complex NP Shift to apply to complement CPs would extend to this scenario as well. CPs start out in the nominative Case-marked position, but are driven from there in order to form an S-structure that meets the Case Resistance Principle. This is shown in the derivation in~(\ref{ex3:crpd}).
\begin{fullwidth}
\medskip
\begin{exe}
	\raggedright
	\ex \label{ex3:crpd} 
	\small
	\begin{forest}
	[IP
		[CP [that pigs fly, roof]]
		[\ibar{I} 
			[\obar{I} [will] ]
			[VP [\ibar{V} 
				[\obar{V} [be] ]
				[AP [\ibar{A} [\obar{A} [obvious]]] ]
			] ]
		]
	]
	\end{forest}
	{\Huge ,}
	\small
	\begin{forest}
	[IP
		[CP [that pigs fly, roof]]
		[IP [\ibar{I} 
			[\obar{I} [will] ]
			[VP [\ibar{V} 
				[\obar{V} [be] ]
				[AP [\ibar{A} [\obar{A} [obvious]]] ]
			] ]
		]
		]
	]
	\end{forest}
\end{exe}
\medskip
\end{fullwidth}

This evidence all points to the fact that finite CPs move either to the left or to the right linear edges of the sentences they are part of. If CPs are prevented from being in Case marked positions in the S-structure, this would get us close to deriving this fact. But is there evidence for the motivation Stowell proposes for this movement? In particular, is there motivation for the claim that finite CPs, like DPs, require Case?

One piece of suggestive evidence comes from the class of verbs that permit both DPs and finite CPs. These are only just those verbs that already exceptionally allow two DPs. These include verbs like \emph{promise}, \emph{tell}, and \emph{show}; they are known as ``double object'' verbs. Some examples indicating that double object verbs also take DP and CP are in~(\ref{ex3:cpfilter1}).
\begin{exe}
	\ex \label{ex3:cpfilter1} 
	\begin{xlist}
	\ex[]{Karen promised me that she would sing.\\ Karen promised me the ring}
	\ex[]{Wendell told me that he can't stand Karen's singing.\\ Wendell told me the story.}
	\ex[]{Annahita showed me that she cares.\\ Annahita showed me her concern.}
	\end{xlist}
\end{exe}
This isn't completely the case, as (\ref{ex3:cpfilter2}) is a counterexample.
\begin{exe}
	\ex \label{ex3:cpfilter2} 
	\begin{xlist}
	\ex[]{Shai persuaded Adam that he should go.}
	\ex[*]{Shai persuaded Adam the fact}
	\end{xlist}
\end{exe}
But, so far as I know, (\ref{ex3:cpfilter2}) is the only counterexample. To the extent that there is a match in the verbs which accept DP+CP and those which accept DP+DP complements, there are grounds for believing that their surface positions are governed by the same, or similar, principles. And to the extent that the dominant principle is the Case Filter, then there is reason to conclude from these data that CPs are subject to the Case Filter as well.

This match between the distribution of DPs and CPs should be completely general if Stowell is correct. Indeed, finite CPs are distributed in sentences much like DPs are. We find them as complements to verbs, as we have seen, and in the subject position of other finite clauses, but not in the subject position of infinitives (as we shall see shortly). These are just the Case-marked positions.

But there are several differences in their distribution. In English, finite CPs are probably never found as the complements to a preposition, though of course DPs are. The only potential counterexample comes from temporal prepositions, as in (\ref{ex3:cpfilter3}).
\begin{exe}
	\ex \label{ex3:cpfilter3}\begin{xlist}
	\ex[]{I left before Cherlon arrived.}
	\ex[*]{I left before that Cherlon arrived.}
	\end{xlist} 
\end{exe}
Similarly, both adjectives and nouns can take CP complements, but not DP complements.\enlargethispage*{18pt}
\begin{exe}
	\ex \label{ex3:cpfilter4} 
	\begin{xlist}
	\ex[]{Shigeto is unhappy that he had to sing.}
	\ex[*]{Shigeto is unhappy that.}
	\end{xlist}
	\ex \label{ex3:cpfilter5}
	\begin{xlist}
	\ex[]{the proof that lemons cure cancer}
	\ex[*]{The proof this fact.}
		\end{xlist}
\end{exe}
If this has a Case Theoretic explanation, then Stowell's theory is in trouble. To my knowledge, an account of the distribution of finite CPs has not yet been formulated. We will see in the next chapter that there are verbs which do not c-select DPs and do not assign accusative Case but do c-select CPs. When this is coupled with the fact that CPs can be the internal arguments of nouns and adjectives---words that never support accusative Case---it seems unlikely that Stowell's explanation for the obligatoriness of CP movement is correct. They move not because the positions they originally occupy are Case-marked ones, and they are somehow incompatible with Case. At the time of writing, there is no consensus about what is responsible for the special placement of CPs.\marginnote[0pt]{See \citet{Moulton:2009,Moulton:2015} for a good proposal.}

Though the evidence suggests that Stowell's method of deriving the linear ordering of verbal complements meets with some success, there are problems too. One which threatens Stowell's Adjacency Condition on Case assignment, and its use in fixing the order of complements, concerns so-called ``double object'' constructions, as in:
\begin{exe}
	\ex \label{ex3:problems1} 
	\begin{xlist}
	\ex[]{Shigeto showed Luis the picture.}
	\ex[]{Shigeto baked Luis a cake.}
	\end{xlist}
\end{exe}
How can the second DP in these examples receive Case? We will have to develop some of the rest of the system that is responsible for linearizing arguments before we can engage this difficulty. So, let me ask for your indulgence and postpone examining this case until chapter~\ref{sec:thetacase}.

\section{Small Clauses and the Derived Subjects Hypothesis}
\label{sec:smallclauses}

In the last chapter, I decided that we should adopt an \ibar{X} Skeleton that forced every XP to contain an \ibar{X}. This narrows the class of possible kinds of phrases that are allowed to just those types that we have evidence for in IPs and DPs. In IPs and DPs, we discovered that there are two kinds of projections: one which contains a phrase that the other cannot. That phrase is the thing that resides in Specifier position. For IPs, the Specifier position can be occupied by either a DP or CP, and the Specifier of DP can be occupied by a DP marked with the genitive ``s.'' \ibar{I}s and \ibar{D}s don't have those phrases, and that's what makes it necessary to distinguish IP from \ibar{I} and DP from \ibar{D}. 

Forcing every phrase to have an XP and \ibar{X} does no harm, but it does raise the puzzle of what distinguishes the XP from the \ibar{X} for the majority of phrases. We haven't found any reason to think that, say, VPs are different from \ibar{V}s and this is puzzling if there really is a distinction between these phrases. Let's face this puzzle now. Let's see if we can identify what distinguishes the XP from \ibar{X} of most phrases. If we can find nothing that does, then we should transit to a different \ibar{X} theory, one that doesn't require every phrase to have an~\ibar{X}.
\begin{exe}
\ex \label{ex3:specquest} What resides in the Specifiers of:
\begin{xlist}
	\ex Adjective Phrase
	\ex Verb Phrase
	\ex Prepositional Phrase
	\ex Determiner Phrase
	\ex Adverb Phrase
\end{xlist}
\end{exe}
The structure of Adverb Phrases is not well-studied, partly because of the apparent poverty of material that can be in adverb phrases. We will generally ignore the structure of these phrases as a result. It's time, though, to see what occupies the Specifier positions of APs, VPs and PPs.

Let's begin with VPs and APs. What are the terms, if any, that stand in the Specifiers of VP and AP? One hypothesis is that VPs and APs can have external arguments in them, and these are the designated Specifiers. This can be seen in certain special circumstances, of which (\ref{ex3:smcla}) are examples.
\begin{exe}
	\ex \label{ex3:smcla} 
	\begin{xlist}
	\ex[]{She will let [them make a cabinet].}
	\ex[]{This will make [them angry at me].}
	\end{xlist}
\end{exe}
As indicated by the brackets, it is thought that the strings following \emph{let} and \emph{make} form a single phrase which is serving as the complement to these verbs.

Why don't we instead believe that these verbs are followed by two complements, the first of which is \emph{them}? Because \emph{them} is not an argument of the verb. We can see from other examples that this isn't the case. In (\ref{ex3:smcla1}), for example, this position can be occupied by an expletive.
\pagebreak
\begin{exe}
	\ex \label{ex3:smcla1} 
	\begin{xlist}
	\ex[]{I will let [it seem that there is a mistake on the handout].}
	\ex[]{This will make [it obvious that there is a mistake on the handout].}
	\end{xlist}
\end{exe}
Expletive \emph{it} can only stand in a Specifier position. Thus, the positions following \emph{let} and \emph{make} in these cases must be Specifier positions.
%
%//todo Check to make sure that at this point in the text, the reader has been taught that expletives cannot be direct objects.
%

Note that the bracketed strings in (\ref{ex3:smcla}) are, indeed, an AP and VP, and not IPs. This, we know, because the strings which follows these specifier positions are characterized by the rules which generate VPs and APs, and not those which characterize IPs. As (\ref{ex3:smcla3}) indicates, an \obar{I} is not found in these strings, for instance.
\begin{exe}
	\ex \label{ex3:smcla3} 
	\begin{xlist}
	\ex[*]{I will let [her should make a cabinet].}
	\ex[*]{I will let [her to make a cabinet].}
	\ex[*]{I will let [her makes a cabinet].}
	\end{xlist}
\end{exe}

From these observations we will want to give the examples in (\ref{ex3:smcla}) the parses in (\ref{ex3:smclatree}).%\enlargethispage*{20pt}
\begin{exe}
\ex \label{ex3:smclatree}
a.
\small
\begin{forest}
[\ibar{V}
	[\obar{V} [let] ]
	[VP
		[DP [them, roof]]
		[\ibar{V}
			[\obar{V} [make] ]
			[DP [a cabinet, roof] ]
		 ]
	 ]
]
\end{forest}
b.
\small
\begin{forest}
[\ibar{V}
	[\obar{V} [make] ]
	[AP
		[DP [them, roof]]
		[\ibar{A}
			[\obar{A} [angry] ]
			[PP [at me, roof] ]
		 ]
	 ]
]
\end{forest}
\end{exe}
These complements to \emph{let} and \emph{make} are sometimes called ``small clauses,'' and it is to Stowell\cite{Stowell:1978}, once again, that we owe this way of viewing them. If this is correct, we have now discovered what stands in the Specifier of AP and VP. These positions can contain a ``subject'' argument.

A similar conclusion can be reached for PPs, though here the evidence is somewhat more meager. Cases such as (\ref{ex3:ppsmcla}) plausibly involve small clause PPs.
\begin{exe}
\ex \label{ex3:ppsmcla} 
\begin{xlist}
	\ex I want those kids in the bathtub!
	\ex She saw them in the next room.
\end{xlist}
\end{exe}
Both of these examples can have a meaning in which it seems that the verb has just one object. In (\ref{ex3:ppsmcla}a), the object of \emph{want} seems to be, roughly, what \emph{them in the bathtub now} describes. On this meaning, the content of my desire is for those damn kids to be in the tub.\marginnote[-20pt]{Of course, (\ref{ex3:ppsmcla}a) has another meaning, one in which the content of my desire is children, and more particularly, those children in the tub. That meaning makes the object of \emph{want} a DP headed by \emph{kids}.} Under our present understanding of how \mbox{$\theta$-roles} are assigned, this requires that (\ref{ex3:ppsmcla}a) have the parse in (\ref{ex3:ppsmaclatree1}).
\begin{exe}
\ex \label{ex3:ppsmaclatree1} 
\small
\begin{forest}
[\ibar{V}
	[\obar{V} [want] ]
	[PP
		[DP [those kids, roof]]
		[\ibar{P}
			[\obar{P} [in] ]
			[DP [the bathtub, roof]]
		 ]
	 ]
]
\end{forest}
\end{exe}

Similarly, among its several meanings, (\ref{ex3:ppsmcla}b) can describe a situation in which what was seen was \emph{them in the room}. This reading arises when \emph{in the room} does not describe where the seeing event took place---that reading is manufactured by letting \emph{in the room} modify \emph{see}---but instead describes the location of the referent of \emph{them}. To produce such a reading, we want \emph{them in the room} to be the object of \emph{see}, and this is achieved by a parse like (\ref{ex3:ppsmaclatree2}).
\begin{exe}
\ex \label{ex3:ppsmaclatree2} 
\small
\begin{forest}
[\ibar{V}
	[\obar{V} [see] ]
	[PP
		[DP [them, roof]]
		[\ibar{P}
			[\obar{P} [in] ]
			[DP [the room, roof] ]
		 ]
	 ]
]
\end{forest}
\end{exe}
In both of the examples in (\ref{ex3:ppsmcla}), then, there is reason to believe that the verbs are followed by one PP, whose Specifier is filled with the DP argument.

Small clauses, incidentally, furnish evidence that the Complex NP Shift solution to the problem we faced in the previous section about complement order is better than relaxing the Projection Principle and Case filter. Recall that we considered as a way of accounting for (\ref{ex3:wrongdpo}) that argument DPs to be positioned higher than non-argument DPs within \ibar{V}, and let whatever forces care about the relative lengths of the DP and non-argument decide which of the two possible positions is best.
\begin{exe}
	\raggedright
	\ex \label{ex3:wrongdpo} 
	She ate on Tuesday all the cookies she had brought.
\end{exe}
But we can see from~(\ref{ex3:smclhnps}) that whatever allows argument DPs to have the exceptional VP-last position in (\ref{ex3:wrongdpo}), should allow argument DPs that aren't arguments of the relevant \ibar{V} to have that position too.
\begin{exe}
	\raggedright
	\ex \label{ex3:smclhnps} 
	She made angry every child that had counted on having cookies for breakfast.
\end{exe}
(\ref{ex3:smclhnps}) can have the derivation shown in~(\ref{ex3:smhpd}), if Complex NP Shift is responsible for these word orders.
\begin{fullwidth}
\medskip
\begin{exe}
	\raggedright
	\ex \label{ex3:smhpd} 
	\small
	\begin{forest}
	[IP
		[DP [she, roof]]
		[\ibar{I} 
			[\obar{I} [past] ]
			[VP [\ibar{V}
				[\obar{V} [make] ]
				[AP
					[DP [every child that had\\ counted on having\\ cookies for breakfast, roof]]
					[\ibar{A} [\obar{A} [angry]]]
				 ]
			] ]
		]
	]
	\end{forest}
	{\Huge ,}
	\small
	\begin{forest}
	[IP
		[IP
		[DP [she, roof]]
		[\ibar{I} 
			[\obar{I} [past] ]
			[VP [\ibar{V} 
				[\obar{V} [make] ]
				[AP [\ibar{A} [\obar{A} [angry]]] ]
			] ]
		]
		]
		[DP [every child that had\\ counted on having\\ cookies, roof]]
	]
	\end{forest}
\end{exe}
\medskip
\end{fullwidth}
Simply relaxing the Projection Principle, or Case filter, would not explain (\ref{ex3:smclhnps}). But Complex NP Shift does.

Consider now the examples in (\ref{ex3:dersub1}).
\begin{exe}
\ex \label{ex3:dersub1} 
\begin{xlist}
	\ex They should make a cabinet.
	\ex They are angry at me.
	\ex They are in the next room.
\end{xlist}
\end{exe}
These sentences contain the same VP, AP and PP that constituted small clauses in our earlier examples. Moreover, the \mbox{$\theta$-role} borne by the subject in (\ref{ex3:dersub1}) seems to be the same as that borne by the DP in the examples with small clauses. Indeed, the meaning of the sentences in (\ref{ex3:dersub1}) seems to be identical to those of the small clauses in the previous examples, except for that meaning introduced by the modal \emph{should} (in (\ref{ex3:dersub1}a)) and the verb \emph{be} (in (\ref{ex3:dersub1}b,c)). A number of linguists independently noted that there could well be a small clause within (\ref{ex3:dersub1}), but that this arrangement is obscured by a movement transformation that relocates the subject.\footnote{See \citet{Kuroda:1988}, \citet{Kitagawa:1986}, \citet{Fukui:1986a}, \citet{Bowers:1993}, and \citet{Koopman:1991} for proposals along these lines that fit into the general theoretical assumptions defended here, and \citet{Fillmore:1968} and especially \citet{McCawley:1968a} for early antecedents to the idea.} Indeed, the Case filter based theory about where argument DPs must be predicts that (\ref{ex3:dersub1}a) could have a small clause in it at D-structure. Suppose that the \mbox{$\theta$-role} that \emph{they} bears in (\ref{ex3:dersub1}a) is assigned to the same position that it is in our small clause example in (\ref{ex3:smcla}a), that is to the Specifier of VP. Then the Theta Criterion would require the D-structure of (\ref{ex3:dersub1}a) to be (\ref{ex3:dersub1t}).
\begin{exe}
\ex \label{ex3:dersub1t} 
\small
\begin{forest}
[\ibar{I}
	[\obar{I} [should] ]
	[VP
		[DP [they, roof] ]
		[\ibar{V}
			[\obar{V} [make] ]
			[DP [a cabinet, roof] ]
		 ]
	 ]
]
\end{forest}
\end{exe}
The subject argument, \emph{they}, bears nominative Case, and yet in this phrase marker, it is not in the position our system assigns nominative Case to. That position is the Specifier of a finite IP, which is the position that \emph{they} is, in fact, spoken in. Imagine, then, that there is a movement transformation that relocates \emph{they} into the Specifier of IP, as in~(\ref{ex3:dersub1t2}).\marginnote[0pt]{I will use this notation involving arrows to represent adjacent parses in a derivation. The tail of the arrow shows where the moved phrase is in the first of the two parses, and the head of the arrow shows where the phrase is after the relevant rule has applied, forming the second parse in the pair.}
\begin{exe}
\ex \label{ex3:dersub1t2} 
\small
\begin{forest}
[IP
	[DP [they, roof, name=binder]]
	[\ibar{I} 
		[\obar{I} [should] ]
		[VP
			[{\ }, name=trace ]
			[\ibar{V}
				[\obar{V} [make] ]
				[DP [a cabinet, roof] ]
			 ]
		 ]
	]
]
\draw[->,>=stealth', ultra thin, gray] (trace) to[out=south west,in=south] node[pos=0.75,xshift=10mm]{} (binder);
\end{forest}
\end{exe}
The same derivations could be present in (\ref{ex3:dersub1}b) and (\ref{ex3:dersub1}c) as well. They would be as indicated in~(\ref{ex3:dersub1tt}).
\begin{exe}
\ex \label{ex3:dersub1tt} 
\hspace*{-25pt}
\small
\begin{forest}
[IP
	[DP [they, roof, name=binder]]
	[ \ibar{I}
		[\obar{I} [present] ]
		[VP [\ibar{V} 
			[\obar{V} [be] ]
			[AP
				[{\ }, name=trace ]
				[\ibar{A}
					[\obar{A} [angry] ]
					[PP [at me, roof]]
				 ]
			 ]
		] ]
	]
]
\draw[->,>=stealth', ultra thin, gray] (trace) to[out=south west,in=south] node[pos=0.75,xshift=10mm]{} (binder);
\end{forest}
\hspace*{-15pt}
\begin{forest}
[IP
	[DP [they, roof, name=binder]]
	[\ibar{I} 
		[\obar{I} [present] ]
		[VP [\ibar{V} 
			[\obar{V} [be] ]
			[PP
				[{\ }, name=trace ]
				[\ibar{P}
					[\obar{P} [in] ]
					[DP [the next room, roof] ]
				 ]
			 ]
		] ]
	]
]
\draw[->,>=stealth', ultra thin, gray] (trace) to[out=south west,in=south] node[pos=0.75,xshift=10mm]{} (binder);
\end{forest}
\end{exe}
The conjecture that the subject arguments of IPs start out inside small clauses and move into the position they are spoken goes by the name of ``the derived subjects hypothesis.'' The movement transformation responsible for relocating the subjects is known as ``Argument Movement,'' or more commonly: \mbox{A-Movement}.
\begin{exe}
\ex \label{ex3:amove} Argument Movement\\
Move an argument XP to a Specifier position.
\end{exe}

We shall examine some evidence in support of the derived subjects hypothesis in the section on ``Quantifier Floating,'' later in this chapter. There are other environments in which A Movement is alleged to occur, which we'll need to examine as well. But before we turn to those tasks, we need to make some modifications to our system in the light of this hypothesis.

The first modification is to sharpen the conditions on Case assignment in order for them to produce the desired effects under the derived subjects hypothesis. Our present conditions on Case assignment are~(\ref{ex3:casecond1}).
\begin{exe}
\ex \label{ex3:casecond1} Conditions on Case assignment
\begin{xlist}
	\ex Nominative Case is assigned by finite \obar{I} to its Specifier position.
	\ex Accusative Case is assigned by verbs and prepositions to an adjacent position.
\end{xlist}
\end{exe}
The description of where accusative Case is assigned does not produce the right results when coupled with the derived subjects hypothesis. To see this, consider the D-structures of some of the examples we've looked at.
\begin{exe}
\ex \label{ex3:govill1} 
\small
a.
\hspace*{-14pt}
\begin{forest}
[\ibar{I}
	[\obar{I} [should]]
	[VP 
		[DP [they, roof]]
		[\ibar{V}
			[\obar{V} [make] ]
			[DP [a cabinet, roof]]
		 ]
	]
]
\end{forest}
\hspace*{-5pt}
b.
\begin{forest}
[\ibar{I}
	[\obar{I} [present]]
	[VP [\ibar{V}
		[\obar{V} [be] ]
		[AP
			[DP [they, roof]]
			[\ibar{A}
				[\obar{A} [angry] ]
				[PP [at me, roof] ]
			 ]
		 ]
	]]
]
\end{forest}
\end{exe}
The subjects in these examples cannot be assigned accusative Case; they only appear with nominative. But the statement of accusative Case assignment in (\ref{ex3:casecond1}b) does not ensure this. The subject is adjacent to the verb \emph{make} in (\ref{ex3:govill1}a), and in (\ref{ex3:govill1}b) it is adjacent to the verb \emph{be}. The case in (\ref{ex3:govill1}b) contrasts minimally with (\ref{ex3:govill2}).
\begin{exe}
\ex \label{ex3:govill2} 
\small
\begin{forest}
[IP
	[DP [she, roof]]
	[\ibar{I} 
		[\obar{I} [past]]
		[VP [\ibar{V}
			[\obar{V} [make] ]
			[AP
				[DP [them, roof]]
				[\ibar{A}
					[\obar{A} [angry] ]
					[PP [at me, roof] ]
				 ]
			 ]
		] ]
	]
]
\end{forest}
\end{exe}
In (\ref{ex3:govill2}), \emph{make} assigns accusative Case to \emph{them}. Because of the parallel structures of (\ref{ex3:govill2}) and (\ref{ex3:govill1}b), the usual conclusion is that the difference in availability of accusative Case must devolve to the difference in the verbs involved. Unlike \emph{make}, \emph{be} does not assign accusative Case. If that's the case, then not every verb can assign accusative Case. Instead, it appears that (\ref{ex3:casevar}) is true.
\begin{exe}
\ex \label{ex3:casevar} The ability to assign accusative Case is part of the lexical specification of a verb.
\end{exe}

The response to (\ref{ex3:govill1}a) needs to be more general. No matter what verb is involved, it is not able to assign accusative Case to its Specifier position. The usual answer given to this problem, one that extends to other cases, as we shall see, is based on the structural relationship between the verb and the positions it can assign Case to. A verb is able to assign accusative Case to its sister, as in (\ref{ex3:govill1}a), or to a position within its sister, as in (\ref{ex3:govill2}). But it's not able to assign Case to anything outside of the phrases it is dominated by. To capture this range of positions, we define the following relation.
\begin{exe}
\ex \label{ex3:ccomdef} $\alpha$\ \textsc{c-commands} $\beta$\ iff 
\begin{xlist}
	\ex every phrase dominating $\alpha$\ also dominates $\beta$, and
	\ex $\alpha$\ does not dominate $\beta$.
\end{xlist}
\marginnote[-25pt]{$\alpha$ dominates $\beta$ if $\beta$ is part of $\alpha$ and $\beta \neq \alpha$. Or, alternatively, $\beta$ is dominated by $\alpha$ if the phrase structure rule expanding $\alpha$ leads to a rule that introduces $\beta$.  The ``c-command'' relation comes from \citet{Reinhart:1976}, whose definition is slightly different from the one produced here. (The central difference is that she uses ``branching node'' in place of my ``phrase.'') The name ``c-command'' is short for ``constituent command.'' Reinhart presents her proposal as a modification of the relation ``command'' in \citet{Langacker:1969}. It is close to the notion ``in construction with'' that is found in \citet{Klima:1964}, which Langacker cites as his model.}
\end{exe}
The c-command relation is useful in describing a wide range of syntactic phenomena, and so it will arise in several other contexts in these lectures. The definition in (\ref{ex3:ccomdef}) anticipates these later uses, and so goes a bit beyond what is necessary for the Case assignment relation. Only (\ref{ex3:ccomdef}a) is needed here. (\ref{ex3:ccomdef}b) prevents things from c-commanding themselves or their contents, and neither of these scenarios are relevant for the Case assignment relation. Armed with c-command, we can restrict the condition on accusative Case assignment to (\ref{ex3:newacc1}) and, thereby, prevent verbs from assigning accusative Case to their Specifier position.
\begin{exe}
\ex \label{ex3:newacc1} Accusative Case is assigned by some verbs and prepositions to positions that they c-command and are adjacent to.
\end{exe}

We need one additional change to the condition on accusative Case assignment. This change is made necessary by the contrast in grammaticality of examples like~(\ref{ex3:govill3}).
\begin{exe}
\ex \label{ex3:govill3}
\begin{xlist}
	\ex[*]{She made him running easy.}
	\ex[]{She saw him running.}
\end{xlist}
\end{exe}
The complement to \emph{saw} in (\ref{ex3:govill3}b) is a kind of non-finite IP called a ``gerund.'' These IPs are headed by the inflectional morpheme \emph{ing} that appears on the following verb. The subject of this IP in (\ref{ex3:govill3}b) receives its accusative Case from \emph{see}; its S-structure parse is (\ref{ex3:govill3t}). 
\begin{exe}
	\ex \label{ex3:govill3t} 
	\small
	\begin{forest}
	[IP
		[DP [she, roof]]
		[\ibar{I} 
			[\obar{I} [past] ]
			[VP [\ibar{V}
				[\obar{V} [see] ]
				[IP
					[DP [him, roof]]
					[\ibar{I}
						[\obar{I} [ing] ]
						[VP [run, roof]]
					 ]
				 ]
			] ]
		]
	]
	\end{forest}
\end{exe}
The condition on accusative Case assignment in (\ref{ex3:newacc1}) correctly captures the grammaticality of this example; \emph{see} c-commands and is adjacent to~\emph{him}.

Consider now the contrasting (\ref{ex3:govill3}a). In this example, the gerund is the subject of a small clause, as indicated in~(\ref{ex3:govill3tt}). 
\begin{exe}
\ex \label{ex3:govill3tt} 
\small
\begin{forest}
[IP
	[DP [she, roof]]
	[\ibar{I} 
		[\obar{I} [past] ]
		[VP [\ibar{V}
			[\obar{V} [make] ]
			[AP
				[IP
					[DP [him, roof]]
					[\ibar{I}
						[\obar{I} [ing] ]
						[VP [run, roof]]
					 ]
				 ]
				[\ibar{A} [easy, roof] ]
			 ]
		] ]
	]
]
\end{forest}
\end{exe}
The condition on accusative Case assignment in (\ref{ex3:newacc1}) would allow \emph{make} to assign Case to \emph{him} in (\ref{ex3:govill3tt}): \emph{make} c-commands, and is adjacent to, \emph{him}. And yet, the ungrammaticality of this example traces back to the availability of Case for \emph{him}, as can be appreciated by considering the fact that this example becomes grammatical without \emph{him}.
\begin{exe}
\ex \label{ex3:govill4} She made running easy.
\end{exe}
The ungrammaticality of (\ref{ex3:govill3}a) should arise because it violates the Case filter, and this indicates that \emph{make} cannot assign accusative to \emph{him}.

A salient difference between (\ref{ex3:govill3}a) and (\ref{ex3:govill3}b) is the depth of embedding of the DP to be assigned Case. There appears to be a condition on how far into its sister a verb can assign Case. The precise statement of this condition is still under investigation; we will consider some of the proposals in what follows. I will add this condition to the formulation of Case assignment in a way that allows us to modify it as we go along. An initial reformulation of the conditions on accusative Case assignment, then, is~(\ref{ex3:newcase2}).
\begin{exe}
\ex \label{ex3:newcase2} Accusative Case is assigned by some verbs and prepositions to positions that they c-command, are klose to, and are adjacent to.\\[7pt]
$\alpha$\ is \textsc{klose} to $\beta$\ iff there is no more than one phrase that dominates $\beta$\ but not $\alpha$.
\end{exe}
This definition allows a verb to assign Case to its sister and to the immediate daughters of its sister, but prevents it from assigning Case to any position more deeply in its sister. In \emph{Lectures on Government and Binding}, and much subsequent work, a term---``government''---is defined that collapses the \mbox{c-command} and klose relations. The condition on accusative Case assignment is often formulated as~(\ref{ex3:newcase3}).
\begin{exe}
\ex \label{ex3:newcase3} Accusative Case is assigned by some verbs and prepositions to positions that they govern and are adjacent to.\\[7pt]
$\alpha$\ \textsc{governs} $\beta$\ iff $\alpha$\ c-commands $\beta$\ and $\alpha$\ is klose to $\beta$.
\end{exe}

Finally, what we see from the derivations involving Argument Movement is that the Case filter will have to be satisfied if an S-structure meets its requirements but a D-structure doesn't. We saw in the Complex NP Shift examples that the Case filter could also be satisfied if its requirements were met by a D-structure but not an S-structure. It looks like the Case filter is satisfied if either a D-structure or an S-structure meet its requirements.
\begin{exe}
	\raggedright
	\ex \label{ex3:csf} Case filter\\
	An argument DP must be in a Case marked position at either D-structure or S-structure.
\end{exe}

With these changes, then, the system of Case assignment and the Case filter correctly work with the derived subjects hypothesis to ensure that subjects of finite clauses show up with nominative Case in the Specifier of IP at S-structure. The existence of small clauses, and the accompanying thesis that subjects of IPs start out in small clauses, requires that the subject \mbox{$\theta$-role} be assigned to the Specifier of the phrase headed by the term that has the \mbox{$\theta$-role}. This brings us to the second change to our system made necessary by the Derived Subject Hypothesis. The Projection Principle determines where \mbox{$\theta$-roles} are assigned, and so we should modify this condition to capture where subject \mbox{$\theta$-roles} are assigned. At present, the Projection Principle requires that words assign their \mbox{$\theta$-roles} to nearby positions within their projection. The Derived Subjects Hypothesis requires that this be broadened to include Specifier positions as well. 
\pagebreak
\begin{exe}
\ex \label{ex3:newproprin} The Projection Principle
\begin{xlist}
	\ex If $\beta$ is c-selected by $\alpha$, then there must be a projection of $\alpha$, $\pi$, that immediately dominates $\beta$ and if $\pi$ immediately dominates some other phrase, $\gamma$, then $\gamma$ must be another projection of $\alpha$ or part of another c-selected argument of~$\alpha$.
	\ex An external $\theta$-role of $\alpha$ must be assigned by $\alpha$ to a phrase in the Specifier position of~$\alpha$.
\end{xlist}
\end{exe}

This view of how $\theta$-roles are assigned makes most D-structures different than their S-structures. It requires that Argument Movement frequently be engaged to bring the subject argument of a verb from the position it gets its $\theta$-role from to the position that determines its Case. In section \ref{sec:argmove}, we'll examine more closely the rule that Argument Movement has in determining the form of sentences. But in the next section, we'll examine a consequence this hypothesis has for the D-structure status of~IPs.

\section{The EPP}
\label{sec:epp}

When we were discovering the phrase structure rules for English, we hit upon a couple oddities of IPs. One of those oddities is that the smallest IP is not a string that consists of just its head. This is a possibility for every other phrase we saw: the smallest NP is just a noun, the smallest VP is just a verb, and so on. But the smallest IP consists of both a phrase in Specifier position and a phrase in complement position, in addition to the head. We recorded this fact by giving to IPs the following phrase structure rules.
\begin{exe}
\raggedright
	\ex \label{ex3:IPrules}
	\begin{xlist}
		\ex IP $\rightarrow$ DP \ibar{I} 
		\ex IP $\rightarrow$ CP \ibar{I}
		\ex \ibar{I} $\rightarrow$ \obar{I} VP
	\end{xlist}
\end{exe}
Every IP, according to these rules, must have a DP or CP in its Specifier position, as well as a VP in its complement position.

If the derived subjects hypothesis is correct, this isn't correct. Instead, as the derivations we've seen in the previous section indicate, IPs have no Specifier at D-structure, and instead gain a Specifier by way of Argument Movement. We need D-structures to allow for the shape indicated in~(\ref{ex3:IPu}).
\begin{exe}
\raggedright
	\ex \label{ex3:IPu} IP $\rightarrow$ \ibar{I}
\end{exe}
This predicts that, in principle, an IP should be able to correspond to an S-structure that has nothing in its Specifier position. This, it turns out, is incorrect.

We have an account for this for every sentence that has a DP argument which does not reside in a Case-marked position at D-structure. In those scenarios, the Specifier of IP will, when it is finite, be the Case-marked position to which that argument DP must move. Thus, for instance, when a DP is an argument of an adjective like \emph{obvious}, it must move in order to satisfy the Case filter in the way indicated in~(\ref{ex3:obvious4}).
\begin{exe}
\raggedright
	\ex \label{ex3:obvious4}
	\small
	\begin{forest}
	[IP
		[DP [this idea, roof, name=binder] ]
		[\ibar{I}
			[\obar{I} [might]]
			[VP [\ibar{V}
				[\obar{V} [be]]
				[AP [\ibar{A}
					[\obar{A} [obvious]]
					[{\ }, name=trace ]
				] ]
			] ]
		 ]
	]
	\draw[->,>=stealth', thin, gray] (trace) to[out=south west,in=south,looseness=1] node[pos=0.75,xshift=-10mm]{} (binder);
	\end{forest}
\end{exe}

But what of IPs in which the D-structures do not violate the Case filter. If CP arguments are not subject to the Case filter, for instance, why is the D-structure shown in~(\ref{ex3:epp1}) not a grammatical S-structure.
\begin{exe}
\raggedright
	\ex \label{ex3:epp1}
	\small
	\begin{forest}
	[IP [\ibar{I}
		[\obar{I} [will]]
		[VP [\ibar{V}
			[\obar{V} [be]]
			[AP [\ibar{A} 
				[\obar{A} [obvious]]
				[CP [that pigs fly, roof] ]
			] ]
		] ]
	]]
	\end{forest}
\end{exe}
Instead, the D-structure in (\ref{ex3:epp1}) is associated with one of the two S-structures in~(\ref{ex3:epp2}).
\begin{exe}
\raggedright
	\ex \label{ex3:epp2}
\footnotesize
\begin{forest}
[IP
	[CP [that pigs fly, name=binder, roof] ]
	[\ibar{I}
		[\obar{I} [will]]
		[VP [\ibar{V}
			[\obar{V} [be]]
			[AP [\ibar{A}
				[\obar{A} [obvious]]
				[{\ }, name=trace ]
			] ]
		] ]
	 ]
]
\draw[->,>=stealth', thin, gray] (trace) to[out=south west,in=south,looseness=1] node[pos=0.75,xshift=-10mm]{} (binder);
\end{forest}
\begin{forest}
	[IP 
	[DP [it, roof]]
	[\ibar{I}
		[\obar{I} [will]]
		[VP [\ibar{V}
			[\obar{V} [be]]
			[AP [\ibar{A} 
				[\obar{A} [obvious]]
				[CP [that pigs fly, roof] ]
			] ]
		] ]
	]]
\end{forest}
\end{exe}
The pronoun \emph{it} in the second of these S-structures is not an argument. In fact, it doesn't refer to anything in the same way that an argument pronoun typically does. If it has a denotation, I don't know what it is. Terms that have no denotations are known as ``expletives,'' and the \emph{it} in (\ref{ex3:epp2}) is one.

In a syntax that embraces derivations, we see from these cases that the requirement that an IP have a DP or CP in its Specifier position is a fact about S-structures. IPs can exist without something in their Specifier position, but the S-structures of IPs cannot. We'll recognize this fact with the following description.
\begin{exe}
\raggedright
	\ex \label{ex3:eppform1}
	The S-structure of an IP must have a DP or CP in its Specifier position.
\end{exe}
We will investigate some of the ideas about what is responsible for this fact. \citet{Chomsky:1981} proposed that (\ref{ex3:eppform1}) is connected to the Projection Principle---the condition that determines where arguments must be at D-structure. He called (\ref{ex3:eppform1}) the ``Extension to the Projection Principle'' to signal this hypothesis. This proposal has been abandoned, and so the name no longer makes sense. Nonetheless, the name persists, though in the more cryptic acronym:~``EPP.''%
%
\marginnote{ The EPP is a specific version of what David Perlmutter and Paul Postal, in their Relational Grammar framework, call the “Final 1 Law.” See \citet{Perlmutter:1983}. The earliest statement of this condition that I know of is \citet{Perlmutter:1971}.}%
%

The EPP can be satisfied by using Argument movement to form an S-structure from a D-structure that does not put a CP or DP in Specifier of IP. This is what happens in (\ref{ex3:epp1}) and the first case of~(\ref{ex3:epp2}). It can also be satisfied by creating a D-structure in which the IP's Specifier is occupied by the expletive \emph{it}. This is what is illustrated by the second case in (\ref{ex3:epp3}), which we can take to be both the D-structure and the S-structure of that sentence. Forming this D-structure for a sentence in which the argument of the adjective is a DP, as in (\ref{ex:epp3}), leads to ungrammaticality because of the Case filter.
\begin{exe}
\raggedright
	\ex \label{ex3:epp3}
	\small
	\begin{forest}
	[IP
		[DP [it, roof] ]
		[\ibar{I}
			[\obar{I} [might]]
			[VP [\ibar{V}
				[\obar{V} [be]]
				[AP [\ibar{A}
					[\obar{A} [obvious]]
					[DP [this idea, roof] ]
				] ]
			] ]
		 ]
	]
	\end{forest}
\end{exe}
If there is only one Specifier position for any XP, then the sole Specifier position for the IP in (\ref{ex3:epp3}) is occupied by the expletive \emph{it}. Let's assume that this prevents Argument movement from forming an S-structure from (\ref{ex3:epp3}) in which \emph{this idea} occupies the Specifier of IP. There is consequently no derivation that includes (\ref{ex3:epp3}) in which the argument DP \emph{this idea} occupies a Case-marked position. This violates the Case filter.

In addition to the Case filter, we'll need to embrace the EPP. Together, they control the kinds of S-structures that are permitted.

\section{Argument Movement}
\label{sec:argmove}

The derived subjects hypothesis conjectures that subjects of garden variety finite IPs are spoken in positions different from their D-structure position. Their D-structure position is conjectured to be the Specifier of the VP belonging to that IP, a position that the existence of small clauses indicates is a possible position for subject arguments. It relies, then, on the existence of a mechanism that gives arguments a different position than their \mbox{$\theta$-marked} one, and we have introduced Argument Movement as that mechanism.

There is independent support for the existence of a mechanism like Argument Movement. There are situations in which we are more confident about where the relevant \mbox{$\theta$-marked} position is, and the argument bearing that \mbox{$\theta$-role} is spoken somewhere else. 

\subsection{Passive}
\label{sub:passive}

Consider, for instance, the pair of examples in (\ref{ex3:passintro}).%\enlargethispage*{21pt}
\begin{exe}
	\ex \label{ex3:passintro}
	\begin{xlist}
	\ex He will make [her angry at me].
	\ex She will be made [angry at me].
\end{xlist}
\end{exe}
The $\theta$-role that \emph{she} bears in (\ref{ex3:passintro}b) is the same one that it bears in (\ref{ex3:passintro}a), and yet \emph{she} shows up in a different Case and in a different position in (\ref{ex3:passintro}b). The VPs are different in these examples. The verb in (\ref{ex3:passintro}b) is said to be in the ``passive voice,'' while the verb in (\ref{ex3:passintro}a) is in the ``active voice.'' Passive verbs in English are embedded under the auxiliary verb \emph{be}. It is a fairly widespread generalization about Passive verbs that they do not assign Accusative Case, even if the Active verbs they are related to do. English conforms to this generalization. An informal way of thinking about the passive is that it is as an operation that derives a passive predicate from an active one by way of the changes described in~(\ref{ex3:passrule}).
\begin{exe}
	\ex \label{ex3:passrule} \textsc{Passive}
	\begin{xlist}
	\ex[]{Add passive morphology to the active verb and embed it under \emph{be}, and}
	\ex[]{Rob the verb of its Accusative Case, and}
	\ex[]{Rob the verb of its external $\theta$-role.}
	\end{xlist}							
\end{exe}
This takes a garden variety ``transitive'' verb and produces a verb which looks ``intransitive,'' as indicated in (\ref{ex3:passexam}).
\begin{exe}
	\ex \label{ex3:passexam} Jerry admires her. $\rightarrow$ She is admired.
\end{exe}
The D-structure representation of the Passive sentence in (\ref{ex3:passexam}) is (\ref{ex3:passexamt}a) from which the S-structure in (\ref{ex3:passexamt}b) is produced.%\enlargethispage{16pt}
\begin{exe}
\ex \label{ex3:passexamt} 
a.
\small
\begin{forest}
[IP [\ibar{I}
	[\obar{I} [pres] ]
	[VP [\ibar{V} 
		[\obar{V} [be] ]
		[VP [\ibar{V}
			[\obar{V} [admired] ]
			[DP [she, roof] ]
		] ]
	] ]
	]
]
\end{forest}
b.
\small
\begin{forest}
[IP
	[DP [she, roof]]
	[\ibar{I} 
		[\obar{I} [pres] ]
		[VP [\ibar{V}
			[\obar{V} [be] ]
			[VP [\ibar{V} [\obar{V} [admired]]] ]
		] ]
	]
]
\end{forest}
\end{exe}
Because Passive has robbed \emph{admired} of its ability to assign Accusative Case, its object will seek out the Nominative Case marked position to satisfy the Case filter. This will cause (\ref{ex3:passexamt}a) to be mapped onto an \mbox{S-structure} in which \emph{she} is positioned in the Specifier of IP.
\pagebreak

The passive in (\ref{ex3:passintro}b) has the derivation in~(\ref{ex3:passintrod}).
\begin{exe}
\ex \label{ex3:passintrod} 
\small 
\begin{forest}
[IP
	[DP [she, roof, name=binder]]
	[\ibar{I} 
		[\obar{I} [will]]
		[VP [\ibar{V}
			[\obar{V} [be] ]
			[VP [\ibar{V}
				[\obar{V} [made] ]
				[AP
					[{\ }, name=trace ]
					[\ibar{A}
						[\obar{A} [angry] ]
						[PP [at me, roof] ]
					 ]
				 ]
			] ]
		] ]
	]
]
\draw[->,>=stealth', ultra thin, gray] (trace) to[out=south west,in=south] node[pos=0.75,xshift=10mm]{} (binder);
\end{forest}
\end{exe}
This mirrors the kind of derivation that the derived subjects hypothesis gives to every finite clause. The subject argument is put into a D-structure position from which its \mbox{$\theta$-role} derives, but cannot remain there and satisfy the Case filter. Argument Movement kicks in and gives it a position in which it is Case marked. Passive Constructions, then, are ones in which we have independent support for displacement of arguments from their \mbox{$\theta$-marked} positions. They even manufacture cases which resemble the situation the derived subjects hypothesis claims to be widespread.

There's a confusing aspect to English morphology which clouds how the Passive works in English. English has a morphological rule by which transitive verbs can be converted into adjectives, and these adjectives look just like Passive participles. I'll call these adjectives, ``adjectival participles.'' This process is responsible for created the adjectival participle in (\ref{ex3:adj}), for instance.
\begin{exe}
	\ex \label{ex3:adj} the admired student  
\end{exe}
It could be, then, that what we have in (\ref{ex3:passexamt}b) is an adjectival phrase headed by \emph{admired} rather than a Passive VP. It might have a derivation like that in~(\ref{ex3:adjt}).
\begin{exe}
\ex \label{ex3:adjt} 
\small
\begin{forest}
[IP
	[DP [she, roof, name=binder]]
	[\ibar{I}
		[\obar{I} [pres] ]
		[VP [\ibar{V}
			[\obar{V} [be] ]
			[AP
				[{\ }, name=trace ]
				[\ibar{A} [admired, roof] ]
			 ]
		] ]
	 ]
]
\draw[->,>=stealth', ultra thin, gray] (trace) to[out=south west,in=south] node[pos=0.75,xshift=10mm]{} (binder);
\end{forest}
\end{exe}
Over a large class of cases, then, we will not be able to tell whether we are looking at a Passive sentence or a \emph{be}+AP (a so-called ``copular construction'') sentence.

In some cases, however, it is possible to tell.\marginnote[0pt]{The \emph{locus classical} for how to distinguish adjectives from verbal passives is \citet{Wasow:1977}.} This is because the morphological process that derives adjectives from verbs operates on the arguments of the verb. In particular, it makes the $\theta$-role that a verb assigns to the DP it c-selects become the external $\theta$-role of the resulting adjective. For instance, the adjective in (\ref{ex3:adjt}) assigns to its Specifier position the same $\theta$-role that the verb it was derived from assigns to its complement.\marginnote{See \citet{Levin:1986} for arguments for this view of the rule.} This is quite general. We can see this by considering the relation these adjectives have to the nouns they modify. The pre-nominal adjective in (\ref{ex3:adj}) is related to the noun it modifies in the same way that the verb it derives from is related to its object. When we figure out how the modification relation works semantically, we will need to derive the fact that the modification relation for adjectives roughly matches the relation these adjectival participles have to their subjects in the copular construction. We can use this correlation, then, to see which $\theta$-role is ``externalized'' by the adjective formation process. What we find when we look is that, in fact, it's always the ``direct object'' $\theta$-role that gets externalized. In each of the cases in (\ref{ex3:.adj}), for instance, the modified noun is related to its adjective in the same way that the direct object is to the verb.
\begin{exe}
	\ex \label{ex3:.adj}
	\begin{xlist}
	\ex[]{Ben introduced the man to Min-Joo.\\ the introduced man}
	\ex[]{Ben placed the book on the table.\\ the carefully placed book}
	\ex[]{Min-Joo baked the cake for Ben.\\ the baked cake}
	\end{xlist}  
\end{exe}
In no case is this relation like the one that holds between the verb and its subject argument or its ``indirect object'' argument.\footnote{For more details about this process, see \citet{Levin:1986}.}

Because of this restriction on the adjective formation process, we can be certain that the example in (\ref{ex3:passintro}b) is a Passive and not an adjective. This is because the argument that shows up as the nominative argument in this example is not an argument of the verb at all. It is the subject of the complement of the verb. If \emph{made} were an adjective in this example, then it's the complement itself that would have appeared as subject, as can be confirmed by considering alternations like those in (\ref{ex3:makepass}).
\begin{exe}
\ex \label{ex3:makepass} 
\begin{xlist}
	\ex[]{Someone made this table}
	\ex[]{this (recently) made table.\\[-3pt]}
	\ex[]{Someone made him sad.}
	\ex[*]{the (recently) made man  }
\end{xlist}
\end{exe}

What we see in the Passive, then, is that an argument DP has a different position than where its $\theta$-role is assigned. It's not that the Passive causes a $\theta$-role to be assigned by a verb to a different position, because sometimes the relocated argument DP isn't even an argument of the relevant verb. A syntax that embraces Argument movement can account for this.

\subsection{Raising}

There are verbs which behave much like the Passive verb in (\ref{ex3:passintro}b), but without needing to be Passivized. An example is (\ref{ex3:raiseintro}).
\begin{exe}
	\ex \label{ex3:raiseintro} She seems [angry at me].
\end{exe}
Just as in the Passive example, \emph{she} is an argument of \emph{angry at me}, not \emph{seems}. And yet, just as in the Passive, the subject argument of \emph{angry at me} shows up in a higher Specifier of IP and bears the nominative Case assigned to this position. This example seems to require the derivation in~(\ref{ex3:raised}).
\begin{exe}
\ex \label{ex3:raised} 
\small
\begin{forest}
[IP
	[DP [she, roof, name=binder]]
	[\ibar{I} 
		[\obar{I} [present] ]
		[VP [\ibar{V}
			[\obar{V} [seem] ]
			[AP
				[{\ }, name=trace ]
				[\ibar{A}
					[\obar{A} [angry] ]
					[PP [at me, roof] ]
				 ]
			 ]
		] ]
	]
]
\draw[->,>=stealth', ultra thin, gray] (trace) to[out=south west,in=south] node[pos=0.75,xshift=10mm]{} (binder);
\end{forest}
\end{exe}
Verbs like \emph{seem} are called ``raising'' verbs. They have properties that force an argument in their complement to ``raise'' into their subject position. Under the proposals we are examining, one of those properties is that it does not assign accusative Case. Raising verbs, then, are like \emph{be} in being among those verbs that do not have an accusative Case associated with them. As a result, in situations like (\ref{ex3:raised}) the argument of the embedded small clause must Argument Move into a position where it may satisfy the Case filter. In these circumstances too we see cases like those that the derived subjects hypothesis claims are common. 

\section{Quantifier Float}\label{sec:qfloat}

The derived subjects hypothesis is, presently, a kind of null hypothesis. It rests on the observation that the existence of small clauses indicates that external \mbox{$\theta$-roles} can be assigned to Specifier of VP, and so the null hypothesis is that they always are. The existence of movement Transformations, or some other means of overcoming violations of Contiguity (i.e., resolving mismatches between constituent structure and linear order), makes it understandable why the surface position of the subject of finite clauses is not the Specifier of VP. And the Case filter provides a reason why a movement transformation would be invoked in these contexts: it gives an account of why subject arguments remain in Specifier of VP sometimes, and move in others. The derived subjects hypothesis, then, fits together independently needed pieces---movement transformations, the Case filter, and a Projection Principle that produces small clauses---into a simple picture of how subject arguments are placed. Nonetheless, it would be encouraging to find empirical support that directly speaks on behalf of the derived subjects hypothesis.

\citet{Sportiche:1988} argues that the phenomenon of ``Quantifier Float'' provides such evidence. Quantifier Float arises when a quantificational term is not in the determiner position normally reserved for quantifiers, but is instead to the right of DP it is associated with. It can be illustrated by the English and French examples in (\ref{ex3:qfloat1}).\footnote{\citet{Miyagawa:1989} provides a similar analysis for Quantifier Float based on Japanese data.}
\begin{exe}
	\ex \label{ex3:qfloat1} \gll Les enfants (*tous) ont (tous) vu ce film (tous). \\ The kids (all) have (all) seen this film (*all).\\
\end{exe}
Sportiche argues that there is a violation of Contiguity in these examples of the sort that movement Transformations are designed to account for. He suggests that the floated quantifier and the DP it is related to start out together at D-structure and then are separated at S-structure by Argument Movement. Note  that there are language particular differences in this process that are reflected in the differing grammaticality judgments for some of the French and English examples. We'll return to these differences in the next chapter.

Sportiche's argument relies on the observation that the relation between the Floated Quantifier and the argument it is related to mimics the relation that Argument movement is designed for. Once this is appreciated, the fact that the floated quantifier appears in just those positions the derived subjects hypothesis predicts it should provides the evidence we are looking for. To see the force of his argument, we will first have to learn some generalizations about Argument Movement. 

One of those generalizations is that the moved phrase can only be relocated to a c-commanding position. I'll call this condition ``Upwards.''
\begin{exe}
\ex \label{ex3:ccommgen} Upwards\\
If $\alpha$\ moves to the position $\beta$, then $\beta$\ must c-command $\alpha$.
\end{exe}
This is a generalization about all movement rules: all clear cases of a term moving to a non-c-commanding position are ungrammatical.%
%
\marginnote[0pt]{We'll encounter examples in the next chapter which requires a reformulation of Upwards, however. The version in (\ref{ex3:ccommgen}) isn't entirely accurate.}%
%
 For instance, there is another movement rule in English that moves certain DPs into Specifier of CP. This rule is part of a suite of operations that, among other things, is responsible for making question sentences in English. Some examples of this rule's operation are in~(\ref{ex3:whmove1}).
\begin{exe}
	\raggedright
	\ex \label{ex3:whmove1} 
	\begin{xlist}
		\ex I wonder who you talked to.
		\ex I asked which problem you will talk about.
		\ex I asked about which problem you will talk.
	\end{xlist}
\end{exe}
The DPs that this rule moves are called ``wh-phrases,'' and an initial, though not complete, definition of wh-phrases is~(\ref{ex3:whphr}).
\begin{exe}
	\raggedright
	\ex \label{ex3:whphr} A wh-phrase is an XP that:
	\begin{xlist}
		\ex is headed by a wh-determiner (e.g., \emph{which})
		\ex is headed by a wh-pronoun (e.g., \emph{who}, \emph{what})
		\ex is a PP with a wh-phrase in complement position.
	\end{xlist}
\end{exe}
We will look closer at the details of this rule later, but an initial formulation of it is in~(\ref{ex3:whmove}).
\begin{exe}
	\raggedright
	\ex \label{ex3:whmove} Wh Movement\\
	Move a wh-phrase to a Specifier of CP.
\end{exe}
\pagebreak
This produces a derivation for (\ref{ex3:whmove1}a) like that shown in~(\ref{ex3:whmoved}).\marginnote[0pt]{One of the peculiarities of CPs that make questions in English is that there is no obvious word in \obar{C} position. We will return to this mystery in a couple of chapters.}
\begin{fullwidth}
\medskip
\begin{exe}
	\raggedright
	\ex \label{ex3:whmoved} I asked...
	\sn
	\footnotesize
	\begin{forest}
	[CP
		[\obar{C} ]
		[IP 
			[DP [you, roof] ]
			[\ibar{I}
				[\obar{I} [past] ]
				[VP [\ibar{V}
					[\obar{V} [talk] ]
					[PP [\ibar{P}
						[\obar{P} [to] ]
						[DP [who, roof] ]
					] ]
				] ]
			 ]
		]
	]
	\end{forest}
	{\Huge ,}
	\small
	\begin{forest}
	[CP
		[DP [who, roof]]
		[\ibar{C} 
			[\obar{C} ]
			[IP 
				[DP [you, roof]]
				[\ibar{I}
					[\obar{I} [past] ]
					[VP [\ibar{V}
						[\obar{V} [talk] ]
						[PP [\ibar{P} [\obar{P} [to]]] ]
					] ]
				 ]
			 ]
		]
	]
	\end{forest}
	\end{exe}
	\medskip
	\end{fullwidth}

Now, to see that Wh Movement is also subject to (\ref{ex3:ccommgen}), consider the sentence in (\ref{ex3:nonccom1}). (\ref{ex3:nonccom1}) is ungrammatical because the rule of Wh Movement has moved the wh-phrase to a Specifier of CP that does not c-command the position the wh-phrase occupies at D-structure.
\begin{exe}
\ex[*]{ \label{ex3:nonccom1} He explained to whom Ethan left.}\marginnote[0pt]{I have placed the CP complement to \emph{explain} inside the VP in (\ref{ex3:nonccom1}), but we have now concluded that CPs can be Complex NP Shifted to the right of VPs. I leave out this detail here just to simplify the derivation involved. This example would still violate (\ref{ex3:ccommgen}) even if I'd given the correct S-structure parse.}
\sn
\small
\begin{forest}
[IP
	[DP [he, roof]]
	[\ibar{I} 
		[\obar{I} [past] ]
		[VP [\ibar{V} 
			[\ibar{V}
				[\obar{V} [explain] ]
				[{\ }, name=trace ]
			 ]
			[CP
				[PP [to whom, roof, name=binder]]
				[\ibar{C}
					[\obar{C} ]
					[IP [Ethan left, roof]]
				 ]
			 ]
		] ]
	]
]
\draw[->,>=stealth', ultra thin, gray] (trace) to[out=west,in=south west] node[pos=0.75,xshift=10mm]{} (binder);
\end{forest}
\end{exe}
\pagebreak
Complex NP Shift also obeys (\ref{ex3:ccommgen}), as the ungrammaticality of (\ref{ex3:nonccom2}) demonstrates.
\begin{exe}
\ex[*]{\label{ex3:nonccom2} That Ethan stood every silly boy bothered.}
\sn
\small
\begin{forest}
[IP
	[CP	
		[\obar{C} [that] ]
		[IP
			[IP
				[DP [Ethan, roof] ]
				[\ibar{I}
					[\obar{I} [past] ]
					[VP [stand, roof] ]
				 ]
			 ]
			[DP [every silly boy, roof, name=binder] ]
		 ]
	 ]
	[\ibar{I} 
		[\obar{I} [past] ]
		[VP [\ibar{V}
			[\obar{V} [bother] ]
			[{\ }, name=trace ]
		] ]
	]
]
\draw[->,>=stealth', ultra thin, gray] (trace) to[out=south,in=south, looseness=2] node[pos=0.75,xshift=10mm]{}  (binder);
\end{forest}
\end{exe}
These are a small sample of the many examples where movement rules would create ungrammatical sentences, if they were not subject to~(\ref{ex3:ccommgen}).

Sportiche points out that Quantifier Float is subject to the parallel constraint in (\ref{ex3:floatgen}). 
\begin{exe}
\ex \label{ex3:floatgen} If $\alpha$\ is a floated quantifier related to the DP $\beta$, then $\beta$\ must c-command~$\alpha$.
\end{exe}
It isn't possible for (\ref{ex3:27}b) to have an interpretation like that which (\ref{ex3:27}a) gets, presumably because it is not possible for Quantifier Float to relate these two structures. And because (\ref{ex3:27}b) can have no other interpretation---there are no other plural DPs for \emph{tous} to be construed with---the sentence is ungrammatical.
\begin{exe}
	\ex \label{ex3:27} 
	\begin{xlist}
	\ex[]{\gll [L'auteur de tous ces livres] a vu ce film. \\ {the author} of all these books has seen this film\\ }
	\ex[*]{\gll [L'auteur de ces livres] a tous vu ce film.\\ {the author} of these books has all seen this film\\ }
	\end{xlist}
\end{exe}
The ungrammaticality of (\ref{ex3:27}b), then, can be credited to the fact \emph{ces livres} fails to c-command \emph{tous}.

What the c-command requirement on Quantifier Float establishes is that however the semantic connection between a floated quantifier and the DP it involves is worked out, it should follow that the DP c-commands the quantifier. There are a variety of ways that could be ensured, but because movement Transformations obey this relation, they provide one. We'll see what Sportiche's suggestion is about the semantic connection between the floated quantifier and the DP in a moment. But first, let's consider how Quantifier Float seems to be diagnostic of Argument movement.

We have seen in the previous section that the subject argument, \emph{the kids}, in examples like (\ref{ex3:raiseqfl}) moves from a position within the small clause that follows the main verb.
\begin{exe}
\ex \label{ex3:raiseqfl} 
\begin{xlist}
	\ex[]{The kids seem happy.}
	\ex[]{The kids were made happy.}
	\ex[]{The kids were on the table.}
\end{xlist}
\end{exe}
And, as expected on Sportiche's proposal, a floated quantifier within the small clause can be related to \emph{the boys}, as in (\ref{ex3:raiseqfl2}).\marginnote[0pt]{Some speakers find (\ref{ex3:raiseqfl2}b) marginal; I don't know why.}
\begin{exe}
\ex \label{ex3:raiseqfl2} 
\begin{xlist}
	\ex[]{The kids seem both happy.}
	\ex[?]{The kids were made both happy.}
	\ex[]{The kids were both on the table.}
\end{xlist}
\end{exe}
The relationship between the floated quantifier and the subject is like that in examples such as (\ref{ex3:raiseqfl3}), then.
\begin{exe}
\ex \label{ex3:raiseqfl3} 
\begin{xlist}
	\ex[]{She made the kids both happy.}
	\ex[]{She found the kids both on the table.}
\end{xlist}
\end{exe}
In (\ref{ex3:raiseqfl3}), unlike (\ref{ex3:raiseqfl2}), the DP is able to remain within the small clause and satisfy the Case filter.

Consider now (\ref{ex3:qfloatlong}).
\begin{exe}
\ex \label{ex3:qfloatlong} The parents found the kids both on the table.
\end{exe}
This is grammatical, and has an interpretation in which \emph{both} is construed with \emph{the kids}. It doesn't have an interpretation, however, in which \emph{both} is construed with \emph{the parents}. Unlike (\ref{ex3:raiseqfl2}), the floated quantifier in the small clause of this example cannot be construed with the argument in the higher subject position. The difference between (\ref{ex3:raiseqfl2}) and (\ref{ex3:qfloatlong}) is that the argument in the higher subject position is \mbox{$\theta$-marked} inside the small clause in (\ref{ex3:raiseqfl2}), but in (\ref{ex3:qfloatlong}) it is \mbox{$\theta$-marked} outside the small clause. In fact, (\ref{ex3:qfloatlong}) is part of the following generalization about the relationship between a floated quantifier and the DP it is construed with.
\begin{exe}
\ex \label{ex3:qfloatgen} A floated quantifier construed with a DP cannot be lower in the phrase marker than the position that DP gets its $\theta$-role.
\end{exe}
Because \emph{the parents} in (\ref{ex3:qfloatlong}) is \mbox{$\theta$-marked} by some phrase that contains \emph{find}, the floated quantifier is necessarily contained within this projection and (\ref{ex3:qfloatgen}) prevents it from being construed with \emph{the parents}.  Sportiche's proposal that Argument Movement relates a floated quantifier's position to the argument it is construed with derives (\ref{ex3:qfloatgen}). The phrases that Argument Movement will apply to must start out in their \mbox{$\theta$-marked} positions; this is a consequence of the Theta Criterion. Therefore, because Argument movement brings arguments to c-commanding positions only, there will be no way for an argument to be placed within the XP that \mbox{$\theta$-marks} it. 

It's clear, now, why examples like (\ref{ex3:qfloat1}) provide evidence for the derived subjects hypothesis.
\begin{exe}
\exr{ex3:qfloat1} The kids have all seen this film.
\end{exe}
For (\ref{ex3:qfloat1}) to conform to (\ref{ex3:qfloatgen}), the position that \emph{the kids} gets its \mbox{$\theta$-role} from must be within the phrase containing \emph{have}. That is precisely what the derived subjects hypothesis claims.

Let's consider now how Sportiche proposes to make the semantic connection between the floated quantifier and the DP it is construed with. His leading idea is to use whatever semantic rules of interpretation derive a meaning from (\ref{ex3:qflsem1}) to give what appears to be the same meaning for (\ref{ex3:qfloat1}).
\begin{exe}
\ex \label{ex3:qflsem1} All the kids have seen this film.
\end{exe}
Imagine that the semantic connection between \emph{all} and the DP it combines with in (\ref{ex3:qflsem1}) is established at D-structure. This connection, then, is produced from the representation in (\ref{ex3:qflsem1d}).
\begin{exe}
\ex \label{ex3:qflsem1d}
\small
\begin{forest}
[\ibar{I}
	[\obar{I} [pres] ]
	[VP [\ibar{V}
		[\obar{V} [have] ]
		[VP 
			[QP
				[\obar{Q} [all] ]
				[DP [the kids, roof] ]
			 ]
			[\ibar{V}
				[\obar{V} [seen] ]
				[DP [this film, roof] ]
			 ]
		]
	] ]
]
\end{forest}
\end{exe}
This D-structure yields the S-structure representation for (\ref{ex3:qflsem1}) when Argument movement relocates the entire QP \emph{all the kids} into the Specifier of IP. It yields the S-structure for (\ref{ex3:qfloat1}) when Argument movement applies to just the DP part of this QP, leaving the quantifier behind, as in (\ref{ex3:qfloat1d}).
\begin{exe}
\ex \label{ex3:qfloat1d} 
\footnotesize
\begin{forest}
[IP
	[DP [the kids, roof, name=binder]]
	[\ibar{I} 
		[\obar{I} [pres] ]
		[VP [\ibar{V}
			[\obar{V} [have] ]
			[VP 
				[QP
					[\obar{Q} [all] ]
					[{\ }, name=trace ]
				 ]
				[\ibar{V}
					[\obar{V} [seen] ]
					[DP [this film, roof] ]
				 ]
			 ]
		] ]
	]
]
\draw[->,>=stealth', ultra thin, gray] (trace) to[out=south,in=south, looseness=2.2] node[pos=0.75,xshift=10mm]{} (binder);
\end{forest}
\end{exe}

There's a popular alternative to Sportiche's account which you should keep in mind. I'll introduce it now, and offer a problem for it, but keep it at hand because it may begin to look more appealing when we discover some of the problems for Sportiche's analysis. This alternative collapses the conditions on so-called ``subject oriented'' (see \citealt{Jackendoff:1972}) adverbs with those of Q-float. Note that these adverbs have a distribution very much like floated quantifiers.
\begin{fullwidth}
\begin{exe}
	\vspace*{\baselineskip}
	\ex \label{ex3:44} The kids (deliberately) have (deliberately) seen (*deliberately) this film (*?deliberately).
	\vspace*{\baselineskip}
\end{exe}
	\end{fullwidth}
To get the interpretation right, we would need a semantics that lets a floated quantifier `distribute' the predicate that follows it over the parts of the subject. This approach to Q-float has many advocates: \citet{Bowers:1993}, for example, and \citet{Kayne:1975}. \citet[chapter 3]{Roberts:1987} discusses the semantics for such an analysis. This alternative can be wedded to Sportiche's proposal that Argument movement is involved; it can be seen as a different way of producing the semantic composition. If we are to maintain Sportiche's suggestion that Argument movement is implicated, we should see the D-structure representation for (\ref{ex3:qfloat1}) as being something like (\ref{ex3:qfloat1da}).
\begin{exe}
\ex \label{ex3:qfloat1da} 
\small
\begin{forest}
[\ibar{I}
	[\obar{I} [pres]]
	[VP [\ibar{V}
		[\obar{V} [have] ]
		[VP
			[DP [the kids, roof]]
			[\ibar{V}
				[QP [all, roof] ]
				[\ibar{V}
					[\obar{V} [seen] ]
					[DP [this film, roof] ]
				 ]
			 ]
		 ]
	] ]
]
\end{forest}
\end{exe}
On this alternative view, Argument movement could still be credited with determining the relationship between an argument DP and a floated quantifier that Sportiche's proposal offers. What we might consider altering is how the quantifier combines syntactically with the DP it is construed with. Sportiche's (best) argument against this alternative is that the distribution of subject-oriented adverbs and floated Q are not always the same cross-linguistically.

There's another kind of argument against this alternative found in \citet{Giusti:1990, Giusti:1990b} which is worth looking at. Her argument is typological, and based on where floated quantifiers can be found related to objects across languages. First, note that floated quantifiers related to objects are, by themselves, something of a puzzle for an approach like the alternative sketched above. That account demands that the floated quantifier combine with a VP, or similar predicate, and steer how that predicate combines semantically with the neighboring DP. The existence of floated quantifiers related to objects, then, is by itself an argument for Sportiche's imagine of the underlying connection between quantifier and DP. What Giusti observes is that the languages which allow quantifiers to be related to objects are those for which we have independent evidence that objects can move.

Thus, German/Dutch and Icelandic, but not English or Danish, can have a floated quantifier related to an object, as shown in~(\ref{ex3:45}) and~(\ref{ex3:46}).
\begin{exe}
	\ex \label{ex3:45} 
	\begin{xlist}
	\ex[]{\gll Der Lehrer hat die Sch\"uler alle gelobt. \\ the teacher has the students all praised\\ `The teacher has praised all the students.' \hspace*{\fill} (German)}
\ex[]{\gll De leraar heeft de kinderen allen geloofd\\ The teacher has the children all praised.\\ `The teacher has praised all the students.' \hspace*{\fill} (Dutch)}
	\ex[]{\gll St\'udentarnir lasu greinina allir. \\ students-the read article-the all\\ `The students read all the articles. \hspace*{\fill} (Icelandic)}
	\end{xlist}
\hspace*{\fill} \citep[(1): 137, (22a) \& (25b): 144]{Giusti:1990b}\hspace*{.5in}
	\ex \label{ex3:46}
	\begin{xlist}
	\ex[*]{The teacher has praised the children all.}
	\ex[*]{\gll Laereren roste eleverne alle \\ teacher-the praised children-the all\\ `The teacher praised all the students.' \hspace*{\fill} (Danish)}
	\end{xlist}
\hspace*{\fill} \citep[(2c): 137]{Giusti:1990b}\hspace*{.5in}
\end{exe}
She argues that this can be related to the fact that definite objects in Dutch, German and Icelandic can move leftwards, whereas this is not possible for objects in English and Danish. That leftwards movement is possible in these languages is, at least partly, indicated by the fact that they can precede adverbs which are normally found at the left edge of VPs in these languages. We have contrasts, then, between examples such as (\ref{ex3:47}) and those in~(\ref{ex3:48}).%\enlargethispage*{16pt}
\begin{exe}
	\ex \label{ex3:47} 
	\begin{xlist}
	\ex[]{\gll Der Lehrer hat die Sch\"uler ohne Zweifel gelobt. \\ the teacher has the students without {a doubt} praised.\\ `The teacher has without a doubt praised the students.' \hspace*{\fill} \parbox{.5in}{(German)}}
	\ex[]{\gll De leraar heeft de kinderen gisteren geloofd.\\ the teacher has the children yesterday praised.\\`The teacher has yesterday praised the students.' \hspace*{\fill} (Dutch)}
	\ex[]{\gll St\'udentarnir lasu greinina ekki.\\ students-the read articles-the not\\ `The students didn't read the articles.' \hspace*{\fill} (Icelandic)}
	\end{xlist}
\hspace*{\fill} \citep[(2): 136, (22a): 144, (13): 141]{Giusti:1990b}\hspace*{.5in}
\pagebreak
	\ex \label{ex3:48}
	\begin{xlist}
	\ex[*]{The teacher has praised the children not.}
	\ex[*]{\gll Laereren roste eleverne uden tvivl.\\ teacher-the praised students-the without {a doubt}.\\ `The teacher without a doubt praised the students.' \hspace*{\fill} (Danish)}
	\end{xlist}
\hspace*{\fill} \citep[(2): 137]{Giusti:1990b}\hspace*{1in}
\end{exe}
This is just what we would expect on a Argument Movement account of Quantifier float, of course, because on that account movement is what is responsible for separating a DP from the quantifier it is construed with. But this is also an argument for Sportiche's particular way of putting the quantifier together with the DP. On Sportiche's proposal, the only way a quantifier can appear after a DP it is construed with is by that DP's movement. The alternative account, by contrast, allows the D-structure arrangement of quantifier and DP to put the quantifier in a position that follows the DP. It's not clear how such an account could capture Giusti's correlation.

On the other hand, there are examples which Sportiche's proposal does not explain. As he notes, his model of Quantifier Float should allow quantifiers to be related to objects in English in just those cases where the object has undergone Argument movement. We've seen that one situation where this happens is in the Passive: an object moves into a higher Case-marked position when the verb it is an object of is Passivized. And yet, Quantifier Float is not possible in these contexts, as the ungrammaticality of (\ref{ex3:badq}) indicates.
\begin{exe}
\ex[*]{\label{ex3:badq} The kids were praised all.}
\end{exe}
For some steps towards a solution, see \citet{Boskovic:2001,Boskovic:2004a}. 

Another problem concerns examples like (\ref{ex3:raiseqfl3}).
\begin{exe}
\exr{ex3:raiseqfl3} She made the kids both happy.
\end{exe}
In this example, \emph{the kids} does not undergo movement. Its D-structure position is also the one where it gets assigned Case. If quantifiers follow the DPs they are construed with only when they move, how is it that the quantifier manages to follow the DP in this example? We will see a solution to this problem in Chapter 6. Keeping these problems in mind, let's nonetheless adopt Sportiche's account of Quantifier Float.\marginnote[0pt]{For other work on Quantifier Float, see \citet{Jaeggli:1982}, \citet{Belletti:1982}, \citet{Deprez:1989}, \citet{Doetjes:1995} and \citet{Bobaljik:2003}.}

One interesting consequence of Sportiche's analysis is that floated quantifiers can be seen as indicators of the positions that arguments occupy not only in the {D-structure} element of the derivation, but in all the {non-S-structure} parses of the derivation. It gives us, in other words, a useful glimpse at these otherwise invisible parses. We learn from this that phrases can take smaller steps than necessary in their journey from a  $\theta$-marked position to a Case marked position. So in (\ref{ex3:multistep}), for example, it appears that \emph{the children} has made an intermediary stop in the Specifier of the VP headed by \emph{have}.
\begin{exe}
	\ex \label{ex3:multistep} The children might all have eaten chocolate.
\end{exe}
Argument Movement is employed twice, as indicated in (\ref{ex3:multistepd}).
\begin{exe}
\ex \label{ex3:multistepd} 
\small
\begin{forest}
[IP
	[DP [the children, roof, name=binder2]]
	[\ibar{I} 
		[\obar{I} [might] ]
		[VP
			[QP , name=binder1
				[\obar{Q} [all] ]
				[{\ }, name=trace2 ]
			]
			[\ibar{V}
				[\obar{V} [have] ]
				[VP 
					[{\ }, name=trace1 ]
					[\ibar{V}
						[\obar{V} [eaten] ]
						[DP [chocolate, roof] ]
					 ]
				 ]
			 ]
		 ]
	]
]
\draw[->,>=stealth', ultra thin, gray] (trace2) to[out=south,in=south,looseness=2] node[pos=0.75,xshift=10mm]{} (binder2);
\draw[->,>=stealth', ultra thin, gray] (trace1) to[out=south west,in=south,looseness=1.5] node[pos=0.75,xshift=10mm]{} (binder1);
\end{forest}
\end{exe}
We learn from these sorts of examples that it's possible for Argument movement to act on a structure that is itself the product of Argument movement. Such derivations are said to involve ``successive cyclic'' movement. The existence of such derivations means that complex interactions of movements are conceivable, and, as a result, the mapping from D-structure to S-structure is potentially quite complex. There are problems that lurk here, but they are for another time.

\section{Infinitives}
\label{sec:infinitives}
% //todo This section should have a discussion of the condition on A movement that blocks movement out of CP. Presently, there is only the condition that prevents A movement from a Case marked position. Does this derive the condition that prevents movement out of CP?

We have seen that ``raising verbs'' select small clauses but do not assign Case to the DP argument within that small clause. As a result, Argument movement brings that DP argument into a Case marked position. In (\ref{ex3:raisenex}), for instance, \emph{Helen} has moved from within the small clause where its \mbox{$\theta$-role} is assigned, into its surface, nominative Case marked, position.
\begin{exe}
\ex \label{ex3:raisenex} Helen seems happy.
\end{exe}
Raising verbs are (some of them) also capable of selecting an IP headed by \emph{to}, as in (\ref{ex3:raiseipex}).
\begin{exe}
\ex \label{ex3:raiseipex} Helen seems to be happy
\end{exe}
In this example too, \emph{Helen} has moved from a position within the small clause into its surface, nominative Case marked, position.

But in this case, there is a question about how the EPP is satisfied.%
%
\marginnote[0pt]{The EPP, recall, is the requirement that a Specifier of IP be filled with a DP or~CP.}%
%
It's clear how the highest IP satisfies the EPP: its Specifier is filled with the subject argument that resides in it at S-structure. But what of the Specifier of the IP headed by \emph{to}? One possibility is that (\ref{ex3:raiseipex}) has a successive cyclic derivation like that in~(\ref{ex3:dereppt}a).
\begin{exe}
\ex \label{ex3:dereppt} 
a. \hspace*{-10pt}
\small
\begin{forest}
[IP
	[\obar{I} [pres]]
	[VP [\ibar{V}
		[\obar{V} [seem] ]
		[IP
			[DP [Jason, roof, name=binder]]
			[\ibar{I}
				[\obar{I} [to] ]
				[VP [\ibar{V}
					[\obar{V} [be] ]
					[AP
						[{\ }, name=trace ]
						[\ibar{A} [happy, roof] ]
					 ]
				] ]
			 ]
		 ]
	]]
]
\draw[->,>=stealth', ultra thin, gray] (trace) to[out=south west,in=south] node[pos=0.75,xshift=10mm]{} (binder);
\end{forest}
\hspace*{-20pt}
b.
\small
\begin{forest}
[IP
	[DP [the girls, roof, name=binder]]
	[\ibar{I} 
		[\obar{I} [pres]]
		[VP [\ibar{V}
			[\obar{V} [seem] ]
			[IP
				[QP, name=binder1
					[\obar{Q} [all] ]
					[{\ }, name=trace ]
				]
					[\ibar{I}
						[\obar{I} [to] ]
						[VP 
							[{\ }, name=trace1 ]
							[\ibar{V}
								[\obar{V} [like] ]
								[DP [chocolate, roof] ]
							 ]
						 ]
					]
				 ]
			 ]
		] ]
	]
]
\draw[->,>=stealth', ultra thin, gray] (trace) to[out=south,in=south,looseness=1.7] node[pos=0.75,xshift=10mm]{} (binder);
\draw[->,>=stealth', ultra thin, gray] (trace1) to[out=south west,in=south] node[pos=0.75,xshift=10mm]{} (binder1);
\end{forest}
\end{exe}
From this representation, A-movement could apply again to form the \mbox{S-structure} we see in (\ref{ex3:raiseipex}). We might imagine, then, that (\ref{ex3:raiseipex}) has a successive cyclic derivation. If this is correct, we should understand the EPP to be a condition that can be satisfied by any phrase marker in a derivation.

That a successive cyclic derivation of this sort is possible is indicated by (\ref{ex3:multistep2}), whose derivation is indicated in (\ref{ex3:dereppt}b).
\begin{exe}
	\ex \label{ex3:multistep2} The girls seem all to like chocolate.
\end{exe}
We see in this example that Argument movement has moved \emph{all the girls} into the Specifier of the infinitive before the DP part is moved into its surface position. The presence of the floated quantifier in the Specifier of the infinitival IP is the tell-tale sign that the subject argument has passed through the Specifier of the infinitival clause.

Infinitival clauses come in a variety of types that can be distinguished from those that follow raising verbs, and we will look at some of these in this section. The example in (\ref{ex3:con1}), for example, though it has a superficial similarity to (\ref{ex3:raiseipex}), is different in some ways.
\begin{exe}
	\ex \label{ex3:con1} She tried [to be happy]
\end{exe}
In (\ref{ex3:raiseipex}), the count of  $\theta$-roles and arguments matched, as expected. But in (\ref{ex3:con1}), this isn't the case. Both \emph{tried} and \emph{happy} assign an external  $\theta$-role, and yet there is only one argument which seems to bear them both: \emph{she}. This is precisely what the Theta Criterion prohibits.  Recall, that the Theta Criterion is designed in this way in an attempt to capture the fact that (\ref{ex3:8}a) doesn't have the meaning that (\ref{ex3:8}b) does.\marginnote[0pt]{How the Theta Criterion applies to cases where an argument moves from one $\theta$-marked position to another $\theta$-marked position shouldn't be clear, since the Theta Criterion is a condition that holds just of D-structures. We come back to look at this more closely in a later section.}
% \enlargethispage*{\baselineskip}
\begin{exe}
	\ex \label{ex3:8} 
	\begin{xlist}
	\ex[]{I showed Chris.}
	\ex[]{I showed Chris himself.}
	\end{xlist}
\end{exe}
So, either we should find another explanation for the fact that (\ref{ex3:8}) illustrates, or we should find a way of resolving (\ref{ex3:con1}) with the Theta Criterion.

Note also that (\ref{ex3:con1}) seems to be a violation of the Extension to the Projection Principle. Unless we determine that (\ref{ex3:con1}) involves Argument movement of \emph{she} through the Specifier of the infinitival IP, as in the raising case, there appears to be nothing that occupies this Specifier position at any point in the derivation. Moreover, it's not just that (\ref{ex3:con1}) seems to counterexemplify these principles, it perversely can't satisfy them. If the Specifier of the embedded IP is given content, and in so doing hosts an argument for the  $\theta$-role that \emph{angry} assigns, the result is ungrammatical.
\begin{exe}
	\ex[*]{ \label{ex3:nconsbj1} She tried [him to be angry].}
\end{exe}
And if a predicate that does not have an external \mbox{$\theta$-role} to assign is put into these infinitives, the result is also ungrammatical.
\begin{exe}
\ex[*]{ \label{ex3:nthsbj1} She tried [ to seem that pigs fly].}
\end{exe}

Infinitives with this peculiar mix of properties are called control infinitives, and a wide assortment of verbs select infinitives of this type. Some are in~(\ref{ex3:conlist}).
\begin{exe}
	\ex \label{ex3:conlist}
	\begin{xlist}
	\ex Jennifer prefers to eat chocolate.
	\ex Jennifer managed to eat chocolate.
	\ex Jennifer attempted to eat chocolate.
	\ex Jennifer wished to eat chocolate.
	\ex Jennifer asked to eat chocolate.
	\ex Jennifer proposed to eat chocolate.
	\ex Jennifer neglected to eat chocolate.
	\ex Jennifer promised to eat chocolate.
	\ex Jennifer decided to eat chocolate.
\end{xlist}
\end{exe}
In all of these examples, there is an argument---here it is \emph{Jennifer}, the subject of the higher verb---that appears to bear the \mbox{$\theta$-role} assigned to the subject of the infinitival clause. We say in these cases that the subject argument ``controls'' the infinitival clause, or ``controls'' the embedded subject. There are instances of control infinitives where it's the object of the higher verb that is the controller; (\ref{ex3:conlist2}) provides some examples.
\begin{exe}
\ex \label{ex3:conlist2} 
\begin{xlist}
	\ex Bernhard told Cecilia to eat chocolate.
	\ex Bernhard asked Cecilia to eat chocolate.
	\ex Bernhard required Cecilia to eat chocolate.
	\ex Bernhard forced Cecilia to eat chocolate.
\end{xlist}
\end{exe}
In these examples, it is the object, \emph{Ceclia}, not the subject, \emph{Bernhard}, that is understood to bear the external \mbox{$\theta$-role} assigned by \emph{eat}.

From what we have seen about control infinitives, we can draw the following conclusions.
\begin{exe}
\ex \label{ex3:contrdesc1} 
\begin{xlist}
	\ex A control infinitive must have a subject argument.
	\ex The subject argument of a control infinitive cannot be pronounced within the infinitive.
	\ex The subject argument of a control infinitive can be understood as the same as an argument in the higher clause: its ``controller.''
\end{xlist}
\end{exe}
Our goal is to understand what is responsible for these properties of control infinitives. Our hope should be that it is consistent with, maybe even follows from, the forces that are responsible for positioning arguments. We'll look at a few ideas in this section that might bring us within range of these desires. I'll start with the question of whether these infinitives are CPs or not, and then turn to issues that more transparently bear on (\ref{ex3:contrdesc1}).

\subsection{The CP/IP status of infinitives}

Despite the lack of a complementizer, there is some evidence that control infinitives can be a CP. Of the verbs that select a control infinitive, those that also s-select a question can combine with a control infinitive that begins with a wh-phrase; (\ref{ex3:conquest}) are some examples.\marginnote{The status of complementizers in infinitival clauses, and the whole array of differences in kinds of infinitival arguments, originates with \citet{Bresnan:1972}, who built upon \citet{Rosenbaum:1967}.}
\begin{exe}
\ex \label{ex3:conquest} 
\begin{xlist}
	\ex Caroline told Helen how to eat chocolate.
	\ex Caroline asked when to leave.
	\ex Caroline decided when to leave.
	\ex Caroline wondered what to eat.
\end{xlist}
\ex \label{ex3:conquest2}  
\begin{xlist}
	\ex Mariko told Paul whether to eat chocolate.
	\ex Mariko asked whether to leave.
	\ex Mariko decided whether to leave.
	\ex Mariko wondered whether eat chocolate.
\end{xlist}
\end{exe}
Recall that wh-phrases are moved into the Specifier of CP by Wh movement. We should conclude, then, that the infinitives in (\ref{ex3:conquest}) are CPs. And in these very contexts, the complementizer \emph{whether} is, as expected, also possible as (\ref{ex3:conquest2}) shows.

In many, not too distantly related, languages, even the non-question form of a control infinitive comes with a complementizer. Icelandic is a particularly clear example of this. It has subject and object control infinitives that closely resemble those of English. And yet these control infinitives appear with the same complementizer that is found in finite clauses in Icelandic: \emph{a\dh}.
\pagebreak
\begin{exe}
\ex \label{ex3:icecon} \gll Mar\'ia lofa{\dh}i \textbf{a{\dh}} lesa b\'okina. \\ Mary promised that read {the book}\\ `Mary promised to read the book.'
\sn \emph{compare:}
\sn \gll Mar\'ia segir \textbf{a{\dh}} {\th}\'u hafir lesi{\dh} b\'okina. \\ Mary says that you have read {the book}\\ `Mary says that you have read the book.'
\end{exe}
We should conclude, then, that a control infinitive can be a CP.

By contrast, the same kind of evidence regarding raising infinitives suggests that they are never CPs. There are no known examples of raising infinitives that have the syntax of questions.
\begin{exe}
\ex \label{ex3:nowhraise} 
\begin{xlist}
	\ex[*]{It seems whether to have left.\\ \emph{compare:} Nancy seems to have left.}
	\ex[*]{It appears when to eat chocolate.\\ \emph{compare:} Nancy appears to eat chocolate.}
	\ex[*]{It is likely who to dance.\\ \emph{compare:} Nancy is likely to dance.}
\end{xlist}
\end{exe}
This could be because infinitival clauses are incapable of being CPs, and therefore do not support the CP-dependent syntax of questions. But it could also be because raising verbs (and adjectives, such as \emph{likely}) do not semantically select questions. Whether a verb can take a question as an object depends on the semantics of the individual verb. A verb like \emph{believe}, for instance, cannot take a question complement, whereas the similar \emph{know} can.
\begin{exe}
\ex \label{ex3:questselect} 
\begin{xlist}
	\ex[*]{Ania believes whether Eva has left.}
	\ex[]{Ania knows whether Eva has left.}
\end{xlist}
\end{exe}
The absence of question complements to raising verbs could, then, simply reflect a fact about the semantics of these predicates, and nothing about the syntax of their infinitival complements. Indeed, raising predicates do not seem able to combine with question complements even when those complements are finite clauses.
\begin{exe}
\ex \label{ex3:nowhraise2} 
\begin{xlist}
	\ex[*]{It seems whether Elliot has left.\\ \emph{compare:} It seems that Elliott has left.}
	\ex[*]{It appears when Mar\'ia has eaten the chocolate.\\ \emph{compare:} It appears that Mar\'ia has eaten the chocolate.}
	\ex[*]{It is likely who will dance.\\ \emph{compare:} It is likely that Anna will dance.}
\end{xlist}
\end{exe}
Finite clauses can always be CPs (it seems), and so we should conclude that the absence of question complements to raising verbs has to do with their meaning.

Although we cannot conclude that raising infinitives are prevented from being CPs, we also cannot learn from their inability to be questions whether they are allowed to be CPs. Icelandic, however, provides a clue. Raising infinitives in Icelandic cannot include the complementizer~\emph{a{\dh}}.
\begin{exe}
\ex[*]{ \label{ex3:iceraise} \gll Mar\'ia haf{\dh}i virst a{\dh} hafa vaska{\dh} upp diskana.\\ Mary had seemed that {have} washed up {the dishes}\\ `Maria had seemed to have washed up the dishes'}
\sn \emph{compare:}
\sn[]{\gll Mar\'ia haf{\dh}i virst hafa vaska{\dh} upp diskana. \\ Maria had seemed have washed up {the dishes}\\ `Mary had seemed to wash up the dishes.'}
\end{exe}
This would be explained if raising infinitives are IPs, but cannot be CPs.

In Icelandic, then, we can conclude that one difference between raising and control infinitives is their categorial status. Raising infinitives are IPs but cannot be CPs, whereas control infinitives are CPs. In English, the picture isn't as clear. There is no direct evidence that raising infinitives cannot be CPs. And though there is evidence that some control infinitives are CPs---the questions---other control infinitives do not differ outwardly from raising infinitives. English internal evidence supports the conclusion that control infinitives can be CPs, but nothing more.

If we consider the evidence from Icelandic and English together, however, and keep in mind the criterion of explanatory adequacy, we can find some help in making decisions about the proper treatment of English infinitives. We conclude from Icelandic that the mapping in (\ref{ex3:infmap}) exists in at least one grammar.
\begin{exe}
\ex \label{ex3:infmap} 
\begin{xlist}
	\sn raising infinitive $\rightarrow$ IP
	\sn control infinitive $\rightarrow$ CP
\end{xlist}
\end{exe}
Assuming that this is not a capricious relationship, we can assume that there are forces, still to be discovered, that produce it. If the mapping in (\ref{ex3:infmap}) does not materialize in English, then we can conclude that those forces do not operate in the same way in English. When we model these forces, our model should allow them to vary from language to language. Consequently, we must credit the child with the ability to detect that variation as that child is acquiring his or her language. Let's call this ``hypothesis~1.''

Hypothesis 2 is that the forces responsible for the existence of (\ref{ex3:infmap}) in Icelandic are also found in English. On this hypothesis, Icelandic and English should be identical with respect to the categorial status of their infinitives. Any differences in the appearance of infinitives in these languages will flow from other ways these languages differ. What are the differences between Icelandic and English that would have to be explained on hypothesis 2? As we've seen, Icelandic and English raising infinitives look the same; nothing would have to be said under hypothesis 2 about how raising infinitives differ in these languages. Control infinitives in Icelandic and English also look the same in all ways except one: declarative control infinitives can appear with a complementizer in Icelandic but not in English. Hypothesis 2 requires that this difference be made compatible with the mapping in (\ref{ex3:infmap}). One way of doing that is with (\ref{ex3:nullc}).
\begin{exe}
\ex \label{ex3:nullc} The complementizers of infinitival clauses in English have no phonetic realization.
\end{exe}
It seems difficult to avoid a conclusion like (\ref{ex3:nullc}) under hypothesis~2. Hypothesis~2 is likely to require, then, that words can be silent. It must credit the child with the ability to detect the existence of silent words in a sentence.

So we have before us two alternatives. English does not conform to the mapping in (\ref{ex3:infmap}) and the forces that are responsible for (\ref{ex3:infmap}) can vary in ways that are detectable by children acquiring English. Or, English does conform to the mapping in (\ref{ex3:infmap}) but has a silent infinitival complementizer, and the existence of this silent complementizer is detectable by children learning English. Framing the analysis of English infinitives in this way allows us to make decisions based on criteria of language acquisition. If the forces behind (\ref{ex3:infmap}) turn out to be too distant from the information available to a child, then we can conclude that they cannot vary in detectable ways and we're stuck with hypothesis~2. If, by contrast, the existence of words in a sentence cannot be detected without their phonetic manifestation, then we're stuck with hypothesis~1. Hypothesis~2 is, perhaps, the most popular in the subsequent literature. One way of ensuring that the silent complementizer is detectable to the language-acquiring child is to elevate the mapping in (\ref{ex3:infmap}) to a language universal. If the child can rely on (\ref{ex3:infmap}), then it can deduce that there is a silent complementizer in the declarative control infinitives of English. Let's adopt hypothesis~2 as our working hypothesis as well.

Before leaving the question of whether infinitives are CPs or not, we should look at one last ``kind'' of infinitive. Some predicates that take control infinitives as objects can take an infinitive which does have an overt complementizer, and in which the subject argument is visible inside the infinitive. Some examples of these predicates and the two types of infinitives they can combine with, are given in~(\ref{ex3:forto}).
\begin{exe}
\ex \label{ex3:forto} 
\begin{xlist}
	\ex Misato wanted to eat durian.\\
		Misato wanted for him to eat durian.
	\ex Junko preferred to eat natto.\\
		Junko preferred for me to eat natto.
	\ex It is possible to eat poi.\\
		 It is possible for some to eat poi.
\end{xlist}
\end{exe}
The subject of the infinitival clauses appears in the accusative Case, as is clear from the form of the pronouns in (\ref{ex3:forto}a) and (\ref{ex3:forto}b). The complementizer is \emph{for}, which might easily be mistaken for the homophonous preposition found in sentences like \emph{She wanted something for Susan.} We think that the \emph{for} in (\ref{ex3:forto}) is a Complementizer for several reasons. First, it forms a constituent with the entire clause that follows it. This is indicated, among other things, by the behavior of these phrases in the context of cleft constructions.
\begin{exe}
	\ex \label{ex3:fortocleft} 
	\begin{xlist}
	\ex[]{It's \textbf{for him to eat chocolate} that Jason would prefer.}
	\ex[*]{It's to him how to eat chocolate that Jason should explain.}
	\end{xlist}
\end{exe}
The ungrammaticality of (\ref{ex3:fortocleft}b) derives from the fact that there is only room for one phrase between the \emph{it's} and the \emph{that} of these clefts, and in (\ref{ex3:fortocleft}b) two things, a PP and an infinitive, have been shoved into this spot. Thus the grammaticality of (\ref{ex3:fortocleft}a) would seem to argue that we do not want to parse the boldfaced string as a PP followed by an infinitive, but instead as something that makes a single phrase. To the extent that an infinitive is an IP, and that the \emph{him} in (\ref{ex3:forto}) is its subject, then one of the few ways of doing this is to let \emph{for} be a complementizer.

Moreover, the \emph{for} that comes with these infinitives has a property that no preposition in English has, but which the complementizer \emph{that} does have. It is optional. Under a certain set of circumstances, the complementizer \emph{that} can go unspoken in English, as in (\ref{ex3:delthat}).
\begin{exe}
	\ex \label{ex3:delthat} 
	\begin{xlist}
	\ex[]{Jason said that he eats chocolate.}
	\ex[]{Jason said  he eats chocolate.}
	\end{xlist}
\end{exe}
This is also possible for the \emph{for} of those infinitives that are objects to verbs, as (\ref{ex3:delto}) shows.
\begin{exe}
\ex \label{ex3:delto} 
\begin{xlist}
	\ex Misato wanted him to eat durian.
	\ex Junko preferred me to eat natto.
\end{xlist}
\end{exe}
Furthermore, the circumstances controlling the optionality of the complementizer \emph{that} are mimicked by the optionality of \emph{for}. In general, the complementizer \emph{that} can go unspoken only in contexts where the CP it heads is a complement to a verb. That is why \emph{that} cannot go missing when it heads a CP that has been clefted, as in (\ref{ex3:14}). A parallel constraint on \emph{for} ellipsis is indicated by (\ref{ex3:14a}).
\begin{exe}
	\ex \label{ex3:14} 
	\begin{xlist}
	\ex[]{It's that he eats chocolate that Jason said.}
	\ex[*]{It's he eats chocolate that Jason said.}
	\end{xlist}
	\ex \label{ex3:14a}
	\begin{xlist}
	\ex[]{It's for him to eat chocolate that Jason would prefer.}
	\ex[*]{It's him to eat chocolate that Jason would prefer.}
	\end{xlist}
\end{exe}
In general, the constraints on where \emph{that} can go missing are parallel to those on where \emph{for} can, suggesting that they are both subject to some process, restricted to complementizers, that allows them to go silent.


\subsection{On the Control relation}

Let's return now to how control infinitives express their arguments. In the cases we've looked at so far, they are rather like raising infinitives in that an argument of their main predicate surfaces in a position in the higher clause. One possibility, then, is that control infinitives use the same mechanism that raising infinitives do to put their arguments in this position. We might imagine that Argument movement relocates the argument DP of (\ref{ex3:movecon1}) in the way indicated in~(\ref{ex3:movecon1d}).
\begin{exe}
\ex \label{ex3:movecon1} Shai tried to sing.
\ex \label{ex3:movecon1d}
\small
\begin{forest}
[IP
	[DP [Shai, roof, name=binder]]
	[\ibar{I} 
		[\obar{I} [past] ]
		[VP
			[{\ }, name=trace2 ]
			[\ibar{V}
				[\obar{V} [try] ]
				[IP
					[{\ }, name=trace ]
					[\ibar{I} 
						[\obar{I} [to] ]
						[VP
							[{\ }, name=trace1 ]
							[\ibar{V} [sing, roof] ]
						 ]
					 ]
				 ]
			 ]
		 ]
	]
]
\draw[->,>=stealth', ultra thin, gray] (trace1) to[out=south west,in=-110,looseness=2] node[pos=0.75,xshift=10mm]{} (trace);
\draw[->,>=stealth', ultra thin, gray] (trace) to[out=-135,in=-105,looseness=2.3] node[pos=0.75,xshift=10mm]{} (trace2);
\draw[->,>=stealth', ultra thin, gray] (trace2) to[out=south west,in=south] node[pos=0.75,xshift=10mm]{} (binder);
\end{forest}
\end{exe}
An instance of object control, like that in (\ref{ex3:movecon2}) could have the derivation indicated.
\begin{exe}
\ex \label{ex3:movecon2} Shai told Florian to sing.\marginnote[10pt]{If this turns out to be correct, then we will either have to change the restriction on movement that prohibits movement to a non-c-commanding position, or we will have to jettison the view that phrases have no more than two daughters. With these two assumptions, the movement, as indicated here violates the c-command requirement.}
\ex \label{ex3:movecon2d}
\small
\begin{forest}
[IP
	[DP [Shai, roof, name=binder]]
	[\ibar{I} 
		[\obar{I} [past] ]
		[VP 
			[{\ }, name=trace2 ]
			[\ibar{V}
				[\ibar{V}
					[\obar{V} [tell] ]
					[DP , name=binder2 [Florian, roof] ]
				 ]
				[IP
					[{\ }, name=trace1 ]
					[\ibar{I}
						[\obar{I} [to] ]
						[VP
							[{\ }, name=trace ]
							[\ibar{V}[sing, roof] ]
						 ]
					 ]
				 ]
			 ]
		]
	]
]
\draw[->,>=stealth', ultra thin, gray] (trace) to[out=south west,in=-110,looseness=2] node[pos=0.75,xshift=10mm]{} (trace1);
\draw[->,>=stealth', ultra thin, gray] (trace1) to[out=-150,in=-35,looseness=1.2] node[pos=0.75,xshift=10mm]{} (binder2);
\draw[->,>=stealth', ultra thin, gray] (trace2) to[out=south west,in=south] node[pos=0.75,xshift=10mm]{} (binder);
\end{forest}
\vspace*{-5pt}
\end{exe}
These derivations differ from the ones we've posited up to now for Argument movement in that they move arguments into \mbox{$\theta$-marked} positions. If this is the correct account of control infinitives, then, we should adjust our definition of D-structure as the sole parse in a derivation that gets semantically inter-preted. That definition prevents a θ-marked position from being empty at D-structure, and also prevents an argument from gaining a θ-role by moving into a θ-marked position. But if we were to make it possible, somehow, to let parses in a derivation other than the D-structure be semantically interpreted, we would have a transparent account of the Control relation. We would also provide an account for how control infinitives satisfy the EPP and that part of the Theta Criterion that requires every \mbox{$\theta$-role} to be assigned to an argument. On this view, the Control relation is just the one established by Argument movement between an argument and its \mbox{$\theta$-marked} position. This view of the Control relation is offered in \citet{Hornstein:1999} and its consequences explored in \citet{Hornstein:2000} and \citet{Boeckx:2003a,Boeckx:2004b,Boeckx:2006c,Boeckx:2006b}.

There are circumstances where control infinitives do not involve control, however; in these circumstances the subject argument is not an argument expressed elsewhere. This is the case when, for instance, the verbs taking a control infinitive are passivized, as in (\ref{ex3:noc1}).
\begin{exe}
\ex \label{ex3:noc1} 
\begin{xlist}
	\ex To eat natto shouldn't be tried.
	\ex To be content is often wanted.
	\ex To be honest should be promised in a wedding vow.
\end{xlist}
\end{exe}
In fact, whenever a control infinitive stands in subject position, its understood subject does not have to be an argument found elsewhere; (\ref{ex3:noc2}) provides some other examples.\enlargethispage*{20pt}
\begin{exe}
\ex \label{ex3:noc2} 
\begin{xlist}
	\ex To remind him of Amherst Winters bothered Pius.
	\ex \label{ex3:noc2durian}To eat durian shouldn't be possible.
	\ex To leave your shoes untied is a danger to those around you.
	\ex To be the largest prime number is to be very large indeed.
\end{xlist}
\end{exe}
In these examples the subject of the infinitive is understood to have a meaning close to that which \emph{one} gets in sentences like (\ref{ex3:exone}), as can be appreciated by comparing this example to (\ref{ex3:noc2durian}).\marginnote[0pt]{There is a subtle, and indescribable, difference in the meanings of (\ref{ex3:exone}) and (\ref{ex3:noc2durian}), but it doesn't seem to me that this difference resides in the meaning of the subjects.}
\begin{exe}
\ex \label{ex3:exone} For one to eat natto shouldn't be possible.
\end{exe}
The understood subject of these infinitives is said to be ``arbitrary.'' 

In (\ref{ex3:noc1}) and (\ref{ex3:noc2}) there is no argument that could be the understood subject of the infinitive and so it is easy to see that it must get another meaning. But there are also examples where a potential controlling argument is present and yet an arbitrary reading for the understood subject of the infinitive is available. That is the case for the subject argument in the examples in (\ref{ex3:noc3}), for example, which may all have the same arbitrary interpretation that is found in~(\ref{ex3:noc2}).
\begin{exe}
\ex \label{ex3:noc3} 
\begin{xlist}
	\ex Elena discovered that to remind Pius of Amherst winters bothers him.
	\ex Jeff thought that to eat durian shouldn't be possible.
	\ex Bart remembered that to leave his shoes untied is a danger to those around him.
	\ex Jill knew that to be the largest prime number is to be very large indeed.
\end{xlist}
\end{exe}
When the infinitive is a question, it may also sometimes support an arbitrary interpretation for its understood subject, as in (\ref{ex3:noc4}).
\begin{exe}
\ex \label{ex3:noc4} 
\begin{xlist}
	\ex Jon knew how to eat live shrimp.
	\ex Satoshi can tell when to add vinegar to rice.
	\ex Julie decided where to put the \emph{Dicentra}.
\end{xlist}
\end{exe}
These can get interpretations roughly parallel to those in (\ref{ex3:noc4a}), in which the subject is arbitrary.
\begin{exe}
\ex \label{ex3:noc4a} 
\begin{xlist}
	\ex Jon knew how one can eat live shrimp.
	\ex Satoshi can tell when one should add vinegar to rice.
	\ex Julie decided where one should put the \emph{Dicentra}.
\end{xlist}
\end{exe}
The precise conditions under which the arbitrary interpretation are available are imperfectly understood. See \citet{Williams:1980,Williams:1989} and \citet{Landau:2004} for some ideas.

Control infinitives with an arbitrary reading pose the same problem for the EPP and Theta Criterion that normal control infinitives do. They seem to violate the EPP---there is nothing in the Specifier of their IP---and because the \mbox{$\theta$-role} normally assigned to the subject has no argument to bear it, they seem to violate the Theta Criterion as well. Moreover, just like normal control infinitives, they are ungrammatical if there is no \mbox{$\theta$-role} for a subject argument, as (\ref{ex3:noctc1}) shows.
\begin{exe}
\ex[*]{ \label{ex3:noctc1} To seem that pigs fly is worrisome.\\ \emph{compare:} For it to seem that pigs fly is worrisome.}
\end{exe}
Just as for normal control infinitives, then, these facts taken together suggest that there is a subject argument present in these infinitives. If such an argument is present, then the EPP and Theta Criterion could be maintained. Moreover, the ungrammaticality of examples like (\ref{ex3:noctc1}) would be explained: if there is a subject argument, then the Theta Criterion requires that there also be a \mbox{$\theta$-role} assigned to it. In these cases, however, this alleged subject argument cannot be found in some unexpected position. We cannot here finger the controller as that argument because there is no controller. 

A common reaction, and the one I shall adopt, is that there is a subject argument syntactically present in these examples but that it is silent.\marginnote[0pt]{For an approach that does not countenance a silent subject, see \citet{Chierchia:1984}.} This argument is indicated by putting ``\textsc{pro}'' in the position our syntax suggests it should be. The representation for (\ref{ex3:noc2durian}), then, would be (\ref{ex3:noc2duriant}).
\begin{exe}
\ex \label{ex3:noc2duriant} 
\small
\begin{forest}
[IP
	[CP, name=binder1
		[\obar{C} [e] ]
		[IP
			[DP [\textsc{pro}, roof, name=binder2] ]
			[\ibar{I}
				[\obar{I} [to] ]
				[VP
					[{\ }, name=trace ]
					[\ibar{V} [eat durian, roof] ]
				 ]
			 ]
		 ]
	 ]
	[\ibar{I}
		[\obar{I} [shouldn't] ]
		[VP [\ibar{V}
			[\obar{V} [be] ]
			[AP 
				[{\ }, name=trace1 ]
				[\ibar{A} [possible, roof] ]
			 ]
		] ]
	 ]
]
\draw[->,>=stealth', ultra thin, gray] (trace) to[out=south west,in=south] node[pos=0.75,xshift=10mm]{} (binder2);
\draw[->,>=stealth', ultra thin, gray] (trace1) to[out=south west,in=east] node[pos=0.75,xshift=10mm]{} (binder1);
\end{forest}
\end{exe}
I've decided that this infinitival clause is a CP with a silent complementizer, represented here with ``e.'' I've also shown \textsc{pro} moving from the position that the external \mbox{$\theta$-role} is assigned to, Specifier of VP, to the Specifier of IP where it satisfies the EPP. I also show the infinitival CP itself moving from Specifier of AP into Specifier of IP; this claims that the infinitival gets an external \mbox{$\theta$-role} from \emph{possible} and moves in order to satisfy the EPP. 

We can think of \textsc{pro} as being a kind of pronoun (hence its name), maybe along the lines of a silent \emph{one}. Of course, if the grammar of English includes a silent pronoun, then where it can be must be heavily constrained. Otherwise, sentences such as (\ref{ex3:proth}) should be grammatical, with Theta Criterion satisfied by the presence of \textsc{pro} in subject position or object of preposition.
\begin{exe}
\ex \label{ex3:proth} 
\begin{xlist}
	\ex[*]{Ran into the garden.}
	\ex[*]{Andre talked to.}
\end{xlist}
\end{exe}
The problem of restricting \textsc{pro}'s distribution remains unsolved. In fact, the question of where \textsc{pro} can be is still controversial. There is no controversy over the question of whether \textsc{pro} can be found in complement positions; it seems clear that it cannot. And it also seems clear that \textsc{pro} cannot be in subject position of finite clauses, at least not in English. But it's not clear yet in which non-finite clauses \textsc{pro} can be a subject.

Perhaps the most popular view presently is that \textsc{pro} can stand as subject to any control infinitive, not just those that get an arbitrary interpretation. On this view, the Control it induces does not reduce to a special instance of Argument Movement, but is instead an interpretation that \textsc{pro} can receive. The representation for sentences involving Control would not be as indicated in (\ref{ex3:movecon1}) and (\ref{ex3:movecon2}), but would instead include \textsc{pro}, as in (\ref{ex3:procon1d}) and (\ref{ex3:procon2d}).
\begin{exe}
\ex \label{ex3:procon1d} Shai tried to sing.
\hspace*{15pt}
\small
\begin{forest}
[IP
	[DP [Shai, roof, name=binder]]
	[\ibar{I} 
		[\obar{I} [past] ]
		[VP
			[{\ }, name=trace ]
			[\ibar{V} 
				[\obar{V} [try] ]
				[CP [\ibar{C} 
					[\obar{C} [e] ]
					[IP
						[DP [\textsc{pro}, roof,name=binder1] ]
						[\ibar{I}
							[\obar{I} [to] ]
							[VP
								[{\ }, name=trace1 ]
								[\ibar{V}[sing, roof] ]
							 ]
						 ]
					 ]
				] ]
			 ]
		 ]
	]
]
\draw[->,>=stealth', ultra thin, gray] (trace) to[out=south west,in=south] node[pos=0.75,xshift=10mm]{} (binder);
\draw[->,>=stealth', ultra thin, gray] (trace1) to[out=south west,in=south] node[pos=0.75,xshift=10mm]{} (binder1);
\end{forest}
\vspace*{-20pt}
\ex \label{ex3:procon2d} Shai told Florian to sing.
\hspace*{10pt}
\small
\begin{forest}
[IP
	[DP [Shai, roof, name=binder]]
	[\ibar{I} 
		[\obar{I} [past] ]
		[VP 
			[{\ }, name=trace ]
			[\ibar{V} 
				[\ibar{V}
					[\obar{V} [tell] ]
					[DP [Florian, roof] ]
				 ]
				[CP [\ibar{C}
					[\obar{C} [e] ]
					[IP
						[DP [\textsc{pro}, roof, name=binder1]]
						[\ibar{I} 
							[\obar{I} [to] ]
							[VP
								[{\ }, name=trace1 ]
								[\ibar{V} [sing, roof] ]
							 ]
						]
					 ]
				] ]
			]
		 ]
	]
]
\draw[->,>=stealth', ultra thin, gray] (trace) to[out=south west,in=south] node[pos=0.75,xshift=10mm]{} (binder);
\draw[->,>=stealth', ultra thin, gray] (trace1) to[out=south west,in=south] node[pos=0.75,xshift=10mm]{} (binder1);
\end{forest}
\end{exe}
Unlike the movement-based analysis of Control that (\ref{ex3:movecon1}) and (\ref{ex3:movecon2}) illustrate, the representations in (\ref{ex3:procon1d}) and (\ref{ex3:procon2d}) obey the Theta Criterion as it is now formulated. What we would need to do if these representations are correct is produce a non-movement account of Control, an account that would emerge as part of a more general account of how \textsc{pro} gets its semantic interpretation.

What we've discovered in this section, then, is that there is a silent pronominal argument in English and that this silent argument is prevented from appearing in a variety of positions otherwise normally available to arguments. Keeping in mind that this argument, \textsc{pro}, cannot be in complement positions, we can partially describe its special distribution with~(\ref{ex3:prorest}).
\begin{exe}
\ex \label{ex3:prorest} \textsc{The \textsc{pro} Restriction}\\
\textsc{pro} can only be in certain Specifier positions at S-structure.
\end{exe}
Which Specifier positions \textsc{pro} can be in is controversial, but it certainly doesn't include the Specifiers of finite IP.

The \textsc{pro} Restriction describes something that is peculiar to the placement of the silent argument: \textsc{pro}. Let's now look at whether the conditions we've been developing to capture the placement of the non-silent arguments also hold of \textsc{pro}.

\subsection{Argument movement and \textsc{pro}}

To this point, Argument movement has been invoked in situations where a DP is \mbox{$\theta$-marked} in a position that is not Case marked. Because argument DPs must be both Case marked and receive a \mbox{$\theta$-role}, there will be no single parse that can satisfy both of these conditions in such a case. Argument movement therefore produces a series of parses: one of which will satisfy the Theta Criterion and Projection Principle and another one of which will satisfy the Case filter. I have expressed the Theta Criterion and Projection Principle as well-formedness conditions on D-structures, and the Case filter as a well-formedness condition on derivations. This places certain conditions on how Argument movement will arise because it imposes requirements on what the input and output parses can be. If Control is an instance of Argument movement, then certain changes to the Theta Criterion will have to be made, and there will be consequent changes to how Argument movement behaves.

In addition to the constraints imposed on Argument movement by these conditions, we have seen that there is a condition imposed directly on Argument movement, and other movement operations: the moved term can only be relocated to a position that c-commands its starting position. This condition is responsible for blocking (\ref{ex3:nonccoma}).
\begin{exe}
\ex[*]{ \label{ex3:nonccoma} It was shown that Mari seems that pigs fly.}
\end{exe}
If Argument movement could move DPs to non-c-commanding positions, then (\ref{ex3:nonccoma}) should be able to have the derivation indicated in (\ref{ex3:nonccomad}).
\begin{exe}
\ex \label{ex3:nonccomad} 
\small
\begin{forest}
[IP
	[DP [it, roof]]
	[\ibar{I} 
		[\obar{I} [past] ]
		[VP [\ibar{V}
			[\obar{V} [be] ]
			[VP [\ibar{V}
				[\ibar{V}
				[\obar{V} [shown] ]
				[{\ }, name=trace ]
			]
			[CP [\ibar{C}
				[\obar{C} [that] ]
				[IP
					[DP [Mari, roof, name=binder] ]
					[\ibar{I}
						[\obar{I} [pres] ]
						[VP [\ibar{V}
							[\obar{V} [seem] ]
							[CP [that pigs fly, roof] ]
						] ]
					 ]
				 ]
			]] 
			]
		]
	]
]]]
\draw[->,>=stealth', ultra thin, gray] (trace) to[out=south,in=south west] node[pos=0.75,xshift=10mm]{} (binder);
\end{forest}
\end{exe}
The derivation satisfies every other condition we've discovered. The EPP is satisfied for both Specifier of IPs. The number of \mbox{$\theta$-roles} and arguments is the same, so the Theta Criterion is satisfied. And every argument stands as a sister to the term that assigns it a \mbox{$\theta$-role} at D-structure, and so the Projection Principle is satisfied as well. In particular, the argument \emph{Mari} stands as a sister to the verb \emph{shown} at D-structure, and there gets the object \mbox{$\theta$-role} from \emph{shown}. Because \emph{show} has been Passivized, however, it does not have an accusative Case to assign to \emph{Mari}, and so Argument movement must move this DP to a Case marked position. The Specifier of the finite IP that follows \emph{Mari} is such a Case marked position, and so movement to this position should satisfy the Case filter. If this position c-commanded the D-structure position of \emph{Mari}, every other condition would be satisfied and this sentence should be grammatical. It's the c-command requirement on Argument movement, then, that is responsible for preventing sentences such as~(\ref{ex3:nonccoma}). This can be appreciated by comparing the grammatical (\ref{ex3:nonccomb}), in which the c-command requirement is satisfied.
\begin{exe}
\ex \label{ex3:nonccomb} Mari was shown that it seems that pigs fly.
\small
\begin{forest}
[IP
	[DP [Mari, roof, name=binder]]
	[\ibar{I} 
		[\obar{I} [past] ]
		[VP [\ibar{V}
			[\obar{V} [be] ]
			[VP [\ibar{V}
				[\obar{V}
					[\obar{V} [shown] ]
					[{\ }, name=trace ]
				 ]
				[CP [\ibar{C}
					[\obar{C} [that] ]
					[IP
						[DP [it, roof]]
						[\ibar{I}
							[\obar{I} [pres] ]
							[VP [\ibar{V}
								[\obar{V} [seem] ]
								[CP [that pigs fly, roof] ]
							] ]
						]
					 ]
				] ]
			] ]
		] ]
	]
]
\draw[->,>=stealth', ultra thin, gray] (trace) to[out=south west,in=south] node[pos=0.75,xshift=10mm]{} (binder);
\end{forest}
\end{exe}

If \textsc{pro} is an argument DP, then we should expect that Argument movement will be responsible for fixing its position and, consequently, we should find a parallel pattern of grammatical and ungrammatical sentences that the successful or unsuccessful application of Argument movement is responsible for. In the case of overt DPs, the success of Argument movement hinges on being able to find an S-structure that satisfies the Case filter that can be produced from a grammatical D-structure through an application of Argument movement that satisfies the c-command requirement. In the case of \textsc{pro}, the grammatical sentences should be ones in which \textsc{pro} can reach the Specifier of an appropriate infinitival IP from its \mbox{$\theta$-marked} D-structure position by way of an application of Argument movement that satisfies the c-command requirement. 

This  match in grammaticality patterns shows up. The ungrammaticality of (\ref{ex3:nonccomapro}), for instance, parallels the ungrammaticality of (\ref{ex3:nonccoma}).
\begin{exe}
\ex[*]{ \label{ex3:nonccomapro} It was shown to seem that pigs fly.}
\end{exe}
If the infinitival clause in (\ref{ex3:nonccomapro}) is of the sort in which \textsc{pro} is permitted, then this example satisfies every condition except the c-command constraint on Argument movement, as (\ref{ex3:nonccomaprod}) indicates.
\begin{exe}
\ex \label{ex3:nonccomaprod} 
\small
\begin{forest}
[IP
	[DP [it, roof]]
	[\ibar{I} 
		[\obar{I} [past] ]
		[VP [\ibar{V}
			[\obar{V} [be] ]
			[VP [\ibar{V}
				[\ibar{V}
					[\obar{V} [shown] ]
					[{\ }, name=trace ]
				 ]
				[CP [\ibar{C} 
					[\obar{C} [e] ]
					[IP
						[DP [\textsc{pro}, roof, name=binder] ]
						[\ibar{I}
							[\obar{I} [to] ]
							[VP [\ibar{V}
								[\obar{V} [seem] ]
								[CP [that pigs fly, roof]]
							] ]
						 ]
					 ]
				] ]
			] ]
		] ]
	]
]
\draw[->,>=stealth', ultra thin, gray] (trace) to[out=south,in=south west,looseness=1] node[pos=0.75,xshift=-10mm]{} (binder);
\end{forest}
\end{exe}
That this infinitive is of the kind that permits \textsc{pro} in its Specifier position is indicated by (\ref{ex3:noncomaproc}), which permits an arbitrary reading for the subject of its infinitive.
\begin{exe}
\ex \label{ex3:noncomaproc} Mari was shown how to fly.
\end{exe}
When the c-command requirement on Argument movement is obeyed, then examples parallel to (\ref{ex3:nonccomaprod}) are grammatical, as (\ref{ex3:noncomaprocd}) indicate.
\pagebreak
\begin{exe}
\ex \label{ex3:noncomaprocd} To be shown how to fly is swell.
\sn
\small
\begin{forest}
[IP
	[CP [\ibar{C}
		[\obar{C} [e] ]
		[IP 
			[DP [\textsc{pro}, roof, name=binder] ]
			[\ibar{I}
				[\obar{I} [to] ]
				[VP [\ibar{V}
					[\obar{V} [be] ]
					[VP [\ibar{V}
						[\obar{V} [shown] ]
						[{\ }, name=trace ]
					]
					[CP [how to fly, roof]] 
					]
				] ]
			 ]
		 ]
	]]
	[\ibar{I}
		[\obar{I} [pres] ]
		[VP [swell, roof] ]
	 ]
]
\draw[->,>=stealth', ultra thin, gray] (trace) to[out=south west,in=south,looseness=1] node[pos=0.75,xshift=10mm]{} (binder);
\end{forest}
\end{exe}
Examples like these, then, suggest that the position of \textsc{pro} is controlled by the same set of constraints that we've seen determine the position of overt argument DPs. To the extent that this is true, it speaks on behalf of the existence of \textsc{pro}. If the ungrammaticality of sentences like (\ref{ex3:nonccoma}) are due to conditions on where DP arguments may appear in S-structures, and those conditions are mimicked by arguments which are semantically present but not spoken, then we should conclude that the syntactic position of those unspoken arguments matters. They must be present in the syntactic representation for that to hold.

There is another condition on Argument movement of overt DPs that we can seek in the Argument movement of \textsc{pro}. This condition we have not yet encountered; it concerns how the Case filter is satisfied. It appears that DPs cannot be Argument moved from positions that are Case marked.\marginnote{This condition can be found in \citet{Chomsky:1981}, and for a more recent formulation: \citet{Chomsky:2000} and \citet{Chomsky:2001a}.}
\begin{exe}
\ex \label{ex3:earlycase} If Argument movement moves something from position $\alpha$, then $\alpha$\ cannot be a Case-marked position.
\end{exe}
This condition is thought to be responsible for the ungrammaticality of~(\ref{ex3:earlyex1}).
\begin{exe}
\ex[*]{ \label{ex3:earlyex1} Elena seems to that pigs fly.\\
(\emph{compare:} It seems to Elena that pigs fly.)}
\end{exe}
The argument DP \emph{Elena} resides in its \mbox{$\theta$-marked} position, which is the complement position of \emph{to}, and from there moves into the Case marked Specifier of IP position, thereby producing a representation that satisfies the EPP. The D-structure and S-structure representations meet the well-formedness requirements that they need to meet, and the position that \emph{Elena} moves into c-commands the position it moves from. Even though it meets all the other requirements on a grammatical sentence, (\ref{ex3:earlyex1}) is ungrammatical because it violates (\ref{ex3:earlycase}). The position that \emph{Elena} resides in at D-structure is assigned accusative Case by \emph{to}, and (\ref{ex3:earlycase}) prohibits Argument movement from a Case-marked position.

Another example whose ungrammaticality relies on (\ref{ex3:earlycase}) is (\ref{ex3:earlyex2}).
\begin{exe}
\ex[*]{ \label{ex3:earlyex2} Elena seems that is happy.}
\end{exe}
The derivation that needs to be blocked here is shown in~(\ref{ex3:earlyex2d}).
\begin{exe}
\ex \label{ex3:earlyex2d} 
\small
\begin{forest}
[IP
	[DP [Elena, roof, name=binder]]
	[\ibar{I} 
		[\obar{I} [pres] ]
		[VP [\ibar{V} 
			[\obar{V} [seem] ]
			[CP [\ibar{C}
				[\obar{C} [that] ]
				[IP
					[{\ }, name=trace	 ]
					[\ibar{I}
						[\obar{I} [pres] ]
						[VP [\ibar{V}
							[\obar{V} [be] ]
							[AP
								[{\ }, name=trace1]
								[\ibar{A} [happy, roof] ]
							 ]
						] ]
					 ]
				 ]
			] ]
		] ]
	]
]
\draw[->,>=stealth', ultra thin, gray] (trace1) to[out=south west,in=-100,looseness=1.6] node[pos=0.75,xshift=-10mm]{} (trace);
\draw[->,>=stealth', ultra thin, gray] (trace) to[out=south west,in=south,looseness=1] node[pos=0.75,xshift=-10mm]{} (binder);
\end{forest}
\end{exe}
For this sentence to satisfy the EPP, it is necessary for Argument movement to bring \emph{Elena} into the Specifier of the embedded IP. But this position is assigned nominative Case by the finite I, and so movement from this position is prohibited by (\ref{ex3:earlycase}). There is consequently no way to derive (\ref{ex3:earlyex2}) that does not violate either the EPP or~(\ref{ex3:earlycase}).

There is another feature of the derivation in (\ref{ex3:earlyex2}) that might be responsible for its ungrammaticality. The argument has moved out of a CP in (\ref{ex3:earlyex2}), and this is the first example we have considered in which that is true. In grammatical examples of Argument Movement out of embedded clauses, the clauses have been infinitival IPs. There is some evidence that Argument Movement out of CPs is independently prevented, and so that could also be the source of the ungrammaticality of~(\ref{ex3:earlyex2}). That evidence comes from consideration of cases like (\ref{ex3:nofin1}).
\begin{exe}
% \raggedright
	 \sn $\star$This fact seems that it is obvious. \textit{and} $\star$ This fact is possible for it to be obvious.
	\ex \label{ex3:nofin1}
	\small
	\begin{forest}
	[IP
		[DP [this fact, name=binder, roof]]
		[\ibar{I}
			[\obar{I} [pres]]
			[VP [\ibar{V}
				[\obar{V} [seem]]
				[CP [\ibar{C}
					[\obar{C} [that]]
					[IP
						[DP [\ibar{D} [\obar{D} [it\\expletive ]]]]
						[\ibar{I}
							[\obar{I} [pres]]
							[VP [\ibar{V}
								[\obar{V} [be]]
								[AP [\ibar{A} 
									[\obar{A} [obvious]]
									[{\ }, name=trace ]
								  ]]
							 ] ]
						 ]
					 ]
				 ]]
			] ]
		 ]
	]
	\draw[->,>=stealth', thin, gray] (trace) to[out=south,in=south,looseness=1.2] node[pos=0.75,xshift=10mm]{} (binder);
	\end{forest}
	\hspace*{5pt}
	\begin{forest}
	[IP
		[DP [this fact, name=binder, roof]]
		[\ibar{I}
			[\obar{I} [be]]
			[VP [\ibar{V}
				[\obar{V} [be]]
				[AP [\ibar{A} 
					[\obar{V} [possible]]
					[CP [\ibar{C} 
						[\obar{C} [for]]
						[IP
							[DP [\ibar{D} [\obar{D} [it]] ]]
							[\ibar{I}
								[\obar{I} [to]]
								[VP [\ibar{V}
									[\obar{V} [be]]
									[AP [\ibar{A} 
										[\obar{A} [obvious]]
										[{\ }, name=trace ]
									 ]]
								 ]]
							 ]
						 ]
					]]
				]]
			] ]
		 ]
	]
	\draw[->,>=stealth', thin, gray] (trace) to[out=south,in=south,looseness=1.3] node[pos=0.75,xshift=10mm]{} (binder);
	\end{forest}
\end{exe}
\vspace*{-40pt}
On the face of it, these derivations should be grammatical. The D-structures violate the EPP and the Case filter, but both of these maladies are remedied by using Argument Movement to move into the Specifier of the largest IP. The fact that these examples are ungrammatical suggests that there must be something that prevents Argument Movement from applying in this way. Cases of Argument Movement out of CPs in English are not attested; (\ref{ex3:tensedS}) seems true. 
\begin{exe}
\raggedright
	\ex \label{ex3:tensedS} If Argument Movement moves XP from position $\alpha$ to position $\beta$, then there cannot be a CP that dominates $\alpha$ but not~$\beta$.
\end{exe}
This constraint, along with (\ref{ex3:earlycase}), limit where Argument Movement can occur. 

There are examples parallel to the ones supporting (\ref{ex3:tensedS}) and (\ref{ex3:earlycase}) that involve \textsc{pro}, and they too speak on behalf of the existence of \textsc{pro} and letting Argument movement play a role in positioning \textsc{pro}. Parallel to (\ref{ex3:earlyex1}) is (\ref{ex3:earlyex1pro}).
\begin{exe}
\ex[*]{ \label{ex3:earlyex1pro} To seem to that pigs fly is swell.}
\end{exe}
This example would be grammatical if it were possible to Argument move \textsc{pro} from its D-structure position following \emph{to} into the Specifier of the infinitival clause, as shown in~(\ref{ex3:procasemove}).
\begin{exe}
\raggedright
	\ex \label{ex3:procasemove}
	\small
	\begin{forest}
	[IP
		[CP [\ibar{C}
			[\obar{C} [e]]
			[IP
				[PRO, name=binder ]
				[\ibar{I}
					[\obar{I} [to] ]
					[VP [\ibar{V}
					[\ibar{V} 
						[\obar{V} [seem]]
						[PP [\ibar{P} 
							[\obar{P} [to]]
							[{\ }, name=trace ]
						] ]
					]
					[CP [that pigs fly, roof]] 
					]
					]
				 ]
			 ]
		] ]
		[\ibar{I}
			[\obar{I} [pres]]
			[VP [be swell, roof]]
		 ]
	]
	\draw[->,>=stealth', thin, gray] (trace) to[out=south east,in=south,looseness=1] node[pos=0.75,xshift=10mm]{} (binder);
	\end{forest}
\end{exe}
To explain its ungrammaticality in a way that relates it to (\ref{ex3:earlyex1}), we should credit the ban against Argument movement from Case marked positions. 

Parallel to (\ref{ex3:earlyex2}) is (\ref{ex3:earlyex2pro}).
\begin{exe}
\ex[*]{ \label{ex3:earlyex2pro} To seem that is happy is swell.}
\end{exe}
The derivation we should be considering for (\ref{ex3:earlyex2pro})---the one that satisfies the EPP---is (\ref{ex3:earlyex2prod}).%\enlargethispage*{13pt}
\begin{exe}
\ex \label{ex3:earlyex2prod} 
\small
\begin{forest}
[IP
	[CP [\ibar{C}
		[\obar{C} [e] ]
		[IP
			[DP [\textsc{pro}, roof, name=binder]]
			[\ibar{I}
				[\obar{I} [to] ]
				[VP [\ibar{V}
					[\obar{V} [seem] ]
					[CP [\ibar{C}
						[\obar{C} [that] ]
						[IP
							[{\ }, name=trace ]
							[\ibar{I}
								[\obar{I} [pres] ]
								[VP [\ibar{V} 
									[\obar{V} [be] ]
									[AP
										[{\ }, name=trace1 ]
										[\ibar{A} [happy, roof] ]
									 ]
								] ]
							 ]
						 ]
					] ]
				] ]
			 ]
		 ]
	]]
	[\ibar{I} 
		[\obar{I} [pres] ]
		[VP [be swell, roof] ]
	]
]
\draw[->,>=stealth', ultra thin, gray] (trace1) to[out=south west,in=-100,looseness=1.5] node[pos=0.75,xshift=-10mm]{} (trace);
\draw[->,>=stealth', ultra thin, gray] (trace) to[out=south west,in=south,looseness=1] node[pos=0.75,xshift=-10mm]{} (binder);
\end{forest}
\end{exe}
Just like (\ref{ex3:earlyex2d}), this violates the ban against Argument movement out of a~CP.

A final condition influencing Argument movement that might be relevant for the placement of \textsc{pro} is the Case filter. If \textsc{pro} is an argument DP, then we should expect it to obey the Case filter. Is Argument movement invoked to put \textsc{pro} into Case marked positions at S-structure just as it is to put overt DPs into Case marked positions?

It's not easy to determine the answer to this question because of the confound that the \textsc{pro} Restriction presents. The \textsc{pro} Restriction describes the fact that \textsc{pro} seems to occur at S-structure in the Specifiers of control infinitives and nowhere else. If \textsc{pro} is also subject to the Case filter, then the narrower distribution that the \textsc{pro} Restriction allows it will make it difficult to detect. We are reduced, in fact, to the question of whether the positions that the \textsc{pro} Restriction allows \textsc{pro} to be in are Case marked positions.

We presently have two ways of determining whether a position is a Case-marked one. If an overt DP, bearing Case, is capable of surfacing in some position, then that position in assigned Case. And if Argument movement is prohibited from some position, then through a process of elimination we can determine whether the ban against moving from a Case-marked position is to blame and in this way discover whether a position is Case marked. By the first criterion, we cannot conclude that the Specifier positions in which \textsc{pro} resides are Case-marked. The Specifiers of control infinitives cannot contain overt DPs.
\begin{exe}
\ex \label{ex3:noovertcontrol} 
\begin{xlist}
	\ex[*]{Her to run seems unlikely.}
	\ex[*]{She to sing is swell.}
\end{xlist}
\end{exe}
From this, we might conclude that the positions \textsc{pro} resides in are not Case marked positions. This would explain why overt DPs cannot surface in this position: the Case filter would be violated if so. This is the conclusion reached in \emph{Lectures on Government and Binding}, and it informs most of the literature on \textsc{pro} from the 1980s and 1990s.

The second criterion gives a somewhat more mixed image. If the movement account of Control is correct, then the ban against moving from a Case-marked position allows us to deduce that the Specifier of the control infinitive that a controller moves from is not Case marked. So, for instance, we can conclude that the Specifier position of the infinitival complement to a verb like \emph{try} is not Case-marked if there is Argument movement of the sort indicated in (\ref{ex3:illcf}).
\begin{exe}
\ex \label{ex3:illcf} \linkto Hotze \under tried [\CP [\IP \from{\ \ } to eat durian ]].
\end{exe}
If the infinitive that controllers move from is the same kind that can host \textsc{pro} in non-control environments, then the simplest assumption would be that the Specifier of these infinitives is not Case-marked when they host \textsc{pro} as well. So, for instance, if the infinitive clause that stands in the subject position of a Passivized \emph{try} is the same as that which stands as the object of an Active \emph{try}, as in (\ref{ex3:illcf2}), then the \textsc{pro} this infinitive contains is standing in a non-Case-marked position.
\begin{exe}
\ex \label{ex3:illcf2} [\CP [\IP \textsc{pro} to eat durian ]] shouldn't be tried.
\end{exe}

On the other hand, if the movement account of Control is incorrect, then we cannot use the prohibition against movement from a Case-marked position to deduce that the subject position of a control infinitive is Case-marked. And there are examples whose ungrammaticality could be derived from this prohibition if the subject position of control infinitives are Case-marked. One of these is~(\ref{ex3:illcf3}).
\begin{exe}
\ex[*]{ \label{ex3:illcf3} \linkto Winnie \under was tried [\CP [\IP \from{\ \ } to eat durian ]].}
\end{exe}
If the prohibition against Argument movement from a Case-marked position is not responsible for the ungrammaticality of (\ref{ex3:illcf3}), then it's not clear why it should be bad. There is only one \mbox{$\theta$-role} for the argument \emph{Winnie}, and so the Theta Criterion is satisfied, as is the Case filter and the c-command restriction on movement. Indeed, the contrast between (\ref{ex3:illcf3}) and (\ref{ex3:illcf}) is problematic for the movement account of Control. The only relevant difference between these examples is the number of \mbox{$\theta$-roles} involved: in (\ref{ex3:illcf3}), Passive has removed \emph{try} of its external \mbox{$\theta$-role}, whereas in (\ref{ex3:illcf}) this \mbox{$\theta$-role} remains. The movement account of Control relies on an as yet undiscovered force that will allow Argument movement from a control infinitive into a \mbox{$\theta$-marked} position but prevent Argument movement from the same position that is not into a \mbox{$\theta$-marked} position. A theory of Control that makes it a special reading for \textsc{pro}, however, can make sense of this contrast with the tools at hand. Let the Specifier of a control infinitive be a Case-marked position, and everything follows. The Theta Criterion as well as the ban against movement from a Case-marked position will prevent Argument movement from a control infinitive into a \mbox{$\theta$-marked} position. The movement that is illustrated in (\ref{ex3:illcf}) will be blocked, and instead the subject of the embedded infinitival clause will be \textsc{pro}. The ban against Argument movement from a Case-marked position will continue to prevent Argument movement from a control infinitive even when that movement is not to another \mbox{$\theta$-marked} position, as in (\ref{ex3:illcf3}). 

On this set of assumptions, then, what makes a control infinitive different from, say, the infinitival complement to a verb like \emph{believe} is how the Specifier of the infinitive is Case-marked. The subject position of the infinitive that follows \emph{believe} is accusative Case-marked by \emph{believe}, and so Argument movement is possible from this position when \emph{believe} is Passivized and thereby robbed of its accusative Case.\marginnote[-70pt]{Configurations like (\ref{ex3:believeill}) are called instances of ``exceptional Case-marking,'' or ``Raising to Object.'' The second name references an analysis of these cases that we will examine in the next chapter. There are a variety of ways in which this construction can manifest itself. The verb \emph{wager}, for instance, does not allow sentences like (\ref{ex3:believeill}a), but it does allow sentences like (\ref{ex3:believeill}b).
\begin{exe}
	\raggedright
	\exi{i.}[*]{Bart wagered him to be good.}
	\exi{ii.}[]{He was wagered to be good.}
\end{exe}
See \citet{Postal:1974} for a thorough examination of the kinds of constructions that verbs of this type allow.}
\begin{exe}
\ex \label{ex3:believeill} 
\begin{xlist}
	\ex Bart believed [\IP him to be happy ].
	\ex He was believed [\IP to be happy ].
\end{xlist}
\end{exe}
The S-structure parse for (\ref{ex3:believeill}a) is~(\ref{ex3:believeillt}).
\begin{exe}
	\raggedright
	\ex \label{ex3:believeillt} 
	\small
	\begin{forest}
	[IP
		[DP [Bart, roof]]
		[\ibar{I} 
			[\obar{I} [past] ]
			[VP [\ibar{V}
				[\obar{V} [believe] ]
				[IP
					[DP [him, roof]]
					[\ibar{I} 
						[\obar{I} [to] ]
						[VP [\ibar{V}
							[\obar{V} [be] ]
							[AP [\ibar{A} [\obar{A} [happy]]] ]
						] ]
					 ]
				 ]
			] ]
		]
	]
	\end{forest}
\end{exe}
If the infinitival clause that is c-selected by \emph{believe} is an IP, then \emph{believe} will govern the Specifier of that IP and thereby be close enough to assign accusative Case to it.\marginnote[0pt]{Note, however, that \emph{believe} c-selects a CP \underline{finite} clause: \emph{She believed that pigs fly.}} If the correspondence between the IP/CP distinction and the clauses that can and can't house \textsc{pro} is correct, then the hypothesis that the infinitival complement of \emph{believe} is an IP, not a CP, will also explain why \textsc{pro} cannot appear in these infinitives.
\begin{exe}
	\raggedright
	\ex[*]{\label{ex3:nop} Bart believed [\IP \textsc{pro} to be happy].}
\end{exe}
Unlike \emph{try}, then, passivizing \emph{believe} causes the position that \emph{him} resides in in (\ref{ex3:believeill}a) to no longer be Case marked, and this is why A Movement is possible.

The subject position of the control infinitive following \emph{try}, by contrast, is Case marked by something other than \emph{try} and so Argument movement from that position will be blocked in both the Active and Passive versions of \emph{try}.
\begin{exe}
\ex \label{ex3:tryill} 
\begin{xlist}
	\ex[*]{\linkto Bart \under tried [\CP [\IP \from{\ \ } to be happy ]].}

	\ex[*]{\linkto Bart \under was tried [\CP [\IP \from{\ \ } to be happy]].}
\end{xlist}
\end{exe}
If we adopt this account, the subject position of control infinitives is Case-marked and since these are the positions that \textsc{pro} resides in, we can let \textsc{pro} be subject to the Case filter.

What's mysterious under this account, though, is why overt DPs cannot stand in the subject position of control infinitives. If these positions are Case-marked, then we should expect not just \textsc{pro} to be able to reside in them; overt DPs should be able to as well. But, as we've seen, this is not possible; examples like (\ref{ex3:noill}) are ungrammatical.
\begin{exe}
\ex \label{ex3:noill} 
\begin{xlist}
	\ex[*]{She to sing is swell.}
	\ex[*]{Elena tried her to sing.}
\end{xlist}
\end{exe}
A popular solution to this problem is proposed in \citet{Chomsky:1993}. They suggest that the Specifier of a control infinitive is assigned a Case that only \textsc{pro} may bear. English, they speculate, has a fourth Case, one whose morphological manifestation only \textsc{pro} is capable of achieving. They call this the ``null Case.'' This would have the desired effect of blocking overt DPs from the subject position of control infinitives. It could also be made to derive the \textsc{pro} Restriction. If the only place that null Case is assigned is to the Specifiers of control IPs, and if we not only restrict null Case to \textsc{pro} but prevent \textsc{pro} from bearing any other Case, then the Case filter alone will properly segregate the positions of \textsc{pro} and overt DPs. Overt DPs will have to surface in a nominative, accusative or genitive Case-marked position, and \textsc{pro} will have to surface in a position marked with null Case. \citet{Martin:1992,Martin:2001} argues that the availability of null Case can be tied to the semantics of the infinitival \emph{to} that comes with control infinitives. He argues that the \emph{to} which heads control infinitives has an interpretation as a kind of tense, whereas the \emph{to} that heads the infinitives that \textsc{pro} cannot appear in does not have a tense interpretation. If correct, this would suggest that null Case is assigned under conditions rather like those that nominative Case is assigned. Nominative Case is assigned to the Specifier position of a finite \obar{I}; on Martin's view, null Case is assigned to the Specifier of another type of tense \obar{I}, the one that appears in control infinitives.

There are problems with both the idea that there is a special Case reserved for \textsc{pro} and the view that this special Case is tied to a particular interpretation that distinguishes control infinitives from others. \citet{Pires:2001, Pires:2006}, \citet{Terzi:1992}, and \citet{Roussou:2001} (among others) have shown that the tensed interpretation Martin claims to correlate with the presence of \textsc{pro} doesn't. Therefore, the supposed Case that only \textsc{pro} can receive does not seem to be predictable from other properties of the infinitives it is alleged to be present in. Moreover, there is evidence that the subject position of control infinitives is capable of being assigned the normal Cases that can be borne by overt DPs. This evidence comes from Icelandic.

In Icelandic, floated quantifiers can be marked for Case and when they are, they have the same Case that the subjects they are related to do. In (\ref{ex3:icepro1}), for instance, the subject \emph{Str\'akarnir} (`the boys') is inflected for nominative Case, and so is the floated quantifier \emph{allir} (`all').
\begin{exe}
\ex \label{ex3:icepro1} \gll Str\'akarnir komust allir \'i sk\'ola. \\ {the boys.nom} got all.nom to school\\ `The boys all managed to get to school.'\\
\hspace*{\fill} \citep[(6a): 331]{Sigurdhsson:1991}\hspace*{.5in}
\end{exe}
The rules of Case assignment differ interestingly from English. Sometimes, the DPs in Specifiers of finite IPs can appear in something other than nominative Case. This seems to depend on the verb involved. For instance, the subject argument of the verb \emph{vanta{\dh}i} (`lacked') appears in the accusative Case, and the subject argument of \emph{leiddist} (`bored') shows up in the dative Case. In these situations, a floated quantifier related to the subject will inflect for the Case that the subject bears. This is illustrated in (\ref{ex3:icepro2}).
\begin{exe}
\ex \label{ex3:icepro2} 
\begin{xlist}
	\ex \gll Str\'akana vanta{\dh}i alla \'i sk\'olann. \\ {the boys.acc} lacked all.acc in {the school}\\ `The boys were all absent from school.'
	\ex \gll Str\'akunum leiddist \"ollum \'i sk\'ola. \\ {the boys.dat} bored all.dat in school\\ `The boys were all bored in school.'
\end{xlist}
\hspace*{\fill} \citep[(6b,c): 331]{Sigurdhsson:1991}\hspace*{.5in}
\end{exe}
In Icelandic, then, we can tell what Case a subject DP has by looking at its floated quantifier; and subject arguments have a wider range of Case options than they do in English. We can use this feature of Icelandic to determine what Case, if any, \textsc{pro} bears. What we find is that \textsc{pro} seems able to bear whatever Case an overt DP would bear in that position.\marginnote[0pt]{But PRO is capable of bearing other Cases as well, and so they do not behave precisely like overt DPs in this regard.} The examples in (\ref{ex3:icepro3}) show this correspondence.
\begin{exe}
\ex \label{ex3:icepro3} \gll Str\'akarnir vonast til \ldots \\ {the boys.nom} hope for \ldots\\ `the boys hope \ldots'
\begin{xlist}
	\ex \gll {[\CP} a{\dh} \textsc{pro} komast allir \'i sk\'ola ]. \\ {[\CP} that \textsc{pro} get all.nom to school ]\\ `to all manage to get to school.'
	\ex \gll {[\CP} a{\dh} \textsc{pro} vanta ekki alla \'i sk\'olann ]. \\ {[\CP} that \textsc{pro} lack not all.acc in school ]\\ `to all not be absent from school.'
	\ex \gll [\CP a{\dh} \textsc{pro} lei{\dh}ast ekki \"ollum \'i sk\'ola ]. \\ [\CP that \textsc{pro} bore not all in school ]\\ `to not all be bored in school.'
\end{xlist}
\hspace*{\fill} \citep[(8a--c): 331]{Sigurdhsson:1991}\hspace*{.5in}
\end{exe}
The controller---\emph{Str\'akarnir} (`the boys')---is marked with nominative Case in each of these examples. But the quantifier floated in the control infinitive inflects in a way that depends on the verb of the infinitival clause. When the verb is the kind that makes its subject appear in the accusative or dative Case, in finite clauses, it makes the floated quantifier in the control infinitives inflect for the accusative or dative Case. If the match between the Case borne by a floated quantifier and the Case borne by the subject it is related to holds in infinitival clauses as well, then we can conclude from these examples that \textsc{pro} bears the nominative, accusative or dative Cases assigned to Specifiers of finite IP in Icelandic.

Using null Case, then, to segregate \textsc{pro} from overt DPs does not seem attractive. The evidence we have does not suggest that \textsc{pro} and overt DPs differ with respect to the Cases they can bear. At present, it seems that the \textsc{pro} Restriction, and the ban against overt DPs in the subject position of control infinitives, must be sought elsewhere.

\subsection{\textsc{pro} and Control: A Summary}

The evidence we've reviewed in this section, and the explanations of it that we've entertained, are complex and open-ended. It might be useful to arrange in one place some of the conclusions we've reached and hypotheses we've looked at. 

We've discovered that there are silent argument DPs---``\textsc{pro}''---which have a semantic interpretation rather like \emph{one}. It shows up in infinitival clauses that function as subjects and as indirect questions; (\ref{ex3:summ1}) are examples.
\begin{exe}
\ex \label{ex3:summ1} 
\begin{xlist}
	\ex \textsc{pro} to eat durian is ill-advised.
	\ex I wondered how \textsc{pro} to eat durian.
\end{xlist}
\end{exe}
These infinitival clauses can be CPs, as the presence of the wh-phrase in (\ref{ex3:summ1}b) shows. In closely related languages, like Icelandic, it's possible to show that they are always CPs. We've entertained the hypothesis that these infinitives are always CPs in English as well, and this requires the conclusion that the complementizers of these CPs are phonetically silent. Let's call these infinitives: ``\textsc{pro}-infinitives.''

The existence of \textsc{pro} means that there are forces which restrict its occurrence to just the places where we have evidence for it. This seems to be the Specifier of the IPs in \textsc{pro}-infinitives. These same positions do not seem to allow overt DPs, as the contrast between (\ref{ex3:summ2}) and (\ref{ex3:summ1}) indicates.
\begin{exe}
\ex \label{ex3:summ2} 
\begin{xlist}
	\ex[*]{She to eat durian is ill-advised.}
	\ex[*]{I wondered how him to eat durian.}
\end{xlist}
\end{exe}
We have entertained the hypothesis that what makes the Specifier positions of \textsc{pro}-infinitives special in this way has to do with the Case they are assigned. But so far as I can tell, there is no evidence that these Specifiers get assigned a special Case, one that would make it different from the Cases assigned to other Specifiers of IPs. I suggest we abandon this hypothesis, therefore.

So, putting these parts together, we have:% \enlargethispage*{22pt}
\begin{exe}
\ex \label{ex3:summlist1} 
\begin{xlist}
	\ex There is a silent DP: \textsc{pro}.
	\ex \textsc{pro} can get a semantic interpretation like that assigned to \emph{one}: this is called its ``arbitrary'' interpretation.
	\ex \textsc{pro} can only appear in the Specifiers of \textsc{pro}-infinitives.
	\ex Overt DPs cannot appear in the Specifiers of \textsc{pro}-infinitives.
	\ex \textsc{pro}-infinitives can be CPs. It's possible that they are always CPs.
\end{xlist}
\end{exe}

We've also encountered infinitival clauses that do not contain a \textsc{pro} with an arbitrary reading, but instead have a silent subject that refers to the same thing that an argument spoken elsewhere does; (\ref{ex3:summ3}) is an example.
\begin{exe}
\ex \label{ex3:summ3} Andre tried to sing.
\end{exe}
In these cases, the overt DP `controls' the understood subject of the infinitive. We have two ideas about how to characterize these situations.

On the first view, control is a special instance of Argument movement, and so the syntax of (\ref{ex3:summ3}) is identical to that posited for (\ref{ex3:summ4}).
\begin{exe}
\ex \label{ex3:summ4} Andre seemed to sing.
\end{exe}
In both cases, \emph{Andre} moves from the position \emph{sing} assigns a \mbox{$\theta$-role} to. In (\ref{ex3:summ4}) this is the only \mbox{$\theta$-role} \emph{Andre} receives, as the verb \emph{seem} does not assign an external \mbox{$\theta$-role}. In (\ref{ex3:summ3}), by contrast, \emph{Andre} also receives a \mbox{$\theta$-role} from \emph{try}, presumably by moving into the position this \mbox{$\theta$-role} is assigned to. If this is the correct treatment of Control, then our present formulation of the Theta Criterion will have to be changed so that it allows on argument to bear more than one \mbox{$\theta$-role}. This account of Control will also have to find a way of capturing the difference between (\ref{ex3:summ3}) and (\ref{ex3:summ5}).
\begin{exe}
\ex[*]{ \label{ex3:summ5} Andre was tried to sing.}
\end{exe}
And, if the Specifiers of these infinitives are Case-marked, as the evidence from Icelandic suggests, then (\ref{ex3:summ3}) would seem to be a violation of the prohibition against Argument movement from Case-marked positions. So, this characterization of Control includes the points in~(\ref{ex3:summlist2}).
\begin{exe}
\ex \label{ex3:summlist2} If Control is Argument movement, then:
\begin{xlist}
	\ex \textsc{pro}-infinitives are prevented from arising in contexts where Control is obligatory.
	\ex The Theta Criterion requires changing.
	\ex There is an unknown force responsible for the contrast in:
	\begin{xlist}
	\ex[]{Andre tried to sing.}
	\ex[*]{Andre was tried to sing.}
\end{xlist}
	\ex If the Specifiers of the IPs from which Control arguments move is Case marked, then the generalization that Argument movement from Case-marked position is ungrammatical is lifted in these contexts.
\end{xlist}
\end{exe}

On the second view, Control is a special meaning that \textsc{pro} is able to receive. If this is correct, those infinitives in which Control arises are \textsc{pro}-infinitives. This characterization of Control brings with it the following points.
\pagebreak
\begin{exe}
\ex \label{ex3:summlist3} If Control is found in \textsc{pro}-infinitives, then:
\begin{xlist}
	\ex \textsc{pro} cannot get an arbitrary reading in those contexts where Control is obligatory.
	\ex There is a semantic rule independent of Argument movement which gives the Control reading.
	\ex If the Specifiers of \textsc{pro}-infinitives are Case-marked, then the ban against Argument movement from Case-marked positions can be used to account for the contrast between:
	\begin{xlist}
	\ex[]{Andre tried to sing.}
	\ex[*]{Andre was tried to sing.}
\end{xlist}
\end{xlist}
\end{exe}

These two approaches to the Control relation are presently hotly contested. I will try to stay neutral between the two views and their consequences in what follows. But if it becomes necessary to choose, I will pick the second view.

\section{The Theta Criterion Reconsidered} % (fold)
\label{sec:reconsidered}

In the previous chapter, we encountered the need to design a syntax that guides, or at least describes, how the semantics of arguments and argument-takers happens. Following the language used in \citet{Chomsky:1981}, we divided this task into two parts. One part simply records the fact that the connection between the arguments in a sentence and the word that these phrases are arguments of must be made. This is the goal of the Theta Criterion.
\begin{exe}
	\raggedright
	\ex \label{ex3:tc2} Theta Criterion
	\begin{xlist}
		\ex Every $\theta$-role must be assigned to a unique syntactic position.
		\ex Every argument XP must be related to exactly one position that is assigned a $\theta$-role.
	\end{xlist}
\end{exe}
The second part is a linking theory, which determines where the syntactic positions are that the $\theta$-roles assigned by some word are. This is the goal of the Projection Principle.
\begin{exe}
	\raggedright
	\ex \label{ex3:pp2} The Projection Principle
	\begin{xlist}
		\ex If $\beta$ is c-selected by $\alpha$, then there must be a projection of $\alpha$, $\pi$, that immediately dominates $\beta$ and if $\pi$ immediately dominates some other phrase, $\gamma$, then $\gamma$ must be another projection of $\alpha$ or part of another c-selected argument of~$\alpha$.
		\ex The Specifier of $\alpha$ is the position to which the external $\theta$-role of $\alpha$ is assigned.
	\end{xlist}
\end{exe}
What we have seen in this chapter is that sentences may satisfy these conditions even when argument XPs are not spoken in $\theta$-marked positions. These cases arise when movement rules put an argument XP into a different position than the one to which their $\theta$-role is assigned. The Theta Criterion and Projection Principle are presented as conditions on D-structure parses, then, and these can differ from S-structure parses in just the ways that Complex NP Shift and Argument Movement permit or require.

What we've just seen in the discussion of the differences in the raising and control infinitives, however, requires a slightly different view of the second clause of the Theta Criterion. If control infinitives do not arise as a result of A-movement, then this is blocked by preventing an Argument from occupying two $\theta$-marked positions in a derivation. On this way of describing things, then, the Theta Criterion is a condition on derivations, not on D-structures. When I presented the Theta Criterion, I presented it as a description of a semantic effect. Its source might come from how the meanings of arguments and argument-takers are combined. Our model makes D-structures the parse that the semantics interprets. There is something inconsistent about the way the Theta Criterion's effect on constraining derivations is described in the previous sections and our model of derivations presently. 

I want to resolve that problem in this section, and I want to take a few steps in the direction of understanding better what the Theta Criterion describes. There are some other problems that were left hanging when this condition was introduced that I'd like to address here.

The idea behind $\theta$-roles is that they represent specialized relations that are part of the mix of meanings that a verb brings to a sentence. That these specialized relations exist has sometimes been argued to afford a handle on a linking theory that can predict where an argument of some verb will be based on the kind of relation it has to the verb's meaning. Recall, for instance, that Ed Keenan argues that there is a cross-linguistic correlation between being an ``agent'' of some action and being the external argument of the verb that describes that action. Claims of this kind require, of course, an independent definition of the relations that $\theta$-roles encode. To test Keenan's hypothesis requires knowing when some argument is an ``agent'' of an action in a way that is independent of its being an external argument. Working definitions of $\theta$-roles, however, have proven largely unsuccessful, and so determining that they have a contribution to the meanings of sentences is still not established.\marginnote[0pt]{See \citet{Jackendoff:1972} for an initial, and searching, attempt to establish an inventory of $\theta$-roles, and \citet{Dowty:1979} for an excellent examination of the problems attendant with defining $\theta$-roles.} Nonetheless, the existence of $\theta$-roles remains a popular working hypothesis. Let us assume their existence.

When we learn the meaning of a verb, we can decompose that meaning into two parts. There is the type of argument that the verb selects, and there is the action or state of affairs involving that argument that the verb describes. The language of $\theta$-roles is generally reserved for the method by which verbs combine with DP, or sometimes PP, arguments. I'll represent the meaning of a verb as a list of these components. To do that, I'll introduce very informally a metalanguage that we'll use to express meanings.

The central ingredient of this metalanguage is a method of building a function. A function relates one set to another set by mapping each element of the first set to exactly one element of the second. So, for instance, we can define a function, let's call it $g$, that relates the set \{a, b, c\} to the set \{1, 2, 3\} in the way indicated in~(\ref{ex3:g1}).
\begin{exe}
	\raggedright 
	\ex \label{ex3:g1} The function $g$:\\
	a $\rightarrow$ 1\\ b $\rightarrow$ 2\\ c $\rightarrow$ 3
\end{exe}
We can represent this function $g$ with the set of ordered pairs in (\ref{ex3:g2}), where the first member of each pair is from the first set (the ``domain'' of the function) and the second member of each pair is the element of the second set that $g$ relates the first element to. (The second set is known as the ``range'' of the function.)
\begin{exe}
	\raggedright
	\ex \label{ex3:g2} 
	$g$ = \{ \tuple{a,1}, \tuple{b,2}, \tuple{c,3} \}
\end{exe}
A third way of representing this function is in~(\ref{ex3:g3}).
\begin{exe}
	\raggedright
	\ex \label{ex3:g3} 
	$g$(a) = 1\\ $g$(b) = 2\\ $g$(c) = 3
\end{exe}
We say that $g$ applies to its argument a and gives 1, applies to its argument b to give 2, and applies to its argument c to give 3. 

We can use the language of mathematics to describe functions that relate numbers, as well as other mathematical objects, to other numbers. For instance, the expression in (\ref{ex3:f1}) describes a function that accepts numbers---let's restrict the domain to positive whole numbers---and maps that number onto its square.
\begin{exe}
	\raggedright
	\ex \label{ex3:f1} 
	$x^2$, where $x$ is a positive whole number
\end{exe}
We understand $x$ here to be an element of the domain set and the ``$^2$'' to be a function that maps $x$ onto its square. The use of the variable $x$ is a compact way of representing a function that relates large sets. In (\ref{ex3:f1}), the domain and range are sets of infinite cardinality so representing the function by way of $^2$ is the only feasible method. This method requires an understanding of what a variable is and a way of defining ``$^2$.'' Informally, ``$^2$'' is defined as a function that multiplies its argument, $x$, to itself. %
$
$\marginnote[0pt]{Of course, to understand this definition of $^2$ requires defining ``multiplies,'' and this will in turn require a definition of addition and whole numbers. But we wish to get back to linguistics, so here is where we stop.}$
$
The function expressed by (\ref{ex3:f1}) can be expressed in the format that (\ref{ex3:g3}) illustrates with~(\ref{ex3:f11}).
\begin{exe}
\raggedright
	\ex \label{ex3:f11}
	\textsc{sqr}($x$) = $x \times x$, where $x$ is a positive whole number
\end{exe}

It is possible to define functions whose range is another function. For instance, (\ref{ex3:f2}) is a compact way of representing a function whose domain is the set \{ $2$, $3$, $4$, $\ldots$ \} and whose range is a function $x^n$ where $n$ is the argument of the function defined in~(\ref{ex3:f2}).
\begin{exe}
	\raggedright
	\ex \label{ex3:f2}%
	%
	\marginnote[0pt]{Informally, ``$^{n}$'' can be defined as the function that multiplies its argument, $x$, to itself $n-1$ times.}%
	%
	$x^n$, where $n$ is a positive whole number greater than 1
\end{exe}
With this understanding, the function in (\ref{ex3:f2}) can be represented by the series in~(\ref{ex3:f3}).
\begin{exe}
	\raggedright
	\ex \label{ex3:f3} 
	$x^n$(2) = $x^2$ = $x \times x$\\ $x^n$(3) = $x^3$ = $x \times x \times x$\\ $x^n$(4) = $x^4$ = $x \times x \times x \times x$\\ \hspace*{50pt}$\vdots$
\end{exe}

When functions return other functions, we will want to keep track of how the values of the variables within them are assigned. This gives us a way of determining the function in the range returned by the first function. A way of doing this involves an operator, $\lambda$, defined informally in~(\ref{ex3:lambda1}).
\begin{exe}
	\raggedright
	\ex \label{ex3:lambda1} 
	$\lambda x$ $f$($x$) applied to a = $f$(a), where $f$ is a function.
\end{exe}
We say that the variable $x$ is bound by the operator $\lambda x$. With this notation, the function described in (\ref{ex3:f1}) is rendered as~(\ref{ex3:l1}).
\begin{exe}
	\raggedright
	\ex \label{ex3:l1} 
	$\lambda x$ $x^2$, where $x$ is a positive whole number
\end{exe}
This gives us the domain to range relations in~(\ref{ex3:l1a}).
\begin{exe}
	\raggedright
	\ex \label{ex3:l1a} 
	{}[$\lambda x$ $x^2$](1) = $1^2$ = 1\\
	{}[$\lambda x$ $x^2$](2) = $2^2$ = 4\\
	{}[$\lambda x$ $x^2$](3) = $3^2$ = 9\\
	\hspace*{40pt} $\vdots$
\end{exe}
The function in (\ref{ex3:f3}) is rendered as~(\ref{ex3:l2}).
\begin{exe}
	\raggedright
	\ex \label{ex3:l2} 
	$\lambda n$ $\lambda x$ $x^n$, where $x$ is a positive whole number and $n$ is a positive whole number greater than 1
\end{exe}
This gives us the domain to range relations in~(\ref{ex3:l2a}).
\begin{exe}
	\raggedright
	\ex \label{ex3:l2a} 
	{} [$\lambda n$ $\lambda x$ $x^n$](2) = [$\lambda x$ $x^2$]\\ 
	{} [$\lambda n$ $\lambda x$ $x^n$](3) = [$\lambda x$ $x^3$]\\
	{} [$\lambda n$ $\lambda x$ $x^n$](4) = [$\lambda x$ $x^4$]\\
	\hspace*{40pt} $\vdots$
\end{exe}
The function that results from applying (\ref{ex3:l2}) to the number ``2'' is the one that (\ref{ex3:l1}) defines. One thing the $\lambda$ notation allows us to do is order the way the variables in a complex function are assigned to elements of the domain. This will be key to expressing the compositionality of syntactic phrases.

I'll adopt a way of presenting compactly cases where the range of one function is another function. To illustrate this method of presentation, consider the functions in~(\ref{ex3:illff}).
\enlargethispage*{15pt}
\begin{exe}
	\raggedright
	\ex \label{ex3:illff} 
	\begin{xlist}
		\ex $\lambda x\ P(x)$
		\ex $\lambda y\ Q(y)$
	\end{xlist}
\end{exe}
Suppose that (\ref{ex3:illff}a) has the domain of positive whole numbers ($x$ ranges over \{1,2,3\ldots\}) and (\ref{ex3:illff}b) has the domain of lower case roman letters ($y$ ranges over \{a,b,c,\ldots\}). Now consider the function in~(\ref{ex3:illff2}), which puts these two functions together.
\begin{exe}
	\raggedright
	\ex \label{ex3:illff2} 
	$\lambda x$ $\lambda y\ Q(y)+P(x)$ 
\end{exe}
This is a function whose domain is positive whole numbers and whose range is the function in (\ref{ex3:illff}b). It could apply to a positive whole number, such as ``$2$,'' and yield the function shown in~(\ref{ex3:illff3}).
\begin{exe}
	\raggedright
	\ex \label{ex3:illff3} 
	{} [$\lambda x$ $\lambda y\ Q(y)+P(x)$]($2$) = $\lambda y \ Q(y)+P(2)$
\end{exe}
This function has letters as its domain, so it can be fed something like ``b'' and give the resulting expression in~(\ref{ex3:illff4}).
\begin{exe}
	\raggedright
	\ex \label{ex3:illff4} 
	{} [$\lambda y\ Q(y)+P(2)$](b) = $Q$(b)+$P(2)$
\end{exe}
Our present notation allows us to represent this situation with~(\ref{ex3:illff5}).
\begin{exe}
	\raggedright
	\ex \label{ex3:illff5} 
	{} [ [$\lambda x$ $\lambda y\ Q(y)+P(x)$]($2$) ](b)
\end{exe}
The function in (\ref{ex3:illff2}) applies to ``$2$'' and forms a new function that applies to ``b.'' The compact way of representing this that I will adopt here is~(\ref{ex3:illff6}).\marginnote[0pt]{Another common convention, found often in the writing of logicians, is to represent (\ref{ex3:illff5}) with (i).
\begin{exe}
	\raggedright
	\exi{i.} {}[$\lambda x$ $\lambda y\ Q(y)+P(x)]$(b,$2$)
\end{exe}
Note that the order of arguments is reversed in this notation. Logicians frequently get things backwards.}
\begin{exe}
	\raggedright
	\ex \label{ex3:illff6} 
	{}[$\lambda x$ $\lambda y\ Q(y)+P(x)]$($2$)(b)
\end{exe}

The domain of a function can also be a set of functions. The function in~(\ref{ex3:complex}) illustrates.
\begin{exe}
	\raggedright
	\ex \label{ex3:complex} 
	$\lambda P$ $\lambda y$ 3 + $P(y)$, where $y$ is a positive whole number and $P$ a function whose domain is positive whole numbers and whose range is positive whole numbers.  
\end{exe}
The function defined in (\ref{ex3:f1}) has a domain and range of positive whole numbers, so it is in the domain of (\ref{ex3:complex}). Thus, (\ref{ex3:complexf}) is well-formed.
\begin{exe}
	\raggedright
	\ex \label{ex3:complexf} 
	{}[$\lambda P$ $\lambda y$ 3 + $P(y)$]($\lambda x$ $x^2$) =\\
	$\lambda y$ 3 + [$\lambda x$ $x^2$]($y$) =\\
	$\lambda y$ 3 + $y^2$
\end{exe}
Note that the second line of these equalities is equal to the last line because the $y$ is an argument of ``$\lambda x$ $x^2$,'' and so the function can apply to that argument to yield $y^2$. The domain of this function is the set of positive whole numbers, so this describes the relations in~(\ref{ex3:addf}).
\begin{exe}
	\raggedright
	\ex \label{ex3:addf} 
	{}[$\lambda y$ 3 + $y^2$](1) = 3 + 1^2 = 4\\
	{}[$\lambda y$ 3 + $y^2$](2) = 3 + 2^2 = 7\\
	{}[$\lambda y$ 3 + $y^2$](3) = 3 + 3^2 = 12\\
	\hspace*{50pt} $\vdots$
\end{exe}

\label{lambdalaw}Because functions can take other functions as arguments, it is necessary to legislate how the variables in the function that is an argument can be manipulated by the function that takes that argument. The law is that no variable in a function can be bound by the $\lambda$ operator of another function. So neither of the following cases are well-formed.
\pagebreak
\begin{exe}
	\raggedright
	\ex \label{ex3:badl} 
	\begin{xlist}
		\ex {} [$\lambda P \lambda x$ $P$] ($x^2$) $\neq$ $\lambda x$ $x^2$
		\ex {} [$\lambda P \lambda x$ $P$] ($\lambda x\ x^2$) $\neq$ $\lambda x$ $\lambda x\ x^2$
	\end{xlist}
\end{exe}
The unbound variable (should such things exist) in the argument of (\ref{ex3:badl}a) cannot become bound in the function it is an argument of. And the bound variable in the argument of (\ref{ex3:badl}b) cannot acquire an additional binder in the function it is an argument of. To get a variable inside one function to be bound by the $\lambda$ operator of another function, the technique illustrated by (\ref{ex3:complexf}) must be employed.

Let me illustrate that method more generally. Consider the two functions in~(\ref{ex3:nofree1}).
\begin{exe}
	\raggedright
	\ex \label{ex3:nofree1} 
	\begin{xlist}
		\ex $\lambda x\ P(x)$
		\ex $\lambda M$ $\lambda y\ M(y)$ 
	\end{xlist}
\end{exe}
Suppose that (\ref{ex3:nofree1}a) is in the domain of (\ref{ex3:nofree1}b). This allows (\ref{ex3:nofree1}b) to apply to (\ref{ex3:nofree1}a) as in~(\ref{ex3:nofree2}) .
\begin{exe}
	\raggedright
	\ex \label{ex3:nofree2} 
	{}[$\lambda M$ $\lambda y\ M(y)$]($\lambda x P(x)$) = $\lambda y\ [\lambda x\ P(x)](y)$
\end{exe}
If $y$ is in the domain of $\lambda x\ P(x)$ (i.e., the function in (\ref{ex3:nofree1}a)), then it can apply to $y$ and this results in~(\ref{ex3:nofree3}).
\begin{exe}
	\raggedright
	\ex \label{ex3:nofree3} 
	$\lambda y\ P(y)$
\end{exe}
In this way, then, the variable inside the function in (\ref{ex3:nofree1}a) has become bound by $\lambda y$ in the function in~(\ref{ex3:nofree1}b).

Now that the $\lambda$ operator has been introduced, we are at last ready to describe how the component parts of the meaning of verbs can be modeled. Recall that we have decided that there are two components to the meaning of a verb. One component describes the state of affairs that the verb refers to and the other lists the relations denoted by the $\theta$-roles, if any, associated with that verb. Let's call the states of affairs that verbs refer to ``eventualities.'' %
%
\marginnote[-30pt]{I am describing here a way of rendering the meaning of verbs that originates with \citet{Davidson:1967}. Davidson called the states of affairs that verbs refer to ``events.'' The term ``eventualities'' comes from \citet{Bach:1986}, who offers this term to include both events, which can be conceived of as dynamic actions, and states, which are static. For an early argument on behalf of importing events into compositional semantics, see \citet{Higginbotham:1985}.}%
%
Using this idea, the meaning of (\ref{ex3:kiyomi}a) can be paraphrased with~(\ref{ex3:kiyomi}b).
\begin{exe}
	\raggedright
	\ex \label{ex3:kiyomi} 
	\begin{xlist}
		\ex Kiyomi danced.
		\ex Kiyomi was the agent of a dancing eventuality
	\end{xlist}
\end{exe}
We'll begin by defining functions that correspond to the ``agent'' and ``dancing'' parts of the paraphrase in~(\ref{ex3:kiyomi}).
\begin{exe}
	\raggedright
	\ex \label{ex3:parap}%
	%
	\marginnote[0pt]{In the semantic literature, the more neutral word ``entities'' is used instead of ``things.''}%
	%
	\begin{xlist}
		\ex \textsc{agent}($x)(\epsilon$) is a function from things, $x$, to a function whose domain is eventualities, $\epsilon$. The range of \textsc{agent} is \{True,False\}. When $x$ is the agent of $\epsilon$, then \textsc{agent}($x$)($\epsilon$) $\rightarrow$ True. Otherwise, \textsc{agent}$(x)(\epsilon$) $\rightarrow$ False. 
		\ex \textsc{dance}($\epsilon$) is a function from eventualities, $\epsilon$, to \{True,False\}. When $\epsilon$ is a dancing, then \textsc{dance}($\epsilon$) $\rightarrow$ True, otherwise \textsc{dance}($\epsilon$) $\rightarrow$ False.
	\end{xlist}
\end{exe}
Loosely speaking, what \textsc{agent} and \textsc{dance} do is describe the conditions under which an eventuality has an agent or is a dancing. They do that by spelling out when the \textsc{agent} and \textsc{dance} functions map to truth or falsity. A perhaps more transparent way of rendering this informally is with the paraphrases in~(\ref{ex3:parapp}).
\begin{exe}
	\raggedright
	\ex \label{ex3:parapp} 
	\begin{xlist}
		\ex \textsc{agent}($x)(\epsilon$)=T $\approx$ $x$ is Agent of $\epsilon$
		\ex \textsc{dance}($\epsilon$)=T $\approx$ $\epsilon$ is a dancing
	\end{xlist}
\end{exe}
In reading the formulas that represent denotations in what follows, you can safely make this substitution.

The meaning of the word \emph{dance} should describe those eventualities that make the functions in (\ref{ex3:parap}) true. We can model this with the function in~(\ref{ex3:deno1r}).\marginnote[0pt]{I abbreviate ``Truth'' with ``T'' in (\ref{ex3:deno1r}). Another common abbreviation is ``$1$''; ``$1$'' abbreviates Truth and ``$0$'' abbreviates False.}
\begin{exe}
	\raggedright
	\ex \label{ex3:deno1r} 
	\denotes{dance} = $\lambda x$ $\lambda \epsilon$ [\underline{\textsc{agent}}($x$)($\epsilon$) $\wedge$ \textsc{dance}($\epsilon$)]=T, where\\ 
	$x$ is a thing and $\epsilon$ is an eventuality.
\end{exe}
The underline signals which $\theta$-role is external; the only $\theta$-role \emph{dance} provides is external. The brackets ``\denotes{\ }'' represent a function that takes the material the brackets enclose and returns the denotation of that material. In (\ref{ex3:deno1r}), the denotation brackets apply to \emph{dance} and return the meaning of \emph{dance}, which is represented by the function to the right of the ``='' sign.

Consider how this meaning of \emph{dance} can be used to form the meaning of the VP in (\ref{ex3:kiyomi}a), which is given in~(\ref{ex3:kiyomivp}).%
%
\marginnote[0pt]{I have dropped the ``=T'' part of the function. Let's understand the functions we are interested in to always be those that map to~T.}%
%

\begin{exe}
	\raggedright
	\ex \label{ex3:kiyomivp} 
	\small
	\begin{forest}
	[$\lambda \epsilon$ \textsc{agent}(\denotes{Kiyomi})($\epsilon$) $\wedge$ \textsc{dance}($\epsilon$)\\ vP
		[DP	[Kiyomi, roof] ]
		[$\lambda x$ $\lambda \epsilon$ \textsc{agent}($x$)($\epsilon$) $\wedge$ dance($\epsilon$)\\ \ibar{V} [$\lambda x$ $\lambda \epsilon$ \textsc{agent}($x$)($\epsilon$) $\wedge$ dance($\epsilon$)=T \\ \obar{V} [dance]] ]
	]
	\end{forest}
\end{exe}
To achieve this, we will need a set of semantic rules that take the meanings of the parts of a phrase and produce from them the meaning of that phrase.%
%
\marginnote[0pt]{These are the rules that ensure semantic compositionality for phrases.} Adopt the following two semantic rules.%
%
\begin{exe}
	\raggedright
	\ex \label{ex3:perco} Rules of Semantic Composition
	\begin{xlist}
		\ex If X is the sole daughter of Z, then \denotes{Z} = \denotes{X}
		\ex If X and Y are daughters of Z, then \denotes{Z} is the result of applying the function \denotes{X} to its argument~\denotes{Y}.
	\end{xlist}
\end{exe}
I've written on the nodes of (\ref{ex3:kiyomivp}) the denotations derived from \denotes{dance} by these rules. %
%
\marginnote[-10pt]{However, I've not written out what the denotation of \emph{Kiyomi} is. If I had done so, I'd have to find a way of expressing how the name \emph{Kiyomi} allows this expression to refer to a particular person whose name is Kiyomi.}%
%
The meaning assigned to the VP is a function from eventualities that are dancings and that have Kiyomi as agent to truth. This is roughly what we want. To get from this the meaning of (\ref{ex3:kiyomi}a), we'll have to work out the semantic contribution of the past tense, and figure out how to get this sentence to refer to an eventuality of the sort that this VP describes. These parts of the meaning of a sentence do not come from the verb and its arguments, however, so we can safely set them aside.

Consider next the meaning of a verb that has two $\theta$-roles, like \emph{eat}.
\begin{exe}
	\raggedright
	\ex \label{ex3:deno2} 
	\denotes{eat} = $\lambda x$ $\lambda y$ $\lambda \epsilon$ \underline{\textsc{agent}}($y$)($\epsilon$) $\wedge$ \textsc{patient}($x$)($\epsilon$) $\wedge$ \textsc{eat}($\epsilon$)=T, where
	\begin{xlist}
	\ex $x$ and $y$ are things and $\epsilon$ is an eventuality
	\ex \textsc{patient} c-selects DP
	\end{xlist}
\end{exe}
Recall that internal arguments are c-selected by the verbs that take them, so I have added this information in (\ref{ex3:deno2}). I've expressed this as a requirement imposed by the $\theta$-role whose denotation is the function that takes the relevant expression as its argument. The c-selection requirement of \textsc{patient} can be thought of as part of the definition of this function's domain. It specifies what the morpho-syntactic properties of the argument for \textsc{patient} must be. 

Our system derives the parse, and meanings, for a VP headed by \emph{eat} as indicated in~(\ref{ex3:eatex}).
\begin{exe}
	\raggedright
	\ex \label{ex3:eatex} 
	\small
	\begin{forest}
	[$\lambda \epsilon$ \textsc{agent}(\denotes{she})($\epsilon$) $\wedge$ \textsc{patient}(\denotes{it})($\epsilon$) $\wedge$ \textsc{eat}($\epsilon$)\\ VP
		[DP [she, roof] ]
		[$\lambda y$ $\lambda \epsilon$ \textsc{agent}($y$)($\epsilon$) $\wedge$ \textsc{patient}(\denotes{it})($\epsilon$) $\wedge$ \textsc{eat}($\epsilon$)\\ \ibar{V} 
			[$\lambda x$ $\lambda y$ $\lambda \epsilon$ \textsc{agent}($y$)($\epsilon$) $\wedge$ \textsc{patient}($x$)($\epsilon$) $\wedge$ \textsc{eat}($\epsilon$)\\ \obar{V} [eat]]
			[DP [it, roof] ]
		]
	]
	\end{forest}
\end{exe}
The meaning assigned to this VP is a function from eventualities that maps those which are eating eventualities that have the referent of \emph{she} as its agent and the referent of \emph{it} as its patient to True. This is roughly what we want: the VP describes eating eventualities in which she consumed it.

Now that we have a semantics in front of us for verbs and their arguments, let's reconsider the Theta Criterion.
\begin{exe}
	\raggedright
	\exr{ex3:tc2} Theta Criterion
	\begin{xlist}
		\ex Every $\theta$-role must be assigned to a unique syntactic position.
		\ex Every argument XP must be related to exactly one position that is assigned a $\theta$-role.
	\end{xlist}
\end{exe}
Let's first focus on the first clause of the Theta Criterion. This clause has two consequences: it requires every $\theta$-role to be given an argument and it prevents any $\theta$-role from being given more than one argument. A problem for the first of these consequences is that there are numerous counter-examples. For instance, (\ref{ex3:eat1}) is perfectly grammatical and it has all the appearances of containing the same verb that (\ref{ex3:eatex}) does.
\begin{exe}
	\raggedright
	\ex \label{ex3:eat1} 
	Shai ate.
\end{exe}
Indeed, the meaning of the verb \emph{eat} in (\ref{ex3:eat1}) should entail that something was eaten---that there is a Patient of the eating eventuality---because the existence of this Patient can be presupposed in a conversation that includes (\ref{ex3:eat1}). This might be seen, for instance, by the appropriateness of B's response in~(\ref{ex3:eatconvo}).
\begin{exe}
	\raggedright
	\ex \label{ex3:eatconvo} 
	\begin{xlist}
		\exi{A:} Shai ate at the function last night.
		\exi{B:} Yeah, and he told me nothing tasted good.
	\end{xlist}
\end{exe}
Compare this to a conversation that starts with a verb that doesn't entail the existence of a Patient:
\begin{exe}
	\raggedright
	\ex \label{ex3:neatconvo} 
	\begin{xlist}
		\exi{A:} Shai danced at the function last night.
		\exi{B:} Yeah, and he told me nothing tasted good.
	\end{xlist}
\end{exe}
B's response in this conversation is a non-sequitur because no Patient of an eating eventuality is invoked by \emph{dance}. By contrast, the word \emph{nothing} in B's response in (\ref{ex3:eatconvo}) is understood to quantify over the Patients of the eating event invoked by \emph{eat}. Considerations such as these push towards the conclusion that there is a Patient $\theta$-role associated with the \emph{eat} in (\ref{ex3:eat1}). And yet, there is no argument associated with this $\theta$-role, in violation of the Theta Criterion.

This ability for some $\theta$-roles to not be assigned to arguments seems to vary with the verb. So whereas the Patient $\theta$-role need not be matched with an argument for \emph{eat}, the Patient $\theta$-role associated with \emph{build} must be matched with an argument, as the ungrammaticality of~(\ref{ex3:put}) indicates.
\begin{exe}
	\raggedright
	\ex[*]{\label{ex3:put} Cherlon built.}
	\sn \emph{compare:}
	\ex[]{Cherlon built them.}
\end{exe}
Moreover, this ability to violate the Theta Criterion seems to be available only for internal arguments. Alternations like (\ref{ex3:inimp}) are common, but so far as I know, alternations like (\ref{ex3:eximp}), where the Agent $\theta$-role remains but doesn't get assigned to an argument, aren't possible.\marginnote[0pt]{There are alternations of the sort that (\ref{ex3:eximp}) illustrates where the external $\theta$-role is lost. And there are also alternations of the sort (\ref{ex3:eximp}) exemplifies in which the verb has undergone some morphological change. Both of these are taken up in the next section. But finding the same verb with an external $\theta$-role it contains both assigned and not assigned doesn't seem to happen. We return to this issue in the final section of this chapter.}
\begin{exe}
	\raggedright
	\ex \label{ex3:inimp} 
	\begin{xlist}
		\ex[]{Shai ate nattoo.}
		\ex[]{Shai ate.}
	\end{xlist}
	\ex \label{ex3:eximp}
	\begin{xlist}
		\ex[]{Shai danced.}
		\ex[*]{It danced.} (\emph{it} an expletive)
	\end{xlist}
\end{exe}

One way of describing this situation is to allow for a process that creates denotations which existentially close one of the argument positions of a verb. This process would create from (\ref{ex3:deno2}), for instance,~(\ref{ex3:deno2b}).\marginnote[0pt]{Understand $\exists x\ P(x)=T$ to be equivalent in meaning to ``there is an x that is P.'' ``$\exists x\ P(x)$'' maps to T if and only if there is an a in its domain for which P(a)=T. ``$\exists$'' is called an ``existential'' operator.}
\begin{exe}
	\raggedright
	\exr{ex3:deno2} 
	\denotes{eat} = $\lambda x$ $\lambda y$ $\lambda \epsilon$ \underline{\textsc{agent}}($y$)($\epsilon$) $\wedge$ \textsc{patient}($x$)($\epsilon$) $\wedge$ \textsc{eat}($\epsilon$)=T
	\ex \label{ex3:deno2b}
	\denotes{eat} = $\lambda y$ $\lambda \epsilon$ $\exists x$ \underline{\textsc{agent}}($y$)($\epsilon$) $\wedge$ \textsc{patient}($x$)($\epsilon$) $\wedge$ \textsc{eat}($\epsilon$)=T
\end{exe}
Let's call this operation ``Existential Closure.''\marginnote{Existential Closure is modeled after an analysis in \citet{Chierchia:2004}.} As noted, Existential Closure doesn't seem to be able to apply to the external $\theta$-role of a verb.
\begin{exe}
	\raggedright
	\ex \label{ex3:noextc} Existential Closure cannot apply to an external $\theta$-role.
\end{exe}
Let's assume that Existential Closure is, like Passive, a rule that makes new verbs from existing ones. As we've seen, the choice of verb determines whether Existential Closure can apply. If there is a pattern that determines which verb is susceptible to Existential Closure or not, I don't know it.

The verb created by Existentially Closing \emph{eat} will produce a meaning for the VP it heads indicated in~(\ref{ex3:shaidt}).
\begin{exe}
	\raggedright
	\ex \label{ex3:shaidt} 
	\small
	\begin{forest}
	[$\lambda \epsilon$ $\exists x$ \textsc{agent}(\denotes{Shai})($\epsilon$) $\wedge$ \textsc{patient}($x)(\epsilon$) $\wedge$ \textsc{eat}($\epsilon$)\\ VP
		[DP [Shai, roof] ]
		[$\lambda y$ $\lambda \epsilon$ $\exists x$ \textsc{agent}($y)(\epsilon$) $\wedge$ \textsc{patient}($x)(\epsilon$) $\wedge$ \textsc{eat}($\epsilon$)\\ \ibar{V} [$\lambda y$ $\lambda \epsilon$ $\exists x$ \textsc{agent}($y)(\epsilon$) $\wedge$ \textsc{patient}($x)(\epsilon$) $\wedge$ \textsc{eat}($\epsilon$)\\ \obar{V} [eat]] 
		]
	]
	\end{forest}
\end{exe}
This is the right meaning for the VP; it says that this VP describes eventualities that are eatings, have a Patient, and have Shai as their Agent.

Now we have a way of explaining how the first clause of the Theta Criterion can be made consistent with cases where a $\theta$-role is present but there is no argument for that $\theta$-role in the syntax. These cases arise because of Existential Closure. When Existential Closure cannot apply---as with external $\theta$-roles and the internal $\theta$-roles of verbs like \emph{build}---the first clause of the Theta Criterion still has the consequences that were highlighted above: (i)~there must be an argument for the $\theta$-role in the syntax, and (ii)~there must not be more than one such argument. Irene Heim and Angelika Kratzer\cite{Heim:1998} argue that the first of these consequences---that there must be an argument in the syntax matched to the $\theta$-role of a verb---can be described in semantic terms under these assumptions. One way of doing that would be with~(\ref{ex3:vptype}).
\begin{exe}
	\raggedright
	\ex \label{ex3:vptype} The denotation of a VP must be a function whose domain is a set of eventualities.
\end{exe}
To see how this works, consider what happens if \emph{build} is put into a VP that does not contain an argument for its Patient $\theta$-role.
\begin{exe}
	\raggedright
	\ex \label{ex3:nobuild} 
	\small
	\begin{forest}
	[$\lambda y$ $\lambda \epsilon$ \textsc{agent}($y)(\epsilon$) $\wedge$ \textsc{patient}(\denotes{Cherlon})($\epsilon$) $\wedge$ \textsc{build}($\epsilon$)\\ VP
		[DP [Cherlon, roof] ]
		[$\lambda x$ $\lambda y$ $\lambda \epsilon$ \textsc{agent}($y)(\epsilon$) $\wedge$ \textsc{patient}($x)(\epsilon$) $\wedge$ \textsc{build}($\epsilon$)\\ \ibar{V} [$\lambda x$ $\lambda y$ $\lambda \epsilon$ \textsc{agent}($y)(\epsilon$) $\wedge$ \textsc{patient}($x)(\epsilon$) $\wedge$ \textsc{build}($\epsilon$)\\ \obar{V} [build]] ]
	]
	\end{forest}
\end{exe}
This violates (\ref{ex3:vptype}), and gets the meaning wrong. By contrast, when the VP headed by \emph{build} supplies an argument of the right kind, the rules of semantic composition can achieve~(\ref{ex3:vptype}).
\begin{exe}
	\raggedright
	\ex \label{ex3:yesbuild} 
	\small
	\begin{forest}
	[$\lambda \epsilon$ \textsc{agent}(\denotes{Cherlon})($\epsilon$) $\wedge$ \textsc{patient}(\denotes{them})($\epsilon$) $\wedge$ \textsc{build}($\epsilon$)\\ VP
		[DP [Cherlon, roof] ]
		[$\lambda y$ $\lambda \epsilon$ \textsc{agent}($y)(\epsilon$) $\wedge$ \textsc{patient}(\denotes{them})($\epsilon$) $\wedge$ \textsc{build}($\epsilon$)\\ \ibar{V} 
			[$\lambda x$ $\lambda y$ $\lambda \epsilon$ \textsc{agent}($y)(\epsilon$) $\wedge$ \textsc{patient}($x)(\epsilon$) $\wedge$ \textsc{build}($\epsilon$)\\\obar{V} [build] ]
			[DP [them, roof] ]
		 ]
	]
	\end{forest}
\end{exe}

Let's now consider how we might achieve the second consequence of the first clause of Theta Criterion: that a $\theta$-role cannot be matched to more than one argument. Heim and Kratzer argue that this also follows from a semantics like the one given here. %
%
\marginnote[0pt]{The semantics I've laid out here follows \citet{Heim:1998}, but is slightly different. Their semantics does not include eventualities. I've made some adjustments to the way they derive the Theta Criterion as a consequence. See \citet{Chung:2004a} for a similar system.}%
%
Consider an imaginary verb, \emph{schkiss}, which assigns the same $\theta$-role to its external and internal arguments.
\begin{exe}
	\raggedright
	\ex \label{ex3:schk} 
	\denotes{schkiss} = $\lambda x$ $\lambda x$ $\lambda \epsilon$ \textsc{patient}($x)(\epsilon$) $\wedge$ \textsc{schkiss}($\epsilon$), where
	\begin{xlist}
	\ex $x$ is a thing and $\epsilon$ an eventuality
	\ex \textsc{patient} c-selects DP
	\end{xlist}
\end{exe}
Suppose that a schkissing is an eventuality just like a kissing, but in which neither participant is an Agent. The function in (\ref{ex3:schk}) is broken because it binds the variable $x$ from two $\lambda$ operators. Consider how our semantic rules will apply this meaning to a VP headed by \emph{schkiss}:
\begin{exe}
	\raggedright
	\ex \label{ex3:schkt} 
	\small
	\begin{forest}
	[VP
		[DP [Pius, roof] ]
		[$\lambda x$ $\lambda \epsilon$ \textsc{patient}(\denotes{Elena})($\epsilon$) $\wedge$ \textsc{schkiss}($\epsilon$)\\ \ibar{V}
			[$\lambda x$ $\lambda x$ $\lambda \epsilon$ \textsc{patient}($x)(\epsilon$) $\wedge$ \textsc{schkiss}($\epsilon$)\\ \obar{V} [schkiss] ]
			[DP [Elena, roof] ]
		 ]
	]
	\end{forest}
\end{exe}
Because there is no rule of semantic composition that can bring the denotation of the \ibar{V} and \emph{Pius} together in (\ref{ex3:schkt}), this will fail to get an interpretation and the result is ungrammatical.

If this thinking is correct, then we can see the first clause of the Theta Criterion as deriving entirely from constraints on what the denotations of verbs can be and how the meanings of the pieces of a phrase compose. What's left of the Theta Criterion, then, is~(\ref{ex3:slimtc}).
\begin{exe}
	\raggedright
	\ex \label{ex3:slimtc} Theta Criterion\\
	Every argument XP must be related to exactly one position assigned a $\theta$-role.
\end{exe}
This follows from our thesis that the only phrase marker in a derivation that is semantically interpreted is D-structure. The \ibar{X} Skeleton requires every XP to occupy just one position in a phrase-marker. Our thesis that D-structure is the only phrase-marker that is semantically interpreted requires that the semantic rule of composition that puts an argument together with its $\theta$-role look at just one phrase-marker. Together, these derive~(\ref{ex3:slimtc}). These are steps, then, towards understanding how the Theta Criterion might be derived from a semantics.

Note that this way of deriving (\ref{ex3:slimtc}) is incompatible with the movement theory of control. That theory requires a syntax that allows more than just a D-structure to be semantically interpreted. There is some evidence that such a syntax is indeed required for other reasons, and we will eventually develop such a syntax. But our present syntax does not allow for movement from one $\theta$-marked position to another for the simple reason that one or the other of the phrase-markers in such a derivation will not be semantically interpretable. 

\section{Theta-roles and Case}
\label{sec:thetacase}

One can think of Argument Movement as a function from a parse that satisfies the Projection Principle and Theta Criterion to a parse that satisfies the Case Filter and~EPP. $\theta$-roles are determined by lexical items and the phrases they project, and where Case is assigned is also controlled by the these lexical items. Lexical items, and verbs in particular, play a big role in setting up the environments in which Argument movement is called upon. The conditions on Argument movement place a cap, then, on the kinds of situations that verbs can create. They won't allow verbs to exist which create conflicts that Argument movement can't fix.

Verbs set up these situations by influencing where the $\theta$-marked and Case marked positions for arguments are. Thus, for example, we have seen that there are verbs which assign one external $\theta$-role, and verbs which assign an external and an internal $\theta$-role and Accusative Case. These are sometimes called intransitive and transitive verbs, respectively; examples are in (\ref{ex3:simple}).\enlargethispage*{13pt}
\begin{exe}
	\ex \label{ex3:simple} 
	\begin{xlist}
	\ex[]{Martin slept. \hfill (intransitive) \hspace*{1in}}
	\ex[]{Martin likes kiwis. \hfill (transitive) \hspace*{1in}}
	\end{xlist}
\end{exe}
A kind of verb related to transitives are exemplified by \emph{believe}, \emph{consider} and \emph{make}, which assign their Accusative Case to something different than the argument they c-select. As we've seen, these verbs can take a ``clausal'' complement---sometimes these are ``small clauses''---and assign their Accusative Case to a DP within this clause.
\begin{exe}
	\ex \label{ex3:ecm}
	\begin{xlist}
	\ex She believes [\IP him to be unhappy ].
	\ex She considers [\AP him happy ].
	\ex She made [\VP him dance ].
	\end{xlist} 
\end{exe}

We might ask whether there is a connection between a verb having accusative Case and c-selecting a DP internal argument. %
%
\marginnote[0pt]{I will say that a word c-selects an argument if one of its $\theta$-roles c-selects an argument.}%
%
Presently our theory makes these two independent attributes that a verb might have. However, empirically they are confounded. If a verb has accusative Case, then it is capable of c-selecting a DP object that we see with accusative Case. And if a verb does not have accusative Case, then it will not be able to c-select a DP object that shows up with accusative Case. One way we can tell that a verb has accusative Case, then, is to discover that it c-selects a DP object that appears with accusative Case. And if a verb does not c-select a DP object, it will be difficult to tell whether that verb can assign accusative Case since it will not be able to occur with an accusative Case-marked object. So we should ask: are assigning accusative Case and c-selecting a DP object independent properties of a verb?

The answer seems to be yes. We have already seen that passive creates verbs that c-select DP objects but do not assign accusative Case. 
\begin{exe}
	\ex \label{ex3:passecm}
	\begin{xlist}
	\ex[]{Martin was believed [\AP unhappy ].}
	\ex[]{Martin was believed.}
	\end{xlist}
\end{exe}
From (\ref{ex3:passecm}), we see that the passive \emph{believe} combines with the argument it c-selects---here a small clause~AP. Because it does not assign accusative Case, however, the argument DP \emph{Martin} inside that small clause AP cannot get Accusative in its D-structure position, and must move to receive Nominative case. In (\ref{ex3:passecm}b), then, we should conclude that the same set of attributes adhere to the passive \emph{believed}. It continues to c-select its DP argument---here \emph{Martin}---but no longer assigns Accusative Case. Thus, just as in (\ref{ex3:passecm}a), the DP that would have gotten accusative from the active \emph{believe}, namely \emph{Martin}, must move to a position where it gets another case. In (\ref{ex3:passecm}b), that position is Specifier of IP, and the case it receives is Nominative. From these examples, then, we can learn that not all verbs that c-select DPs need be verbs that assign Accusative Case. 

We can reach this same conclusion by considering the internal arguments of English adjectives. Adjectives are a class of word that never assigns accusative Case, and so it is no surprise that no adjective c-selects a DP that shows up in the accusative Case. We find contrasts like those in~(\ref{ex3:acase}).
\pagebreak
\begin{exe}
	\raggedright
	\ex \label{ex3:acase} 
	\begin{xlist}
		\ex[*]{Elena is happy the fact.}
		\ex[]{Elena is happy that pigs don't fly.}
		\ex[]{Elena is happy about the fact that pigs don't fly.}
	\end{xlist}
\end{exe}
Adjectives can c-select PPs and CPs that stand in complement position, partly because these phrases do not need to be assigned Case. However, there is some evidence that suggests that adjectives can c-select a DP if it is possible for that DP to get Case from some term other than the adjective. To see this, reconsider the alternation in~(\ref{ex3:acasecp}).
\begin{exe}
	\raggedright
	\ex \label{ex3:acasecp} 
	\begin{xlist}
		\ex It is obvious [\CP that pigs don't fly].
		\ex{} [\CP That pigs don't fly] is obvious.
	\end{xlist}
\end{exe}
These sentences have a D-structure that contains the \ibar{I} in~(\ref{ex3:acasecpd}).
\begin{exe}
	\raggedright
	\ex \label{ex3:acasecpd} 
	\small
	\begin{forest}
	[\ibar{I}
		[\obar{I} [pres] ]
		[VP [\ibar{V}
			[\obar{V} [be] ]
			[AP [\ibar{A} 
				[\obar{A} [obvious] ]
				[CP [that pigs don't fly, roof] ]
			] ]
		] ]
	]
	\end{forest}
\end{exe}
In (\ref{ex3:acasecp}a), the IP that this \ibar{I} is embedded in contains an expletive in its Specifier position, and this is how (\ref{ex3:acasecp}a) satisfies the EPP. In (\ref{ex3:acasecp}b), by contrast, the IP that embeds this \ibar{I} does exist (or is empty), and the EPP is satisfied by A Movement creating an S-structure in which the CP \emph{that pigs don't fly} occupies this position.

Consider now the contrast in~(\ref{ex3:acase2}).
\begin{exe}
	\raggedright
	\ex \label{ex3:acase2} 
	\begin{xlist}
		\ex[*]{It is obvious this fact.}
		\ex[]{This fact is obvious.}
	\end{xlist}
\end{exe}
This alternation is understood if both of the examples in (\ref{ex3:acase2}) have the \ibar{I} in~(\ref{ex3:acase2d}) at D-structure.
\begin{exe}
	\raggedright
	\ex \label{ex3:acase2d} 
	\small
	\begin{forest}
	[\ibar{I}
		[\obar{I} [pres] ]
		[VP [\ibar{V} 
			[\obar{V} [be] ]
			[AP [\ibar{A} 
				[\obar{A} [obvious] ]
				[DP [this fact, roof] ]
			] ]
		] ]
	]
	\end{forest}
\end{exe}
If the D-structure has an IP that embeds this \ibar{I} whose Specifier is filled with \emph{it}, then there is a Case filter violation and no way of forming an S-structure that remedies that violation. But if the D-structure containing this \ibar{I} embeds it in an IP without an \emph{it} in Specifier position, then A Movement can form an s-structure that satisfies the Case filter; this s-structure corresponds to the grammatical~(\ref{ex3:acase2}b). If this thinking is correct, then the class of words that c-select DPs cannot be just the same class that assigns accusative Case. Adjectives seem to be able to c-select DPs, and they never assign Accusative case.

Let's consider one last scenario that leads to this same conclusion. Consider the alternations in~(\ref{ex3:verbndp2}).
\begin{exe}
	\raggedright
	\ex \label{ex3:verbndp2} 
	\begin{xlist}
		\ex[*]{She hopes the outcome.}
		\ex[]{She hopes that this will be the outcome.}
		\ex[*]{She insisted the claim.}
		\ex[]{She insisted that this is true.}
		\ex[*]{She felt the hypothesis.}
		\ex[]{She felt that this is likely.}
	\end{xlist}	
\end{exe}
These are verbs that, like the adjective \emph{obvious}, can have a c-selected CP in complement position, but not a c-selected~DP. There are two ways of describing these verbs that will account for their inability to show up with an argument DP in complement position. It could be that these verbs don't assign accusative Case, but do c-select DPs. Or it could be that these verbs don't c-select a DP internal argument. One way we can tell is to look at what happens when these verbs are passivized. If, like \emph{obvious}, they c-select DPs and all that's gone wrong is that those DPs are not able to get Case from these verbs, then the passive versions of these verbs should be fine. In the passive, the Specifier of IP will no longer be occupied by the subject argument, and A movement should be able to put the putative DP object into this Case-marked position. As it happens, however, the passive versions of these verbs still do not allow for DP arguments.
\pagebreak
\begin{exe}
	\raggedright
	\ex \label{ex3:verbndp3} 
	\begin{xlist}
		\ex[*]{The outcome was hoped.}
		\ex[*]{The claim was insisted.}
		\ex[*]{The hypothesis was felt.}
	\end{xlist}
\end{exe}
This follows if these verbs do not c-select a DP complement. The passive is possible of these verbs when their complements are CPs, as in~(\ref{ex3:verbndp4}).
\begin{exe}
	\raggedright
	\ex \label{ex3:verbndp4} 
	\begin{xlist}
		\ex It was hoped that this will be the outcome.
		\ex It was insisted that this is true.
		\ex It was felt that this is likely.
	\end{xlist}
\end{exe}
So the problem with (\ref{ex3:verbndp3}) isn't something about whether the passive can apply to these verbs. A reasonable conclusion, then, is that these verbs do not c-select~DPs. In general, whether a word -- verb or adjective -- assigns accusative Case or c-selects a DP argument appears to be independent.

One interesting mystery that arises with the verbs \emph{hope}, \emph{insist}, and \emph{felt} is that A movement cannot apply to complement CPs in the way we would expect. Examples like (\ref{ex3:verbndp5}), which should be parallel to (\ref{ex3:verbndp4}), are ungrammatical.
\enlargethispage*{\baselineskip}
\begin{exe}
	\raggedright
	\ex \label{ex3:verbndp5} 
	\begin{xlist}
		\ex[*]{That this will be the outcome was hoped.}
		\ex[*]{That this is true was insisted.}
		\ex[*]{That this is likely was felt.}
	\end{xlist}
\end{exe}
What distinguishes (\ref{ex3:verbndp5}) from (\ref{ex3:verbndp4}) is merely that the EPP is satisfied in (\ref{ex3:verbndp5}) by using A Movement to move the CP into Specifier of~IP. The ungrammaticality of these examples is part of a larger pattern concerning A movement in English, which can be described with~(\ref{ex3:dpposA}).
\begin{exe}
	\raggedright
	\ex \label{ex3:dpposA} 
	A movement can only move an XP that is in a position that a DP could stand in.
\end{exe}
Because the complement position of these verbs is not a position that a DP is c-selected for, A movement from this position is prevented by (\ref{ex3:dpposA}). By contrast, because the adjective \emph{obvious} does c-select a DP, A movement from the complement position of \emph{obvious} is possible, and examples like (\ref{ex3:acasecp}b) are grammatical. Why (\ref{ex3:dpposA}) is true strikes many syntacticians as mysterious. %
%
\marginnote[-65pt]{That CPs can only move from DP positions is argued in \citet{Alrenga:2005}, who builds on \citet{Koster:1978}. (The observations about where CPs can and can't move from originate, largely, from \citealt{Emonds:1976}.) For recent discussions of the situation involving CPs, see \citet{Moulton:2009} and \citet{Elliott:2017}. And for a discussion of the general issues concerning (\ref{ex3:dpposA}), see \citet{Poole:2017}.}%
%
An important point about accusative Case, then, is that its assignment depends entirely on the syntactic proximity that some position has to a word that licenses accusative. If a position is adjacent to and governed by a word with accusative, then that position is assigned accusative Case. In this respect, accusative is like the other Cases of English. Nominative and genitive are similarly Cases that are assigned to positions that are in a certain syntactic configuration relative to a term that licenses them. Their occurrence is determined entirely by the syntactic structures they are in.

Not all Cases of the world's languages have this property. For instance, the dative Case in German seems to be tied more directly to the meaning of a word. The dative in German shows up on DPs that are c-selected by certain verbs. For instance, the internal argument of the verb \emph{helfen} (`help') is in the dative Case. Compare the situation with \emph{sehen} (`see'), whose internal argument is in the accusative.
\begin{exe}
	\raggedright
	\ex \label{ex3:german1} 
	\begin{xlist}
		\ex \gll Sie hilft mir. \\ she.\textsc{nom} help.\textsc{pres} me.\textsc{dat} \\ `She helps me.'
		\ex \gll Sie sieht mich. \\ she.\textsc{nom} sees.\textsc{pres} me.\textsc{acc} \\ `She sees me.'
	\end{xlist}
\end{exe}
Unlike the accusative Case, the dative Case in German does seem to correlate with c-selection. If a DP bears dative Case in German, then it will get it from the very verb that it is c-selected by. Cases that are tied to the c-selection properties of the word that assigns that Case are sometimes called ``inherent Case.'' And the Cases that English has, which are assigned in a way that is independent of the c-selection properties of the Case assigner, are called ``structural Cases.'' In German, as in many languages, this distinction in Case can be revealed by seeing whether A movement can affect the Case borne by the relevant DP. For instance, the verb that passive creates from \emph{helfen} does not cause the DP it c-selects to ``change'' its Case. This \underline{is} what happens with the passivized version of \emph{sehen}, however. Just as in English, passive makes from \emph{sehen} a verb that no longer supports Accusative case, and consequently its c-selected DP shows up in the nominative Case. This is seen in~(\ref{ex3:german2}).
\begin{exe}
	\raggedright
	\ex \label{ex3:german2} 
	\gll Ich wurde gesehen. \\ I.\textsc{nom} become.\textsc{pres} see.\textsc{pass} \\ `I was seen.'
	\marginnote[-30pt]{The verb that embeds a passive participle in German is not the same as the verb used to embed adjectives in German, i.e. \emph{sein} (`be'). Instead the verb that embeds passive participles is \emph{werden}, which is often glossed with the English \emph{become}.}
\end{exe}
German syntax differs from English in a variety of ways, and the movement of \emph{Ich} in (\ref{ex3:german2}) is achieved by a different set of rules than in English.\marginnote[10pt]{We will learn about this in the next chapter.} What's key now is that the object of \emph{sehen} shows up in the nominative Case, and that is required. Just as in English, it cannot be accusative. By contrast, the passive version of \emph{helfen} leaves the DP c-selected by \emph{helfen} in the dative Case; it cannot show up in the nominative.
\begin{exe}
	\raggedright
	\ex \label{ex3:german3} 
	\begin{xlist}
		\ex[]{\gll Mir wurde geholfen. \\ I.\textsc{dat} become.\textsc{pres} help.\textsc{pass} \\ `I was helped' }
		\ex[*]{\gll Ich wurde geholfen. \\ I.\textsc{nom} become.\textsc{pres} help.\textsc{pass} \\ `I was helped.' }
	\end{xlist}
\end{exe}
\pagebreak

In view of these differences, we can render our rule for the passive with~(\ref{ex3:newpassr1}).
\begin{exe}
	\raggedright
	\ex \label{ex3:newpassr1} Passive
	\begin{xlist}
		\ex Add passive morphology to the active verb and embed it under \emph{be}, and
		\ex Rob the verb of its structural Case, and
		\ex Rob the verb of its external $\theta$-role.
	\end{xlist}
\end{exe}
The earlier version of Passive had (\ref{ex3:newpassr1}b) rob the verb of its Accusative Case. If Accusative is always structural, and it's the only structural Case, then this new version will do what the old version did.

In IndoEuropean languages, there is a strong correlation between whether a Case is structural or inherent on the one hand and whether the Case is nominative or accusative on the other. Quite generally, accusative and nominative are structural Cases and the others---dative, genitive, instrumental, etc.---are inherent Cases. But there are corners of IndoEuropean where this correlation appears to break down. For instance, in Faroese it is possible to find passive versions of dative Case assigning verbs whose c-selected DPs show up in the nominative. Einar Sigur\dh sson provides the example in~(\ref{ex3:faro}).
\begin{exe}
	\raggedright
	\ex \label{ex3:faro} 
	\begin{xlist}
		\ex \gll Teir r\'ostu henni alti\dh. \\ they.\textsc{nom} praise.\textsc{past} her.\textsc{dat} always \\ `They always praised her.'
		\ex \gll Hon {var\dh} {alti\dh} r\'ost. \\ she.\textsc{nom} be.\textsc{past} always praise.\textsc{pass}. \\ `She was always praised.'
	\end{xlist}
\end{exe}
One way of accounting for these cases is to allow Passive to refer to structural Case, as it does in the new passive rule in (\ref{ex3:newpassr1}), and let which Cases are structural or inherent vary.\marginnote[-30pt]{The so-called ``insular'' North Germanic languages (Icelandic and Faroese) are notorious for having a complicated relationship between the morphological forms of Case and their status as structural or inherent. See \citet{Zaenen:1985}, and references therein.}

\subsection{Unaccusatives}

We have seen two types of cases, then, in which a term c-selects a DP but does not assign accusative Case. There are certain adjectives of this kind (e.g., \emph{obvious}) and there are verbs created by the passive operation from transitive verbs. There is evidence that some verbs which might otherwise look like intransitive verbs fall into this last class too. For instance some verbs have two versions, if you like, one in which they are transitive and the other in which they are intransitive. Some examples are~(\ref{ex3:unacc1}).
\begin{exe}
	\ex \label{ex3:unacc1} 
	\begin{xlist}
		\ex She bounced them.
		\ex They bounced.
		\ex She melted them.
		\ex They melted.
		\ex She opened them.
		\ex They opened.
		\ex She grew them.
		\ex They grew
	\end{xlist}
\end{exe}
In their transitive guise, they c-select a DP and assign a certain $\theta$-role to that object, and they also assign Accusative Case to it. In their intransitive guise, that very same c-selected argument DP appears in the Nominative Case and consequently appears in Specifier of IP. Pairs of sentences like (\ref{ex3:unacc1}) are sometimes said to involve a ``causative/inchoative'' alternation.\marginnote{Perhaps this is done first in \citet{Hall:1965}.} The name ``causative'' expresses a hypothesis about the underly semantics of the transitive guise: \denotes{she melted them} $\approx$ ''she caused them to become melted.'' And the name ``inchoative'' references a hypothesis about the intransitive guise: \denotes{they melted} $\approx$ ``they became melted.'' I'll use this term to refer to this alternation, even though the underlying semantics of the transitive and intransitive guises is probably not equally well suited to it.\marginnote{For instance, \textit{she bounced} them doesn't correspond to ``she caused them to become bounced'' and \textit{they bounced} doesn't correspond to ``they became bounced.''}

It certainly seems as if the $\theta$-role borne by \emph{they/them} is the same in the transitive and intransitive versions, and a simple way of understanding this is to imagine that the verb has assigned that $\theta$-role to the complement position in both of its guises. The difference in the intransitive and transitive versions is whether or not the verb has also assigned Accusative Case and also comes with an external $\theta$-role. That is, the intransitive versions might involve a derivation like that indicated in~(\ref{ex3:derunacc1}).
\begin{exe}
\raggedright
	\ex \label{ex3:derunacc1}
	\small
	\begin{forest}
	[IP
		[DP [they, name=binder, roof]]
		[\ibar{I}
			[\obar{I} [past]]
			[VP [\ibar{V}
				[\obar{V} [bounce]]
				[{\ }, name=trace ]
			 ] ]
		 ]
	]
	\draw[->,>=stealth', thin, gray] (trace) to[out=south west,in=south,looseness=1] node[pos=0.75,xshift=-10mm]{} (binder);
	\end{forest}
\end{exe}
If this is correct, then, it appears that there might be two kinds of single argument verbs---ones that assign their single $\theta$-role externally and another that assign their sole $\theta$-role internally.

David Perlmutter\cite{Perlmutter:1978} was the first to conjecture that some intransitive verbs could assign just an internal $\theta$-role, and no external $\theta$-role. In \citet{Perlmutter:1984}, the name ``unaccusative'' is given to verbs of this kind. The intransitives whose sole argument is external are known as ``unergative.'' I will use the term ``intransitive'' to refer to verbs whose sole argument gets either an external $\theta$-role or an internal $\theta$-role.
\pagebreak
\begin{exe}
\raggedright
	\ex \label{ex3:intrans} Intransitives\marginnote{\citet{Burzio:1986}, who gives an extended defense of the view that Italian has both unaccusative and unergative intransitives, uses an alternative terminology which is also common. He calls those verbs that have just an external $\theta$-role ``intransitive,'' and those verbs whose sole $\theta$-role is internal ``ergative.'' Some of the literature follows Burzio's terminology.}
	\begin{xlist}
		\ex A ghost should sleep. \hfill (unergative)
		\ex A ghost should bounce. \hfill (unaccusative)
	\end{xlist}
\end{exe}
You will find different ways of referring to these single argument verbs, but this terminology I believe is the most common. In this book, I'll use Burzio's term (``unergative'') for the intransitives whose argument is external, and Perlmutter's term (``unaccusative'') for the intransitive whose argument is internal.\marginnote{Another early defense of the syntax of intransitives described here is in \citet{Rosen:1984}. \citet{Levin:1995} is recommended as an introduction.}

There are a variety of phenomena that suggest this is the right way of thinking about the two guises of the verbs in (\ref{ex3:unacc1}). These phenomena also suggest that there are intransitive verbs that don't have a transitive guise which are nonetheless unaccusatives. I'll describe some of these processes here. 

One is the morphological rule that forms adjectival participals from verbs discussed in section 3.3.1. This rule forms an adjective whose relation to the noun it modifies is the same as that given by the internal $\theta$-role of the corresponding verb. (\ref{ex3:deadj}) illustrates.
\begin{exe}
\raggedright
	\ex \label{ex3:deadj}
	\begin{xlist}
		\ex a smoothed blanket
		\begin{xlist}
			\ex = a blanket which something smoothed
			\ex $\neq$ a blanket which smoothed something\\
			(compare: \textit{She smoothed the blanket.}
		\end{xlist}
		\ex some mixed nuts
		\begin{xlist}
			\ex = some nuts which someone mixed
			\ex $\neq$ some nuts which mixed something\\
			(compare: \textit{She mixed the nuts}.)
		\end{xlist}
		\ex the painted man
		\begin{xlist}
			\ex = the man which someone painted
			\ex $\neq$ the man which painted something\\
			(compare: \textit{She painted the man.} )
		\end{xlist}
		\ex the known astronaut
		\begin{xlist}
			\ex = the astronaut someone knows
			\ex $\neq$ the astronaut who knows someone\\
			(compare: \textit{She knows an astronaut.})
		\end{xlist}
	\end{xlist}
\end{exe}
As expected, adjectives formed in this way from intransitives are largely not possible.
\begin{exe}
\raggedright
	\ex \label{ex3:nodeadj}
	\begin{xlist}
		\ex[*]{the sung alto}
		\ex[*]{the barked dog}
		\ex[*]{the slept woman}
		\ex[*]{the cried child}
	\end{xlist}
\end{exe}
If the intransitive verbs these adjectives are formed from have just an external $\theta$-role, this is expected. The verbs that participate in the causative/inchoative alternation in (\ref{ex3:unacc1}) are among the verbs that can become adjectives with this morphology.
\begin{exe}
\raggedright
	\ex \label{ex3:inchounacc}
	\begin{xlist}
		\ex a bounced ball
		\ex a melted car
		\ex an opened door
		\ex a grown child
	\end{xlist}
\end{exe}
But this shouldn't be surprising, as it's conceivable that the adjectives in (\ref{ex3:inchounacc}) are formed from the transitive versions of the verbs. 

There's a small class of intransitive verbs which are able to become adjectival participles, even though they have no transitive version. Some cases are in~(\ref{ex3:unaccadj}).\marginnote{See \citet{Hoekstra:1984} and \citet{Levin:1986} for this generalization.}
\begin{exe}
\raggedright
	\ex \label{ex3:unaccadj}
	\begin{xlist}
		\ex the dead roach
		\ex the recently arrived package
		\ex the risen divers
		\ex the surfaced submarine
	\end{xlist}
\end{exe}
The verbs these adjectives derive from (\textit{die}, \textit{arrive}, \textit{rise}, and \textit{surface}) could be unaccusatives. If the verbs that adjectival participles are formed from must have an internal $\theta$-role, then we should conclude that the verbs from which the adjectives in (\ref{ex3:unaccadj}) derive from are unaccusative.

Another morphological process which could be used to distinguish the two intransitives is the suffix \textit{er} which forms nouns from verbs. Some examples are in~(\ref{ex3:ern}).
\begin{exe}
\raggedright
	\ex \label{ex3:ern}
	\begin{xlist}
		\ex (the) singer
		\ex (the) maker
		\ex (the) baker
		\ex (the) manufacturer
		\ex (the) stapler
	\end{xlist}
\end{exe}
Suffixation of \textit{er} onto a verb creates a noun that refers to things that are whatever the external $\theta$-role of the verb is. Thus, a \textit{singer} is something that sings, and a \textit{stapler} is something that staples.\marginnote{Chapter 5.3 sketches an anaylsis.} What we might expect, then, is that \textit{er} can only attach to verbs that have an external $\theta$-role. That means we should find nouns formed from transitive verbs and unergative verbs, but not from unaccusative and raising verbs. Note first that \textit{er} cannot form a noun from the intransitive guise of the verbs that participate in the causative/inchoative alternation in (\ref{ex3:unacc1}). This is what (\ref{ex3:erex}) shows. 
\pagebreak
\begin{exe}
\raggedright
	\ex \label{ex3:erex}
	\begin{xlist}
		\ex the bouncer
		\begin{xlist}
			\sn = the one who bounced something
			\sn $\neq$ the one who is bounced\\
			(\textit{she bounced it} $\sim$ \textit{it bounced})
		\end{xlist}
		\ex the melter
		\begin{xlist}
			\sn = the one who melted something 
			\sn $\neq$ the one who is melted\\
			(\textit{she melted it} $\sim$ \textit{it melted})
		\end{xlist}
		\ex the grower
		\begin{xlist}
			\sn = the one who grows something
			\ex $\neq$ the one who is grown\\
			(\textit{she grew it} $\sim$ \textit{it grew})
		\end{xlist}
	\end{xlist}
\end{exe}
This supports the view that the intransitive versions of these verbs are unaccusative. If the intransitive version of these verbs doesn't have an external $\theta$-role, then they should not be able to form nouns with \textit{er}. Instead, the nouns can only be formed from the transitive versions of these verbs. (\ref{ex3:erex}) is just what we should expect on that view. 

Note that \textit{er} can suffix onto verbs that are intransitive.
\begin{exe}
\raggedright
	\ex \label{ex3:erunerg}
	\begin{xlist}
		\ex (the) runner
		\ex (the) dancer
		\ex (the) sleeper
	\end{xlist}
\end{exe}
The reason \textit{er} cannot attach to the intransitive version of verbs which participate in the causative/inchoative alternation is not because it cannot attach to intransitives, then. Presumably, the difference in the intransitives that can and can't attach to \textit{er} is whether the sole argument of that intransitive gets an external $\theta$-role. 

As (\ref{ex3:raiseer}) indicates, there are other intransitive verbs that cannot be made into nouns with \textit{er}.
\begin{exe}
\raggedright
	\ex \label{ex3:raiseer}
	\begin{xlist}
		\ex * (the) seemer
		\ex * (the) appearer
		\ex * (the) arriver
		\ex * (the) dier
		\ex * (the) surfacer
		\ex * (the) riser
	\end{xlist}
\end{exe}
These include the raising verbs \textit{seem} and \textit{appear}, and this is expected as we've already seen evidence that these verbs do not have an external $\theta$-role. But they also include the intransitive verbs we saw could form adjectival participles in (\ref{ex3:unaccadj}). To the extent that there is a correlation between the intransitive verbs that cannot form nouns with \textit{er} but can form adjectival participles, there is evidence for the existence of unaccusative verbs because their existence predicts this correlation. 

Consider next the prefix \textit{re}, which forms verbs from other verbs. Some examples are in~(\ref{ex3:reunacc}).\marginnote{Using \textit{re} prefixation to diagnose unaccusatives comes from \citet{Horn:1980}.}
\begin{exe}
\raggedright
	\ex \label{ex3:reunacc}
	\begin{xlist}
		\ex I resanded the table.
		\ex She repainted the wall.
		\ex He rethought the problem.
	\end{xlist}
\end{exe}
Roughly speaking, what \textit{re} does is create a verb whose meaning is identical to the one it attaches to, but adds that whatever the V containing the verb and object describes is something that happened previously.\marginnote{See Chapter 5.2.1 for an analysis of the similar meaning that again has.} Thus, for instance, (\ref{ex3:reunacc}a) communicates that \textit{said the table} is something that was true previously. I cannot say (\ref{ex3:reunacc}a) if I am sanding a table that has never been sanded before. If, however, it was previously sanded, even if that sanding was by someone other than me, I can say (\ref{ex3:reunacc}a). Note in particular that (\ref{ex3:reunacc}a) does not communicate that I sanded something other than the table previously. It's only if the previous sanding were of the table that (\ref{ex3:reunacc}a) can be used. So very rough paraphrases of the meanings that are delivered by in (\ref{ex3:reunacc}) are given in~(\ref{ex3:reunaccp}).
\begin{exe}
\raggedright
	\ex \label{ex3:reunaccp}
	\begin{xlist}
		\ex \denotes{I resanded the table}\\ $\approx$ I sanded the table and there was a prior resanding of it
		\ex \denotes{She repainted the wall}\\ $\approx$ She painted the wall and there was a prior painting of it
		\ex \denotes{He rethought the problem}\\ $\approx$ He thought about the problem and there was a prior thinking about it.
	\end{xlist}
\end{exe}
The meaning of \textit{re}, then, targets the meaning of the phrase that is formed by putting the internal argument together with the verb. If there is no internal argument, using \textit{re} is odd. Consider~(\ref{ex3:resang}).
\begin{exe}
\raggedright
	\ex \label{ex3:resang}
	\begin{xlist}
		\ex I resang.
		\ex I rebaked.
		\ex I reslept.
	\end{xlist}
\end{exe}
This matches the oddness of the paraphrases that we might expect their meanings to correspond to:
\begin{exe}
\raggedright
	\ex \label{ex3:resang2}
	\begin{xlist}
		\ex I sang and there was a prior singing
		\ex I baked and there was a prior baking
		\ex I slept and there was a prior sleeping
	\end{xlist}
\end{exe}
The oddness of (\ref{ex3:resang2}) might have to do with how trivial the meaning that \textit{re} provides is. It might be taken for granted that I am not the first one to sing, bake or sleep, but this is all that \textit{re} is adding to the sentences in (\ref{ex3:resang}).

Now, interestingly, there are some intransitive verbs that \textit{re} can attach to without complaint, and the resulting meaning treats the argument of that intransitive as an internal argument. The intransitives that participate in the causative/inchoative alternation are among these.
\begin{exe}
\raggedright
	\ex \label{ex3:inchore}
	\begin{xlist}
		\ex The ball rebounced.\\ $\approx$ the ball bounced and there was a prior bouncing of the ball
		\ex The ice remelted.\\ $\approx$ the ice melted and there was a prior melting of the ice
		\ex The flowers regrew.\\ $\approx$ the flower grew and there was a prior growing of them\marginnote{Note that the referent of the flowers in (\ref{ex3:inchore}c) cannot be the same in the two conjucts of its meaning. We need to let a definite description like \textit{the flowers} be able to refer to different tokens of the same thing.}
		\ex The door reopened.\\ $\approx$ the door opened and there was a prior opening of it
	\end{xlist}
\end{exe}
This is another piece of evidence, then, that the intransitive use of these verbs involves an internal $\theta$-role, not an external $\theta$-role. If has a meaning that references the phrase which combines verb and direct object, then we should treat the arguments in (\ref{ex3:inchore}) as direct objects.

The prefix \textit{re} can also attach to the intransitives that \textit{er} suffixation and the rule that forms adjectival participles suggest are unaccusative. This is what (\ref{ex3:unacccore}) shows.
\begin{exe}
\raggedright
	\ex \label{ex3:unacccore}
	\begin{xlist}
		\ex She reappeared.\\ $\approx$ She appeared and there was a prior appearance of her
		\ex It rearrived.\\ $\approx$ It arrived and there was a prior arrival of it
		\ex It redied.\\ $\approx$ It died and there was a prior death of it\marginnote{(\ref{ex3:unacccore}c) describes situations that are more frequent in fiction than they are in reality.}
		\ex It resurfaced.\\ $\approx$ It surfaced and there was a prior surfacing of it
		\ex It rerose.\\ $\approx$ It rose and there was a prior rising of it
	\end{xlist}
\end{exe}
Again, to the extent that there is a correlation between the class of intransitives that \textit{re} can attach to, producing a meaning that includes the argument of that intransitive, and the class of verbs that cannot form nouns with \textit{er} but can form adjectival participles, we have support for the presence of unaccusatives. This correlation is just what would be expected if this class of verbs were unaccusatives. 

I'll describe one last phenomenon in English which is relevant to determining whether an intranstivie is an unaccusative or an unergative. This phenomenon involves a syntactic construction called a ``resultative,'' of which (\ref{ex3:result}) is an example.
\begin{exe}
\raggedright
	\ex \label{ex3:result} The rod pounded the lump flat.
\end{exe}
The adjective \textit{flat} in this example is understood to describe the result of the pounding. And the property it describes must be understood as holding of the object argument, not the subject argument.%
%
\marginnote[-28pt]{This generalization about resultatives is introduced in \citet{Simpson:1983c}. See \citet[chapter 2]{Levin:1995} for a discussion of its utility in diagnosing unaccusatives. Some possible counter-examples are in (i), from \citet{Wechsler:1997}.
\begin{exe}
	\exi{i.}
	\begin{xlist}
		\ex The wise men followed the star out of Bethlehem.
		\ex The sailors caught a breeze and rode it clear of the rocks.
		\ex He followed Lassie free of his captors
	\end{xlist}
\end{exe}
See \citet[chapter 3.7]{Rothstein:2004} and \citet[chapter 5]{Levin:1995} for some discussion.}%
%
\ From (\ref{ex3:result}) we can conclude that the pounding changed the lump into a flat lump, but not that the pounding changed the rod into a flat rod. This sensitivity to being an external or internal argument in the resultative seems to be quite general. In all of the examples of (\ref{ex3:resultobj}), the adjective describes a state that the object acquires, not a state that the subject acquires.\marginnote[10pt]{Resultatives should be distinguished from constructions which look similar but have a different meaning. In the sentence \textit{Kaden met Shay drunk}, the adjective \textit{drunk} can describe a property that either Kaden or Shay have. But this isn't a resultative because that property is not understood to be the result of the meeting. These constructions are what \citet{Halliday:1967} calls ``depictives.''}
\begin{exe}
\raggedright
	\ex \label{ex3:resultobj}
	\begin{xlist}
		\ex She stabbed me bloody.
		\ex She spun me dizzy.
		\ex The grill cooked the meat black.
		\ex The machine shook the bottle empty.
	\end{xlist}
\end{exe}
As expected, then, resultatives aren't felicitous unless there is an internal argument for it to describe. This is why resultatives are generally not possible with intransitive predicates.\marginnote[20pt]{These are grammatical only as depictives.}
\begin{exe}
\raggedright
	\ex \label{ex3:nointres}
	\begin{xlist}
		\ex[*]{I walked exhausted}
		\ex[*]{I sang hoarse.}
		\ex[*]{It flew unscared.}
		\ex[*]{The child cried tired.}
		\ex[*]{I baked sweaty.}
		\ex[*]{I slept awake.}
	\end{xlist}
\end{exe}
But—--perhaps unsurprisingly by now--—resultatives are possible with the intransitive version of the sentences that participate in the causative/inchoative alternation.
\begin{exe}
\raggedright
	\ex \label{ex3:resincho}
	\begin{xlist}
		\ex It bounced flat.
		\ex It melted away.
		\ex It opened wide.
		\ex It grew tall.
	\end{xlist}
\end{exe}
Like the behavior of the prefix \textit{re}, the suffix \textit{er}, and adjectival participles, this supports the conclusion that these intransitives are unaccusatives.

Unforuntatly resultatives do not combine with all of the other verbs we've seen could be unaccusative. The difficulty forming resultatives of these verbs is indicated by~(\ref{ex3:nogen}).
\begin{exe}
\raggedright
	\ex \label{ex3:nogen}
	\begin{xlist}
		\ex The roach died quiet.
		\ex The man arrived warm.
		\ex The submarine surfaced lower than the plane.
		\ex The bird rose dizzy.
	\end{xlist}
\end{exe}
To the extent that any of these sentences are grammatical, it's because the adjective is not clearly understood to be a result of the event described by the verb. Thus, for instance, (\ref{ex3:nogen}a) is a grammatical sentence if it describes a scenario in which the roach died while quiet. It cannot describe a scenario in which the roach's death resulted in it being quiet. The adjectives in (\ref{ex3:nogen}) are depictives, not resultatives. 

The events described by \textit{die} and \textit{arrive} are punctual—--that is they describe events that have vanishingly short durations—--and this might make them ineligible to have results. Intuitively, resultatives describe a state that unfolds as part of the process that the verb describes. Perhaps a state cannot be part of something as punctual as an arrival or a death. That reasoning would not extend to \textit{rise} and \textit{surface}, though. And this brings us to a problem lurking in these various arguments for the existence of unaccusatives. All of the processes I've presented as relevant for determining when an intransitive is unergative or unaccusative aren't productive. That is, there are other criteria that determine whether these processes can apply. This is shown below, where each process and the enironment it (putatively) requires are paired with examples of that process in that environment that are nonetheless ungrammatical.
\begin{exe}
\raggedright
	\ex \label{ex3:nocorrel}
	\begin{xlist}
		\ex 
		\begin{tabular}{ll}
			generalization: & adjectival participles are formed from verbs that have a direct object.\\ 
			counter-examples: & *\textit{the hoped fact}\\ 
			& *\textit{the costed amount}\\
			& *\textit{the worded letter}
		\end{tabular}
		\ex 
		\begin{tabular}{ll}
			generalization: & \textit{er} nouns are formed from verbs that have an external argument.\\ 
			counter-examples: & *\textit{(the) meaner} (from \textit{she means it})\\ 
			& *\textit{(the) equaler} (from \textit{this equals that})\\
			& *\textit{(the) coster} (from \textit{this costs a lot})
		\end{tabular}
		\ex 
		\begin{tabular}{ll}
			generalization: & \textit{re} prefixes to verbs that have a direct object.\\ 
			counter-examples: & *\textit{reknow this fact}\\ 
			& *\textit{recost this amount}\\
			& *\textit{re\"entertain the idea}
		\end{tabular}
		\ex 
		\begin{tabular}{ll}
			generalization: & resultatives combine with verbs that have a direct object.\\ 
			counter-examples: & *\textit{She carried the cake stale.}\\ 
			& *\textit{She threw the child dizzy.}\\
			& *\textit{She studied this conjecture understood.}
		\end{tabular}
	\end{xlist}
\end{exe}
Unless we know how to give both necessary and sufficient conditions for these processes, we can expect that they will not necessarily pick out the same class of intransitives. In particular, then, we should not expect to find that any given unaccusative will form an adjectival participle, be amenable to \textit{re} prefixation, and license a resultative. Nor should we expect any given unergative to be able to form a noun with \textit{er}. The other factors relevant to determining when these processes can occur might apply. Since the argument that being unaccusative or unergative relies on these phenomena correlating---if \textit{re} and resultatives picked arbitrarily different verbs, we would have no reason to think that they are referencing the same attribute—--it is weakened when the correlations aren't perfect. For this reason, identifying the verbs that are unaccusative in any particular language has been difficult. And, consequently, finding generalizations about the kinds of verbs that can be unaccusative has also been slow-going. Perhaps the best candidates are the intransitives of the causative/inchoative, and other similar, alternations. 

That there are generalizations to be found in which intransitives are unergative and which are unaccusative is suggesed by the fact that the verbs in the causative/inchoative alternation do not always permit an intransitive version. There are instances of the transitive version which have no matching intransitive version, as (\ref{ex3:nocausincho}) shows.
\begin{exe}
\raggedright
	\ex \label{ex3:nocausincho}
	\begin{xlist}
		\ex[]{He broke his promise.}
		\ex[*]{His promise broke.}
		\ex[]{He broke the world record.}
		\ex[*]{The world record broke.}
		\ex[]{This book will open your mind.}
		\ex[*]{Your mind will open.}
		\ex[]{The waiter cleared the table.}
		\ex[*]{The table cleared.}
		\sn \textit{compare:}
		\ex[]{The wind cleared the sky.}
		\ex[]{The sky cleared.}
	\end{xlist}
	\hfill \citep[(9), (11)-(13):~85-6]{Levin:1995}\hfill\hspace*{.15in}
\end{exe}
So far as I know, this asymmetry doesn't go in the other direction. \citet{Levin:1995} argue that this indicates that the intransitive versions of these pairs are derived from the transitive versions.\marginnote{See also \citet{Dowty:1991}, \citet{Haspelmath:1993}, \citet{Hoekstra:1990}, and \citet{Goldberg:1994}.} The process that forms the intransitive from the transitive is sensitive to the particulars of the relation that holds between the verb and its internal argument. That sensitivity, however it might turn out to be formulated, might reflect constraints on the kinds of intransitives that can be unaccusative.

What I've laid out here are some of the ways in which we might discover the existence of unaccusative verbs in English. Many of these ways exist in other IndoEuropean languages as well. In each of these cases, the phenomena I've used involves some syntactic/semantic property that singles out direct objects or internal $\theta$-roles. The prefix \textit{re} and the resultative, for instance, require a direct object for its interpretation. And the adjectival participal requires a verb with an internal $\theta$-role while the nominalizing suffix \textit{er} requires a verb with an external $\theta$-role. There is another common method of distinguishing unaccusative and unergative verbs that, like \textit{er}, targets verbs with external $\theta$-roles, but which isn't found in English. This is the passive. 

The passive in English only applies to transitive verbs, but in many languages it can apply to intransitive verbs as well. This is the case in German, for instance. The intransitive verb \emph{tanzen} (`dance'), for instance, can be turned into a passive verb. The alternations in~(\ref{ex3:impersonal}) illustrate.
\begin{exe}
	\raggedright
	\ex \label{ex3:impersonal} \marginnote{These data are all from \citet{Grewendorf:1989}. Note that German forms the passive with the aid of a verb (\textit{werden} `become') that is different than that used in situations where adjectival participles arise. With adjectival participles, German uses the verb \textit{sein} (`be'), just as English does.}
	\begin{xlist}
		\ex \gll Sie tanzt. \\ she.\textsc{nom} dance.\textsc{pres}  \\ `she dances.'
		\ex \gll Es wurde getantzt. \\ it.\textsc{nom} become.\textsc{pres} dance.\textsc{pass}  \\ `it was danced.'
		\ex \gll Sie arbeitet.\\ she.\textsc{nom} work.\textsc{pres}\\ `she works.'
		\ex \gll Es wurde gearbeitet.\\ it.\textsc{nom} become.\textsc{pres} work.\textsc{pass}\\ `it was worked.'
		\ex \gll Sie schl\"aft.\\ she.\textsc{nom} sleep.\textsc{pres}\\ `She sleeps.'
		\ex \gll Es wurde geschlafen.\\ it.\textsc{nom} become.\textsc{pres} sleep.\textsc{pass}\\ `it was slept.'
	\end{xlist}
\end{exe}
The \emph{es} (`it') in these examples is an expletive. Passives of intransitive verbs are sometimes called ``impersonal passives.'' Not all intransitives can undergo the passive, however. Some cases are~(\ref{ex3:impersonal2}).
\begin{exe}
	\raggedright
	\ex \label{ex3:impersonal2}
	\begin{xlist}
		\ex[*]{ \gll Es wurde angekommen.\\ it.\textsc{nom} become.\textsc{pres} arrive.\textsc{pass}\\ `It was arrived.' }
		\ex[*]{\gll Es wurde gewaschen.\\ it.\textsc{nom} become.\textsc{pres} grow.\textsc{pass}\\ `It was grown.'}
		\ex[*]{\gll Es wurde gefallen.\\ it.\textsc{nom} become.\textsc{pres} fall.\textsc{pass}\\ `It was fallen.'}
	\end{xlist}
\end{exe}
The intransitives that can undergo the passive in German are classed as unergatives by independent criteria, and those which cannot undergo the passive are classed as unaccusatives by those criteria. For instance, German has the same rule that forms nouns from suffixing \textit{er} to verbs, and just as in English, it attaches only to verbs with an external $\theta$-role. In German, this suffix can attach to the intranstivies that passivize, but cannot attach to the intransitives that don't.
\begin{exe}
\raggedright
	\ex \label{ex3:germaner}
	\begin{xlist}
		\ex Arbeiter\\`worker'
		\ex T\"anzer\\`dancer'
		\ex Schl\"afer\\`sleeper'
	\end{xlist}
	\ex \label{ex3:germaner2}
	\begin{xlist}
		\ex[*]{Ankommer\\ `arriver'}
		\ex[*]{Wachser\\ `grower'}
		\ex[*]{Faller\\`faller'}
	\end{xlist}
\end{exe}
In the (many) languages that allow intransitives to passivize, there is a contrast of this sort between those intransitives that seem to be unaccusatives, and those that seem to be unergatives. David Perlmutter and Paul Postal\cite{Perlmutter:1984a} have suggested that it is a language universal that unaccusative verbs cannot passivize, and at the time of writing this remains a live possibility.

\subsection{Burzio's Generalization}

Now that we've established that accusative Case assignment and c-selection of DPs is independent, let's consider the predictions this makes about the types of verbs we should expect to find. We should find, and have, verbs that c-select a DP but do not assign accusative Case: passive participles and unaccusatives. We should also find, and have, verbs that c-select DPs and assign accusative Case: transitive verbs. We should expect to find verbs that do not c-select DPs but do assign accusative Case. One candidate for this type of verb is \emph{let}, which cannot combine with an accusative Case-marked DP complement, but can, plausibly, assign accusative Case to the subject of a small clause that it c-selects.
\begin{exe}
	\raggedright
	\ex \label{ex3:let} 
	\begin{xlist}
		\ex[*]{She let the dance.}
		\ex[]{She let [\VP the dance happen].}
	\end{xlist}
\end{exe}
If \emph{let} is responsible for assigning accusative Case in (\ref{ex3:let}b), then the ungrammaticality of (\ref{ex3:let}a) cannot be credited to the Case filter. We should conclude that \emph{let} does not c-select~DPs. Other possible examples of verbs of this kind are ones like \textit{hope}. Recall that these verbs do not have DP arguments, but they can be passivized.
\pagebreak
\begin{exe}
\raggedright
	\ex \label{ex3:hope}
	\begin{xlist}
		\ex[*]{This wish was hoped.}
		\ex[]{It was hoped that you'd understand this section.}
	\end{xlist}
\end{exe}
Since passivization generally happens in English only to verbs that assign accusative Case, we might be warranted to conclude from (\ref{ex3:hope}b) that \textit{hope} assigns accusative Case. Like \textit{let}, then, this might be a verb that assigns accusative Case but does not c-select a~DP. 

Finally, we should expect to find verbs that do not c-select DPs and do not assign accusative Case. A candidate for this type of verb is \emph{seem}. We have decided that \emph{seem} does not assign accusative Case, since this gives us an explanation for the contrast in~(\ref{ex3:raise}).
\begin{exe}
	\raggedright
	\ex \label{ex3:raise} 
	\begin{xlist}
		\ex[*]{It seems them to be happy.}
		\ex[]{They seem to be happy.}
	\end{xlist}
\end{exe}
The table in (\ref{ex3:table}) shows why structural accusative Case and c-selection of a DP are independent.
\begin{exe}
	\ex \label{ex3:table}
	\begin{tabular}{l|ll}  
 & Accusative Case & No Accusative Case \\ \hline
c-select DP  & \emph{like} (transitives) &  \emph{melt}, \emph{liked} (unaccusatives/passives) \\ %\hline
 $\neg$ c-select DP  & \emph{let, hope} & \emph{seem}\\ %\hline 
%\hline
\end{tabular}
\end{exe}

Interestingly, if we look for the same independence of accusative Case assignment and \underline{external} $\theta$-role assignment, it does not emerge. 
\begin{exe}
	\ex \label{ex3:tablen}
	\begin{tabular}{l|ll}  
\   & no external $\theta$-role & external $\theta$-role \\ \hline
accusative Case and c-select DP & ?? & \emph{like} (transitives) \\ %\hline
accusative Case and $\neg$ c-select DP  & ?? & \emph{let, hope}\\ %\hline 
no accusative Case and c-select DP & \emph{melt} (unaccusative) & ??\\ %\hline
no accusative Case and $\neg$ c-select DP & \emph{seem} & \emph{sleep}?\\ 
%\hline
\end{tabular}
\end{exe}
Some cells of this table are empty, suggesting that the availability of an external $\theta$-role type is not independent of the way that verbs can take ``objects.'' The cell occupied by \emph{sleep?} is not clear. We don't see \textit{sleep} combine with a DP complement nor do we, like \textit{let}, see it combine with a small clause whose subject argument is assigned accusative Case. This is consistent with it neither c-selecting DP nor assigning accusative Case. But it's also consistent with \textit{sleep} not c-selecting a DP complement but being able to assign accusative Case. In general, it's very difficult to tell if verbs of this kind exist. It could be that some unergative verbs are of this type; but unergatives could also be verbs that assign accusative Case, but do not c-select DP. They could occupy the cell occupied by \textit{let} and \emph{hope}. I'll ignore the status of this cell, given the difficulty in determining whether verbs of the type it describes exist.

The other gaps, however, indicate that there is a correlation between having accusative Case and having an external $\theta$-role. Burzio discovered these gaps in this paradigm where there appear to be no verbs, and formulated generalizations which describe these absences.
\pagebreak
\begin{exe}
	\ex \label{ex3:bg} Burzio's Generalization
	\begin{xlist}
	\ex[]{If a verb has structural Case, then it has an external {$\theta$-role}.}
	\ex[]{If a verb has an external $\theta$-role (and c-selects a DP?), then it has structural Case.}
	\end{xlist}
\end{exe}
Why aren't these verbs easy to find? If they exist, they have been very difficult to identify. Putative counter-examples are extremely rare. Among the best, perhaps, are German sentences like~(\ref{ex3:hungert}).
\begin{exe}
	\raggedright
	\ex \label{ex3:hungert} 
	\gll Dich hungert. \\ you.acc hunger.pres \\ `you are hungry.'
\end{exe}
There are a small handful of German verbs that allow their single argument to be accusative like \emph{hungern} does, in violation of the first clause of Burzio's Generalization. One conjecture is that in these situations the normal relationship between accusative and structural Case is lost; the accusative in (\ref{ex3:hungert}) is instead an inherent Case. 

Burzio's first generalization might reflect a relationship between Accusative Case and $\theta$-role assignment for which we should find a source.

The second of these generalizations might be derivable from the conditions we have seen on Argument movement, at least when it involves verbs that c-select DPs. Consider the sentences in which such a verb could exist. One of those is~(\ref{ex3:bgder1}). 
\begin{exe}
	\ex \label{ex3:bgder1} \ [\IP\ should [\VP\  Smith V Jones ]].
\end{exe}
If V assigns a $\theta$-role, but no Case to \emph{Jones}, then there is no way both of the DPs it selects are going to be able to satisfy the Case Filter. There are more DP arguments in (\ref{ex3:bgder1}) than there are Case-marked positions. So, if such a verb is to survive, the only environment it will be able to appear in are sentences which have two Case marked positions. These two positions will both have to be Specifiers, because these are the only positions reachable by Argument movement. Thus, we're looking for positions like V_2 in (\ref{ex3:bgder2}) for our verb to appear.
\begin{exe}
	\ex \label{ex3:bgder2} \ [\IP -- should V_1 [\XP\ -- [\VP\ Smith V_2 Jones]],\\
								 where both ``--'' positions are Case marked.
\end{exe}
Now the first part of Burzio's Generalization tells us that V_1 cannot assign Accusative Case. If it did, then it would also assign an external {$\theta$-role}, and that's going to bring the count of things that need Case to one more than there are Case marked positions. As a consequence, the Case marked position inside XP is going to also have to get its Case from some place other than V_1. So far, the only ways we have seen for this to be done are if XP is in a~CP, as in~(\ref{ex3:bgder3}).
\pagebreak
\begin{exe}
	\ex \label{ex3:bgder3}
	\begin{xlist}
	\sn {[\IP} -- should V_1 [\CP\ for [\IP\ -- to [\VP\ Smith V_2 Jones ]]]].
	\sn {[\IP} -- should V_1 [\CP\ that [\IP\ -- \obar{I} [\VP\ Smith V_2 Jones ]]]]
	\sn {[\IP} -- should V_1 [\CP\ e [\IP\ -- to [\VP\ \textsc{pro} V_2 Jones ]]]]
\end{xlist}	
\end{exe}
For these configurations to lead to a grammatical S-structure, one of the arguments in the embedded CP is going to have to A-move into the Specifier of the higher IP. Recall, however, that there is evidence that suggests that A-movement in English is blocked from taking an argument out of the CP it resides in. Consider, for instance, the derivation in~(\ref{ex3:cpisl}).
\begin{exe}
\raggedright
	\ex \label{ex3:cpisl}
	\small
	\begin{forest}
	[IP
		[DP [this, name=binder, roof]]
		[\ibar{I}
			[\obar{I} [past]]
			[VP [\ibar{V}
				[\obar{V} [seem]]
				[CP [\ibar{C}
					[\obar{C} [that]]
					[IP
						[DP [\ibar{D} [\obar{D} [it]]] ]
						[\ibar{I}
							[\obar{I} [should] ]
							[VP [\ibar{V}
								[\obar{V} [be]]
								[AP [\ibar{A}
									[\obar{A} [obvious]]
									[{\ }, name=trace ]
								] ]
							] ]
						 ]
					 ]
				] ]
			] ]
		 ]
	]
	\draw[->,>=stealth', thin, gray] (trace) to[out=south west,in=south,looseness=1] node[pos=0.75,xshift=-10mm]{} (binder);
	\end{forest}
\end{exe}
We've seen that \textit{obvious} c-selects a DP argument, but does not assign it accusative Case. In (\ref{ex3:cpisl}), A-movement has taken the c-selected DP (=\textit{this}) into the nominative Case marked position in the higher IP. This should lead to a grammatical S-structure, but it doesn't. This contrasts with the grammatical derivation in~(\ref{ex3:cpislg}).
\begin{exe}
\raggedright
	\ex \label{ex3:cpislg}
	\small
	\begin{forest}
	[IP
		[DP [\ibar{D} [\obar{D} [it]]] ]
		[\ibar{I}
			[\obar{I} [past]]
			[VP [\ibar{V}
				[\obar{V} [seem]]
				[CP [\ibar{C}
					[\obar{C} [that]]
					[IP
						[DP [this, name=binder, roof]]
						[\ibar{I}
							[\obar{I} [should]]
							[VP [\ibar{V}
								[\obar{V} [be]]
								[AP [\ibar{A}
									[\obar{A} [obvious]]
									[{\ }, name=trace ]
								] ]
							] ]
						 ]
					 ]
				] ]
			] ]
		 ]
	]
	\draw[->,>=stealth', thin, gray] (trace) to[out=south west,in=south,looseness=1] node[pos=0.75,xshift=-10mm]{} (binder);
	\end{forest}
\end{exe}
The difference between (\ref{ex3:cpisl}) and (\ref{ex3:cpislg}) is whether the A-moved DP has been moved out of the embedded CP. It's facts like these that led us to posit~(\ref{ex3:tensedS}).
\begin{exe}
\raggedright
	\exr{ex3:tensedS} If Argument Movement moves XP from position $\alpha$ to position $\beta$, then there cannot be a CP that dominates $\alpha$ but not~$\beta$.
\end{exe}
This condition will prevent any of the scenarios in (\ref{ex3:bgder3}) from having a grammatical S-structure. So this is also an environment in which a verb of the sort we are missing could not survive.

Finally, consider the situation that arises if such a verb were to select a CP external argument:\marginnote[0pt]{\citet{Hartman:2012} argues that configurations like (\ref{ex3:bgder4}) arise for adjectives, however. He argues that sentences like \emph{Marcin is annoyed that pigs fly} involve an external argument CP and a internal argument DP that has moved to the nominative Case-marked position. There is more about the syntax of these sentences needed to get the word-order right, of course. But if he is correct, then the absence of verbs of this kind cannot be credited to constraints on A Movement in the way sketched here.}
\begin{exe}
	\ex \label{ex3:bgder4} \ [\IP\ should [\VP\  CP V Jones ]].
\end{exe}
Here, only the argument DP \emph{Jones} needs to find a Case-marked position. From what we presently know about A movement, this should be possible: \emph{Jones} should be able to move into the Specifier of IP by A Movement and there satisfy the Case filter. If there is a place that such verbs could exist, this is where they would survive. We will see in the next chapter evidence for a constraint on A movement that would prevent these environments from being grammatical as well. That constraint prevents one argument from moving past another by A Movement under certain circumstances. That constraint might be responsible for blocking this environment as well.

For this reason, the usual conclusion is that the second of Burzio's generalizations derives from constraints on Argument movement in concert with the Case filter. What remains is~(\ref{ex3:finalbg}).
\begin{exe}
\ex \label{ex3:finalbg} Burzio's Generalization\\
If a verb assigns structural Case, then it assigns an external \mbox{$\theta$-role}.
\end{exe}
If Burzio's Generalization cannot be derived from constraints on the syntactic form that such verbs would have to live in, then it suggests that the assignment of Case and \mbox{$\theta$-roles} is not independent in the way that we presently think it is.

Indeed, one way of describing Burzio's Generalization is that the external \mbox{$\theta$-role} \underline{is} the assigner of structural (accusative) Case. It isn't, as we originally thought, that the categories of words that assign accusative Case in English are prepositions and verbs. The accusative Case assigners of English are instead prepositions and whatever part of speech that external $\theta$-roles belong to. A way of giving a syntax for that thesis involves making the external $\theta$-role syntactically external to the verb that it is associated with. This is the position taken by Angelika Kratzer,\cite{Kratzer:1996} who uses the symbol ``v'' (`little \textipa{[vij]}') to indicate the category of an external \mbox{$\theta$-role}.
\begin{exe}
	\raggedright
	\ex \label{ex3:lv1} 
	\small
	\begin{forest}
	[vP
		[DP [Pius, roof]]
		[\ibar{v} 
			[\obar{v} [agent] ]
			[VP [\ibar{V} 
				[\obar{V} [eat] ]
				[DP [them, roof] ]
			] ]
		]
	]
	\end{forest}
\end{exe}
This syntax puts the external $\theta$-role higher than the verb it is associated with, and consequently puts the argument the external $\theta$-role combines with higher too.\marginnote{That external $\theta$-roles are related to a head that is higher than the VP is an idea can also be found in \citet{Williams:1980}, \citet{Hornstein:1987}, \citet{Bowers:1993}, and \citet{Dikken:1992}.} By placing the external $\theta$-role higher than the VP, this syntax also allows it to c-command the DP it allegedly assigns accusative Case to. In these ways, the proposed syntax in (\ref{ex3:lv1}) fits with the syntax we've developed.

But there are several ways it does not. Our present understanding of accusative Case assignment is that it requires the Case assigner to not only c-command the position assigned Case, but to also be adjacent to it and no farther from that position than (\ref{ex3:gg}) describes.
\begin{exe}
	\raggedright
	\ex \label{ex3:gg} The position assigned accusative Case by \obar{X} cannot be dominated by more than one phrase that does not also dominate~\obar{X}.
\end{exe}
(\ref{ex3:gg}) is one of the requirements of ``government.'' In (\ref{ex3:lv1}), \obar{v} meets neither of these conditions: it is not adjacent to \emph{them} and there is more one phrase that dominates \emph{them} that doesn't also dominate~\obar{v}. If (\ref{ex3:lv1}) is correct, then we will have to revise our Case theory. 

The semantics of $\theta$-role assignment, and how it interacts with the syntax, also does not fit (\ref{ex3:lv1}). Our semantics for \emph{eat} gave it the denotation in~(\ref{ex3:eatda}).
\begin{exe}
	\raggedright
	\ex \label{ex3:eatda} 
	\denotes{eat} = $\lambda x$ $\lambda y$ $\lambda \epsilon$ \textsc{\underline{agent}}($y)(\epsilon$) $\wedge$ \textsc{patient}($x)(\epsilon$) $\wedge$ \textsc{eat}($\epsilon$)=T, where
	\begin{xlist}
	\ex $x$ and $y$ are things and $\epsilon$ is an eventuality
	\ex \textsc{patient} c-selects DP
	\end{xlist}
\end{exe}
On this view, $\theta$-roles are functions from things to functions that describe eventualities. If we are to separate the Agent $\theta$-role from (\ref{ex3:eatda}) in the way that (\ref{ex3:lv1}) requires, then the denotations we get from simply doing this to (\ref{ex3:eatda}) is~(\ref{ex3:eatda2}).
\begin{exe}
	\raggedright
	\ex \label{ex3:eatda2} 
	\begin{xlist}
		\ex \denotes{agent} = $\lambda y$ $\lambda \epsilon$ \textsc{agent}($y)(\epsilon$)=T, where
		\begin{xlist}
		\ex $y$ is a thing and $\epsilon$ is an eventuality
		% \ex \denotes{\textsc{agent}} c-selects VP
		\end{xlist}
		\ex \denotes{eat} = $\lambda x$ $\lambda \epsilon$ \textsc{patient}($x)(\epsilon$) $\wedge$ \textsc{eat}($\epsilon$)=T, where
		\begin{xlist}
		\ex $x$ is a thing and $\epsilon$ is an eventuality
		\ex \textsc{patient} c-selects DP
		\end{xlist}
	\end{xlist}
\end{exe}
The denotation of Agent in (\ref{ex3:eatda2}) remains a function from things, $x$, to functions that describe eventualities. But if we plug these denotations into (\ref{ex3:lv1}), the rules of semantic composition we adopted in (\ref{ex3:perco}) do not work.
\begin{exe}
	\raggedright
	\exr{ex3:perco} 
	\begin{xlist}
		\ex If X is the sole daughter of Z, then \denotes{Z} = \denotes{X}
		\ex If X and Y are daughters of Z, then \denotes{Z} is the result of applying the function \denotes{X} to its argument~\denotes{Y}.
	\end{xlist}
	\ex \label{ex3:eatda2bt}
	\small
	\begin{forest}
	[vP
		[DP [Pius, roof]]
		[??\\ \ibar{v} 
			[$\lambda y$ $\lambda \epsilon$ \textsc{agent}($y)(\epsilon$)\\ \obar{v} [agent] ]
			[$\lambda \epsilon$ \textsc{patient}(\denotes{it})($\epsilon$) $\wedge$ \textsc{eat}($\epsilon$)\\ VP [$\lambda \epsilon$ \textsc{patient}(\denotes{it})($\epsilon$) $\wedge$ \textsc{eat}($\epsilon$)\\ \ibar{V} 
				[$\lambda x$ $\lambda \epsilon$ \textsc{patient}($x)(\epsilon$) $\wedge$ \textsc{eat}($\epsilon$)\\ \obar{V} [eat] ]
				[DP [it, roof] ]
			] ]
		]
	]
	\end{forest}
\end{exe}
There is no rule in (\ref{ex3:perco}) that can combine the meanings of \obar{v} and VP in~(\ref{ex3:eatda2bt}).

Kratzer's solution to this is a new rule of semantic composition, which she calls ``Event Identification.''
\begin{exe}
	\raggedright
	\ex \label{ex3:ev} Event Identification\\
	In
	\footnotesize
	\begin{forest}
	[h
		[f ]
		[g ]
	]
	\end{forest}\normalsize\hspace*{4pt} \parbox[t]{3in}{if \denotes{f} is $\lambda x.$$\lambda \epsilon.P(x)(\epsilon)$ and \denotes{g} is $\lambda \epsilon.Q(\epsilon)$,\\[3pt] then \denotes{h} is $\lambda x.$$\lambda \epsilon.[P(x)(\epsilon) \wedge Q(\epsilon)]$.}	
\end{exe}
Making use of this rule allows for the semantics for~(\ref{ex3:eatda2bt}) indicated in~(\ref{ex3:eatda2bd}).
\begin{exe}
	\ex \label{ex3:eatda2bd}
	\small
	\begin{forest}
	[$\lambda \epsilon$. \textsc{agent}(\denotes{Pius})($\epsilon$) $\wedge$ \textsc{eat}($\epsilon$) $\wedge$ \textsc{patient}(Jones)($\epsilon$)\\ vP
		[DP [Pius, roof] ]
		[$\lambda x.\lambda \epsilon$ \textsc{agent}($x)(\epsilon$) $\wedge$ \textsc{eat}($\epsilon$) $\wedge$ \textsc{patient}(\denotes{them})($\epsilon$)\\ \ibar{v} 
			[$\lambda x.\lambda \epsilon$ \textsc{agent}($x)(\epsilon$)//\obar{v} [agent]]
			[$\lambda \epsilon.$ \textsc{eat}($\epsilon$) $\wedge$ \textsc{patient}(\denotes{them})($\epsilon$)\\ VP [\ibar{V} 
				[$\lambda x.\lambda \epsilon$. \textsc{eat}($\epsilon$) $\wedge$ \textsc{patient}($x$)($\epsilon$)\\ \obar{V} [eat] ]
				[DP [them, roof] ]
			] ]
		]
	]
	\end{forest}
\end{exe}

Another possibility is that we could imagine that $\theta$-roles which are put into the syntax independently of the verb have a different denotation. One that is found often in the literature is~(\ref{ex3:agentfn}).
\begin{exe}
	\ex \label{ex3:agentfn}
	\denotes{agent} = $\lambda P.\lambda x.\lambda \epsilon$ \textsc{agent}($x)(\epsilon$) $\wedge$ $P(\epsilon)=T$, where
	\begin{xlist}
	\ex $P$ is a function whose domain is eventualities, $x$ is a thing, and $\epsilon$ is an eventuality
	\ex agent c-selects VP.
	\end{xlist}
\end{exe}
With this denotation, things compose up as indicated in~(\ref{ex3:eatdent})
\begin{exe}
	\ex \label{ex3:eatdent}
	\small
	\begin{forest}
	[$\lambda \epsilon$. \textsc{agent}(\denotes{Pius})($\epsilon$) $\wedge$ \textsc{visit}($\epsilon$) $\wedge$ \textsc{patient}(\denotes{it})($\epsilon$)// vP
		[DP [Pius, roof] ]
		[$\lambda x.\lambda \epsilon$ \textsc{agent}($x)(\epsilon$) $\wedge$ \textsc{visit}($\epsilon$) $\wedge$ \textsc{patient}(\denotes{it})($\epsilon$)\\ \ibar{v} 
			[$\lambda P.\lambda x.\lambda \epsilon$ \textsc{agent}($x)(\epsilon$) $\wedge\ P(\epsilon)$\\ \obar{v} ]
			[$\lambda \epsilon.$ \textsc{visit}($\epsilon$) $\wedge$ \textsc{patient}(\denotes{it})($\epsilon$)\\ VP [\ibar{V} 
				[$\lambda x.\lambda \epsilon$. \textsc{eat}($\epsilon$) $\wedge$ \textsc{patient}($x)(\epsilon$)\\ \obar{V} [eat] ]
				[DP [it, roof] ]
			] ]
		]
	]
	\end{forest}
\end{exe}
This view doesn't require Event Identification; the rule of function composition is sufficient. Both of these techniques yield the same denotation for vP, and both are found in the literature.

Giving external $\theta$-roles denotations like those in (\ref{ex3:eatda2}a) or (\ref{ex3:agentfn}) has certain interesting consequences when combined with this new syntax. It gives an explanation for why external $\theta$-roles bearers are found higher than other arguments. Consider, for instance, what would happen if \obar{v} were to be placed between a verb and another DP it assigns a $\theta$-role to and its $\theta$-role has a denotation like that in~(\ref{ex3:agentfn}). A way of doing that which the \ibar{X} Skeleton permits is~(\ref{ex3:nolo}).
\begin{exe}
	\raggedright
	\ex \label{ex3:nolo} 
	\small
	\begin{forest}
	[vP
		[DP [it, roof]]
		[\ibar{v} 
			[??\\ \ibar{v} 
				[$\lambda P$ $\lambda x$ $\lambda \epsilon$ \textsc{agent}($x)(\epsilon$) $\wedge$ $P(\epsilon)$\\ \obar{v} [agent]]
				[$\lambda y$ $\lambda \epsilon$ \textsc{patient}($y)(\epsilon$) $\wedge$ \textsc{eat}($\epsilon$)\\ VP [\ibar{V} [\obar{V} [eat]]] ]
			]
			[DP [Pius, roof] ]
		]
	]
	\end{forest}
\end{exe}
The denotation for VP is not of the right kind to be an argument of \obar{v}, nor is the denotation of \obar{v} the right kind to be an argument of VP. In general, if the verb heading VP has a $\theta$-role that hasn't yet been assigned, then it won't be of the right semantic kind to be an argument of~\obar{v}. The same consequence arises if we use the denotation for external $\theta$-roles of the sort given in (\ref{ex3:eatda2}), though that case is complicated by the existence of other semantic composition rules we have not yet encountered. Let's somewhat arbitrarily adopt the denotation for external $\theta$-roles given in~(\ref{ex3:agentfn}).

With this semantics, then, we can now simplify the Projection Principle. We need no longer stipulate that the external $\theta$-role is assigned in a different way than all the others. This will follow from the thesis that external $\theta$-roles are placed syntactically in a different position than the verbs they are associated with, and the rules of semantic composition we've adopted. The Projection Principle can now be rendered as~(\ref{ex3:PPfin}).
\begin{exe}
	\raggedright
	\ex \label{ex3:PPfin} The Projection Principle\\
	If $\beta$ is c-selected by $\alpha$, then there must be a projection of $\alpha$, $\pi$, that immediately dominates $\beta$ and if $\pi$ immediately dominates some other phrase, $\gamma$, then $\gamma$ must be another projection of $\alpha$ or part of another c-selected argument of~$\alpha$.
\end{exe}

Noam Chomsky\cite{Chomsky:1965} argued that the Projection Principle is part of the process of lexical insertion. Lexical Insertion puts into \obar{X} positions the morphemes that will be used to make the sentence. In our system, Lexical Insertion also determines the category of the \obar{X} position it is inserted into, as well as all of the projections of that \obar{X}. Let's derive this and the Projection Principle, then, by writing out how Lexical Insertion has these effects.
\pagebreak
\begin{exe}
	\raggedright
	\ex \label{ex3:LI} Lexical Insertion\\
	Put morpheme $\mu$ into \obar{X}, only if\marginnote{These conditions have effects which are frequently assumed in the literature. But the conditions themselves are completely parochial. They are my own modernization of \citet{Chomsky:1965}.}
	\begin{xlist}
		\ex The morpho-syntactic category of \obar{X} and its projections is the same as the morphosyntactic category of $\mu$, and
		\ex Every projection of \obar{X} satisfies the c-selection requirements of $\mu$.
		\sn Let $c$ be a c-selection requirement of $\mu$. $\alpha$ \underline{satisfies} $c$, iff
		\begin{xlist}
			\ex $\alpha$ doesn't have an adopted daughter whose morpho-syntactic category is different from what c requires, or ii. 
			\ex c is satisfied by the phrase that projects $\alpha$.
		\end{xlist}
		\sn $\alpha$ is an adopted daughter of $\beta$ if $\alpha$ is a daughter of $\beta$ and $\beta$ is not a projection of~$\alpha$.
		\end{xlist}
\end{exe}
What an ``adopted daughter'' is can be graphically described. Consider a case like~(\ref{ex3:adopted}).
\begin{exe}
\raggedright
	\ex \label{ex3:adopted}
	\begin{forest}
	[XP
		[YP ]
		[\ibar{X}
			[\obar{X} ]
			[ZP ]
		 ]
	]
	\end{forest}
\end{exe}
The adopted daughter of \ibar{X} is ZP; \obar{X} is not an adpted daughter of \ibar{X} because \ibar{X} is a projection of~\obar{X}. Similarly, of the two daughters that XP has, only YP is an adopted one because XP is a projection of \ibar{X}.

The definition of ``satisfaction'' in (\ref{ex3:LI}b) has complex consequences. It lets the c-selection requirements that comes with a verb's $\theta$-role be satisfied in two ways. One way is to make the phrase that is c-selected an adopted daughter of one of the verb's projections. Thus, if XP in (\ref{ex3:satisfy}) is the category that the c-selection requirement of one of the verb's $\theta$-roles calls for, then that c-selection requirement is satisfied.
\begin{exe}
\raggedright
	\ex \label{ex3:satisfy}
	\begin{forest}
	[\ibar{V}
		[\obar{V} ]
		[XP ]
	]
	\end{forest}
\end{exe}
The definition of satisfaction in (\ref{ex3:LI}) requires that if a projection of a verb has an adopted daughter, here XP, then it cannot be a different category than what the c-selection requirements call for. If, by contrast, XP isn't the category that the c-selection requirements of that $\theta$-role have, \ibar{V} will not satisfy the c-selection requirements. Since (\ref{ex3:LI}b) requires that all projections of the verb satisfy the c-selection requirements, if XP is not the category that one of the $\theta$-roles c-selects, the \ibar{V} will be ungrammatical. This has the effect of allowing XP to be a phrase that isn't the category required by a $\theta$-role of the verb only if the phrases within \ibar{X} have satisfied those requirements. This derives the generalization we had previously about arguments of a verb: they must be found in a projection of that verb that contains no non-arguments.

But the definition of satisfaction in (\ref{ex3:LI}) permits another way for the c-selection requirements of a $\theta$-role associated with a verb to be satisfied. This way is illustrated in~(\ref{ex3:satisfy2}).
\begin{exe}
\raggedright
	\ex \label{ex3:satisfy2}
	\begin{forest}
	[\ibar{V}^*
		[\ibar{V} [\obar{V}] ]
	]
	\end{forest}
\end{exe}
Suppose that \obar{V} has one $theta$-role that is accompanied by a c-selection. \ibar{V} satisfies this c-selection requirement because it has no adopted daughter whose category differs from what that c-selection requirement calls for. In this scenario, the c-selection requirements are satisfied because \ibar{V} has no adopted daughter. Because this is a way that the c-selection requirements can be satisfied, and because the graph in (\ref{ex3:satisfy2}) is always available, this means that c-selection requirements are not responsible for ensuring that verbs (and other argument takers) are supplied with the arguments that c-select. What c-selection requirements do is merely ensure that when those arguments are provided, they be in the right positions and of the right category. That the arguments must be provided will come from the semantics, not the c-selection requirements.

Let's run through how this works in a couple of cases. Our first case involves inserting \emph{put} into the phrase built by the \ibar{X} Skeleton in~(\ref{ex3:template1}).
\begin{exe}
	\raggedright
	\ex \label{ex3:template1} 
	\small
	\begin{forest}
	[XP [\ibar{X}
		[\ibar{X} 
			[\obar{X} ]
			[DP ]
		]
		[PP ]
	]]
	\end{forest}
\end{exe}
Assume that the denotation of \emph{put} is~(\ref{ex3:putden}).
\begin{exe}
	\raggedright
	\ex \label{ex3:putden} 
	\denotes{put} = $\lambda x$ $\lambda y$ $\lambda \epsilon$ \textsc{theme}($x$)($\epsilon$) $\wedge$ \textsc{location}($y$)($\epsilon$) $\wedge$ \textsc{put}($\epsilon$)=T, where
	\begin{xlist}
		\ex $x$ and $y$ are things and $\epsilon$ is an eventuality
		\ex \textsc{theme} c-selects DP
		\ex \textsc{location} c-selects PP
	\end{xlist}
\end{exe}
Lexical Insertion of \emph{put} into the \obar{X} position of (\ref{ex3:template1}) is allowed if all the X's are re-written as V. Because \emph{put} c-selects a DP and a PP, the sisters of \obar{X} and the smallest \ibar{X} must be DP and PP, if they exist, and this requirement is met in (\ref{ex3:template1}). If either DP or PP were different morpho-syntactic categories, the c-selection requirements of \emph{put} would not be satisfied and lexical insertion would be blocked.
\pagebreak 

(\ref{ex3:template1}) is also a structure that \emph{eat} could be lexically inserted into. 
\begin{exe}
	\raggedright
	\ex \label{ex3:eatdena} 
	\denotes{eat} = $\lambda x$ $\lambda \epsilon$ \textsc{patient}($x$)($\epsilon$) $\wedge$ \textsc{eat}($\epsilon$)=T, where
	\begin{xlist}
		\ex $x$ is a thing and $\epsilon$ an eventuality
		\ex \textsc{patient} c-selects DP
	\end{xlist}
\end{exe}
The c-selection requirement of \emph{eat} is DP, and this is satisfied by virtue of DP being a sister to \obar{X}. Because the DP satisfies the one c-selection requirement of \emph{eat}, the sister of the smallest \ibar{X} need not satisfy the c-selection of \emph{eat}. As a consequence any kind of phrase, including PP, can reside in this position. Indeed, nothing need reside in this position.

Consider next what happens if \emph{eat} is inserted into the structure in~(\ref{ex3:template2}).
\begin{exe}
	\raggedright
	\ex \label{ex3:template2} 
	\small
	\begin{forest}
	[VP[\ibar{V}
		[$\lambda \epsilon$ $\exists x$ \textsc{patient}($x$)($\epsilon$) $\wedge$ \textsc{eat}($\epsilon$)\\ \obar{V} [eat]]
		[PP ]
	]]
	\end{forest}
\end{exe}
Because Existential Closure is a rule that can change the denotation of \emph{eat}, we can create a VP headed by \emph{eat} that satisfies the Theta Criterion without including a phrase that gets \emph{eat}'s \textsc{patient} $\theta$-role. I've applied Existential Closure to \emph{eat} in (\ref{ex3:template2}) so that we can consider how things work under this scenario. As it happens, Lexical Insertion is violated in (\ref{ex3:template2}). The c-selection requirement of \emph{eat} is not satisfied by the XP that is immediately dominated by \ibar{V}, and this is what (\ref{ex3:LI}b) requires.

Consider, by way of contrast, what happens if the Existentially Closed \emph{eat} is inserted into the structure in~(\ref{ex3:template3}).
\begin{exe}
	\raggedright
	\ex \label{ex3:template3} 
	\small
	\begin{forest}
	[VP[\ibar{V}
		[$\lambda \epsilon$ $\exists x$ \textsc{patient}($x$)($\epsilon$) $\wedge$ \textsc{eat}($\epsilon$)\\ \ibar{V} [$\lambda \epsilon$ $\exists x$ \textsc{patient}($x$)($\epsilon$) $\wedge$ \textsc{eat}($\epsilon$)\\\obar{V} [eat]]]
		[PP ]
	]]
	\end{forest}
\end{exe}
This \underline{does} satisfy the conditions on Lexical Insertion. The immediate projection of \obar{V} does not contain an XP whose morpho-syntactic category differs from the c-selection of \emph{eat}, and this meets the requirements of (\ref{ex3:LI}b). In general, this definition of Lexical Insertion derives the following generalization.
\begin{exe}
	\raggedright
	\ex \label{ex3:ppp} 
	The sister to the \obar{X} that a morpheme, $\mu$, is lexically inserted into, will have to meet a c-selection requirement of $\mu$ if there is one.
\end{exe}
It is common practice to regard only phrase markers that satisfy (\ref{ex3:ppp}) to be well-formed. Because the difference between (\ref{ex3:template2}) and (\ref{ex3:template3}) is extremely difficult to detect, I do not know if this practice is sound. I will adopt it in this book.\marginnote{There may be some evidence for this conclu-sion from prosody. In particular, the verb before an argument PP has the prosodic status of a word, whereas the verb before a non-argument PP has the prosodic status of a phrase. See \citet{Gussenhoven:1984}.}

If Existential Closure does not apply in (\ref{ex3:template3}), the result will still satisfy the conditions on Lexical Insertion, but the semantics will fail.
\begin{exe}
	\raggedright
	\ex \label{ex3:template4} 
	\small
	\begin{forest}
	[VP[\ibar{V}
		[$\lambda x$ $\lambda \epsilon$ \textsc{patient}($x$)($\epsilon$) $\wedge$ \textsc{eat}($\epsilon$)\\ \ibar{V} [$\lambda x$ $\lambda \epsilon$ \textsc{patient}($x$)($\epsilon$) $\wedge$ \textsc{eat}($\epsilon$)=T\\ \obar{V} [eat]]]
		[PP ]
	]]
	\end{forest}
\end{exe}
The PP will become the argument of \textsc{patient} and we might safely assume that the denotations of PPs do not allow them to be patients.

With a semantics that has the external $\theta$-role c-select a VP, we have been able to simplify the Projection Principle. And this simplified Projection Principle can be derived from the conditions on Lexical Insertion in (\ref{ex3:LI}). Let's therefore regard the semantic issues that arise with the new syntax of external $\theta$-roles solved, and consider some its consequences. We now have an account of alternations like that in~(\ref{ex3:melta}).
\begin{exe}
	\raggedright
	\ex \label{ex3:melta} 
	\begin{xlist}
		\ex Janina melted the ice.
		\ex The ice melted.
	\end{xlist}
\end{exe}
This is a situation where the external $\theta$-role associated with \emph{melt} is simply optional. There are two D-structure configurations available, and they give rise to the two derivations in~(\ref{ex3:vaccill}).
\begin{exe}
	\ex \label{ex3:vaccill}
	\begin{xlist}
		\ex
		\small
		\begin{forest}
		[IP
			[DP [Janina, roof, name=binder]]
			[\ibar{I} 
				[\obar{I} [past] ]
				[vP
					[{\ }, name=trace ]
					[\ibar{v}
						[\obar{v} [agent] ]
						[VP [\ibar{V} 
							[\obar{V} [melt] ]
							[DP [the ice, roof] ]
						] ]
					 ]
				 ]
			]
		]
		\draw[->,>=stealth', ultra thin, gray] (trace) to[out=south west,in=south,looseness=1] node[pos=0.75,xshift=-10mm]{} (binder);
		\end{forest}
		\vspace*{-35pt}
		\ex
		\small
		\begin{forest}
		[IP
			[DP [the ice, roof, name=binder]]
			[\ibar{I} 
				[\obar{I} [past] ]
				[VP [\ibar{V} 
					[\obar{V} [melt] ]
					[{\ }, name=trace ]
				] ]
			]
		]
		\draw[->,>=stealth', ultra thin, gray] (trace) to[out=south west,in=south,looseness=1] node[pos=0.75,xshift=-10mm]{} (binder);
		\end{forest}
	\end{xlist}
\end{exe}
In (\ref{ex3:vaccill}a), the D-structure comes equipped with an Agent $\theta$-role, and this \obar{v} assigns structural (accusative) Case. For this reason, a subject argument (=\emph{Janina}) must be included in the D-structure and the object DP sits at D-structure in a Case-marked position. In (\ref{ex3:vaccill}b), by contrast, the D-structure comes without an external $\theta$-role, and consequently there is no subject argument in this D-structure. Without the external $\theta$-role, there is no structural Case available to the DP object, and so it will not satisfy the Case filter in its D-structure position. This is remedied by A Movement bringing that DP into a Case-marked position, in this example that is the nominative Case-marked Specifier of~IP.

Consider next the Passive alternation in~(\ref{ex3:meltpass}).\marginnote[0pt]{Recall that forms like (\ref{ex3:meltpass}b) are ambiguous. The word \emph{melted} can be a passive verb or it can be an adjective derived morphologically from the verb \emph{melt}. When an adjective, the reading it gives is stative. The eventualities that adjectives describe are states. We are interested here only in the Passive version of (\ref{ex3:meltpass}b). The passive has an eventive reading -- a non-stative interpretation -- and I have tried to help disambiguate in that direction with the adverb \emph{quickly}.}
\begin{exe}
	\raggedright
	\ex \label{ex3:meltpass} 
	\begin{xlist}
		\ex Janina melted the ice.
		\ex The ice was melted (quickly).
	\end{xlist}
\end{exe}
Our present conception of Passive is that it is a rule that builds a lexicon. It creates new verbs from old ones in the way that (\ref{ex3:passold}) describes. (I've now expressed the Passive as a function whose domain and range are verbs.)
% \enlargethispage*{\baselineskip}
\begin{exe}
	\raggedright
	\ex \label{ex3:passold} Passive(V) = V$'$
	\begin{xlist}
		\ex Put V into the morphological form appropriate for passive, and
		\ex Suppress the external $\theta$-role of V, and
		\ex Remove the structural Case assigned by V.
	\end{xlist}
\end{exe}
In languages like English, the change of morphological form must somehow be tied to the fact that the Passive forms of verbs are always embedded inside a VP headed by \emph{be}.\marginnote[0pt]{Passives, like the one that English has, that involve more than a simple change to the morphology of the verb are sometimes called ``periphrastic'' Passives.} This no longer fits our syntax. Under the present hypothesis, verbs never assign structural Case and they do not come with external $\theta$-roles either. We will need to revise our view of the Passive.

In doing this, let's focus on what the ``suppression'' of the external $\theta$-role part of the rule consists of. It is useful here to compare the passive version of \emph{melt} to its intransitive version in~(\ref{ex3:melta}b).
\begin{exe}
	\raggedright
	\ex \label{ex3:passinc} 
	\begin{xlist}
		\ex The ice melted. (=(\ref{ex3:melta}b))
		\ex The ice was melted (quickly).
	\end{xlist}
\end{exe}
Our account of the intransitive \emph{melt} involved giving it a D-structure with no external $\theta$-role. The passive \emph{melt}, however, arguably does contain an external $\theta$-role. One difference between the sentences in (\ref{ex3:passinc}) is that the passive (\ref{ex3:passinc}b) entails that there is an Agent of the melting eventuality it describes, whereas that isn't the case for the intransitive \emph{melt} in~(\ref{ex3:passinc}a). We can see this by considering the sentences in~(\ref{ex3:entails}) and~(\ref{ex3:entailsb}).
\begin{exe}
	\raggedright
	\ex \label{ex3:entails} 
	\begin{xlist}
		\ex All the ice melted yesterday because everyone worked hard.
		\ex All the ice was melted yesterday because everyone worked hard.
	\end{xlist}
	\ex \label{ex3:entailsb}
	\begin{xlist}
		\ex The ice melted yesterday, but I didn't do it.
		\ex The ice was melted yesterday, but I didn't do it.
	\end{xlist}
\end{exe}
The second clause in (\ref{ex3:entails}a) and (\ref{ex3:entailsb}a) sound disconnected from the first clause. But in (\ref{ex3:entails}b) and (\ref{ex3:entailsb}b), they sound connected to the first clause. This is because \emph{everyone} can be understood to be quantifying over the Agents of the melting eventualities in (\ref{ex3:entails}b) but not (\ref{ex3:entails}a), and because \emph{do it} can be understood to refer to an eventuality with an Agent in (\ref{ex3:entailsb}b) but not (\ref{ex3:entailsb}a). We can also see this difference by considering how adverbs that modify Agents work. One of those adverbs is \emph{intentionally}; in(\ref{ex3:hug}), \emph{intentionally} describes the attitude the Agent, but not the Patient, has to the hugging eventuality.
\begin{exe}
	\raggedright
	\ex \label{ex3:hug} 
	Elena hugged Pius intentionally
\end{exe}
This adverb also detects an Agent $\theta$-role in the passive, but not the intransitive, use of \emph{melt}.\marginnote[0pt]{To the extent that (\ref{ex3:intentmelt}a) is grammatical, it requires an interpretation in which \emph{the ice} is an Agent of the melting eventuality; it means something roughly like \emph{the ice melted itself intentionally}. That is not the intransitive \emph{melt} we are concerned with.}
\begin{exe}
	\raggedright
	\ex \label{ex3:intentmelt} 
	\begin{xlist}
		\ex[*]{The ice melted intentionally.}
		\ex[]{The ice was melted intentionally.}
	\end{xlist}
\end{exe}

We do not want the Passive to remove the Agent $\theta$-role, in the same way that has been done in forming the intransitive. This is the point made by Mark Baker and Ian Roberts in a paper from the 1980's.\cite{Baker:1989a} Instead, what it seems we need is a way for the Agent $\theta$-role to undergo Existential Closure. Recall that Existential Closure is an operation that takes certain verbs and makes them ``intransitive'' by binding the argument position of their internal $\theta$-role with an existential operator. From transitive \emph{eat}, for instance, it can create intransitive~\emph{eat}.
\begin{exe}
	\raggedright
	\ex \label{ex3:eatec}
	\begin{xlist}
		\ex 
		\denotes{eat} = $\lambda x$ $\lambda \epsilon$ \textsc{patient}($x$)($\epsilon$) $\wedge$ \textsc{eat}($\epsilon$) \hfill{transitive}
		\ex \small
		\denotes{eat} = $\lambda \epsilon$ $\exists x$ \textsc{patient}($x$)($\epsilon$) $\wedge$ \textsc{eat}($\epsilon$) \hfill{intransitive}
	\end{xlist}
\end{exe}

We noted at the time that this operation must be limited to just internal $\theta$-roles. We don't want to create intransitive verbs with this process by letting it apply to external $\theta$-roles. If that were possible, we would expect to find transitive versions of \emph{melt} with their Agent $\theta$-role suppressed. These verbs would appear in sentences like~(\ref{ex3:nobg2}).
\begin{exe}
	\raggedright
	\ex[*]{\label{ex3:nobg2}%
	It melted the ice.} {\textsc{nb}: \emph{it} is an expletive.}
\end{exe}
These are, in fact, just the verbs that Burzio's Generalization claims are missing.

One way this goal could be achieved using the present system is to rely on the difference between the denotation of verbs that have $\theta$-roles and the denotation of external $\theta$-roles that is offered in~(\ref{ex3:agentfn}). 
\begin{exe}
	\exr{ex3:agentfn}
	\denotes{agent} = $\lambda P.\lambda x.\lambda \epsilon$ \textsc{agent}($x)(\epsilon$) $\wedge$ $P(\epsilon)$, where
	\begin{xlist}
	\ex $P$ is a a function whose domain is eventualities, $x$ is a thing, and $\epsilon$ is an eventuality;
	\ex $P$ is VP.
	\end{xlist}
\end{exe}
External $\theta$-roles, on this view, are functions whose domains are the meanings that VPs have---that is, external $\theta$-roles have VPs as their objects---while internal $\theta$-roles make the meanings of the verbs containing them functions whose domains are things---that is, verbs take DPs as objects. If this is the right denotation for external $\theta$-roles, then Existential Closure won't apply to them because they are not of the right semantic type. Let's adopt this account of how Existential Closure is restricted to internal $\theta$-roles.

Here, then, are desiderata for a good theory of Passive.
\begin{exe}
	\raggedright
	\ex \label{ex3:passivedesi} 
	\begin{xlist}
		\ex The Passive makes the external $\theta$-role susceptible to Existential Closure
		\ex The Passive prevents the external $\theta$-role from assigning accusative Case to the complement of~VP.
	\end{xlist}
\end{exe}
I don't know of an account of the Passive that meets these desiderata. I'll make a proposal here, however, that is not too distant from some of the proposals one can find in the literature. What I propose is that the Passive morpheme attaches to external $\theta$-roles---things of morphosyntactic type ``v''---and changes the way that the external $\theta$-role it attaches to combines with its arguments. I'll express this idea under the assumption that external $\theta$-roles have denotations like that illustrated by~(\ref{ex3:agentfn}). Here's the denotation for the passive morpheme,~\textsc{pass}.
\begin{exe}
	\raggedright
	\ex \label{ex3:pass} \marginnote{So far as I know, there is no proposal in the literature about passive like this one. What I have done here is try to derive a common description of what the passive does in the framework that has been developed in this chapter.}
	\denotes{\textsc{pass}} = $\lambda \mathscr{P}$ $\lambda y$ $\lambda Q$  $\lambda \epsilon$ $\mathscr{P}$($Q(\epsilon))(y)(\epsilon$), where\\
 $\mathscr{P}$ are functions whose domain is a function whose domain is things, $y$ is things, and $Q$ is a function whose domain is things. 
\end{exe}
How this formula applies to the denotations of external $\theta$-roles is somewhat complex, and I won't step through it here. Instead, I'll indicate the change it creates with~(\ref{ex3:passivrule}).
\pagebreak
\begin{exe}
	\raggedright
	\ex \label{ex3:passivrule} Passive (\textsc{pass})\\
	If \denotes{$\theta$-role} = $\lambda P$ $\lambda x$ $\lambda \epsilon$ \textsc{$\theta$-role}($x)(\epsilon$) $\wedge$ $P(\epsilon)$=T, where
	\begin{xlist}
		\sn \textsc{$\theta$-role} c-selects DP
	\end{xlist}
	 \emph{then}\\
	\denotes{\textsc{pass}($\theta$-role)} = $\lambda x$ $\lambda P$ $\lambda \epsilon$ \textsc{$\theta$-role}($x)(\epsilon$) $\wedge$ $P(\epsilon)$=T, where 
	\begin{xlist}
		\sn \textsc{$\theta$-role} c-selects DP
	\end{xlist} 
\end{exe}
\marginnote[-20pt]{From this point forward, I will not include in the denotations what the domain of the functions are. Instead, I will adopt the following conventions. Lower case letters ($x$, $y$,\ldots) will indicate variables over things. Upper case letters ($P$, $Q$,\ldots) will indicate variables over functions whose domain is the set of things or eventualities. And $\epsilon$ will indicate a variable over eventualities. I will also suppress the c-selection information when I think it is irrelevant.}What the Passive does it takes a $\theta$-role whose first argument is a VP and whose second argument is the $\theta$-marked subject, and creates a new $\theta$-role that takes the $\theta$-marked subject argument first and the VP argument second. The $\theta$-roles that have this kind of meaning are just the external $\theta$-roles.

If Passive in general is a process that applies to external $\theta$-roles, then we have an explanation for the fact that Passives do not seem to be able to apply to unaccusative verbs in those languages in which the Passive can apply to intransitive verbs. Passive doesn't seem to be able to apply to intransitive verbs in English, but in many languages this is not the case. If Passive is an operation that applies to external $\theta$-roles, and not the verbs that those $\theta$-roles combine with, this is explained. Unaccusative verbs are just those verbs that do not find themselves in a syntax with external $\theta$-roles. This, then, is a fact which supports the view that the Passive is sensitive to the presence of external $\theta$-roles, and this is built directly into~(\ref{ex3:passivrule}).

The Passive makes an external $\theta$-role become semantically the kind of thing that Existential Closure can apply to. So if we allow Existential Closure to apply to \obar{v} as well as \obar{V}, then we have an explanation for why the Passive allows Existential Closure to apply to external $\theta$-roles. This is shown for the Agent $\theta$-role in~(\ref{ex3:ecext}).

\begin{fullwidth}
\begin{exe}
	\raggedright
	\ex \label{ex3:ecext} 
	\begin{xlist}
		\ex 
		\denotes{agent} = $\lambda P$ $\lambda x$ $\lambda \epsilon$ \textsc{agent}($x$)($\epsilon$) $\wedge$ $P(\epsilon)$=T 
		\ex 
		\denotes{\textsc{pass}(agent)} = $\lambda x$ $\lambda P$ $\lambda \epsilon$ \textsc{agent}($x$)($\epsilon$) $\wedge$ $P(\epsilon)$=T.
		\ex 
		\denotes{\textsc{exist closure}(\textsc{pass}(agent))} = $\lambda P$ $\lambda \epsilon$ $\exists x$ \textsc{agent}($x)(\epsilon$) $\wedge$ $P(\epsilon)$=T.
	\end{xlist}
\end{exe}
\end{fullwidth}

Consider, now, how an external $\theta$-role fits into a D-structure when a passive morpheme has combined with it.
\begin{exe}
	\raggedright
	\ex \label{ex3:ecpasst} 
	\small
	\begin{forest}
	[vP[??\\\ibar{v}
		[$\lambda x$ $\lambda P$ $\lambda \epsilon$ \textsc{agent}($x$)($\epsilon$) $\wedge$ $P(\epsilon)$\\ \obar{v} [\textsc{pass}(agent)] ]
		[$\lambda \epsilon$ \textsc{patient}(\denotes{it})($\epsilon$) $\wedge$ \textsc{melt}($\epsilon$) \\ VP[\ibar{V}
			[\obar{V} [melt] ]
			[DP [it, roof] ]
		] ]
	]]
	\end{forest}
\end{exe}
Our semantic composition rules do not allow \obar{v} and VP to semantically combine -- neither is a function that can take the other as its argument. This could be remedied if we applied Existential Closure to the external $\theta$-role, as in~(\ref{ex3:ecpasstt}).
\begin{exe}
	\raggedright
	\ex \label{ex3:ecpasstt} 
	\small
	\begin{forest}
	[vP [\ibar{v}
		[$\lambda P$ $\lambda \epsilon$ $\exists x$ \textsc{agent}($x$)($\epsilon$) $\wedge$ $P(\epsilon)$\\ \obar{v} [\textsc{pass}(agent)] ]
		[$\lambda \epsilon$ \textsc{patient}(\denotes{it})($\epsilon$) $\wedge$ \textsc{melt}($\epsilon$)\\ VP [\ibar{V}
			[\obar{V} [melt] ]
			[DP [it, roof] ]
		] ]
	]]
	\end{forest}
\end{exe}
But this structure also violates the c-selection restrictions of ``\textsc{pass}(agent).'' The \textsc{agent} $\theta$-role c-selects a DP, in both its active and passive guises. The conditions in Lexical Insertion prevent ``\textsc{pass}(agent)'' from being inserted in the position it occupies in~(\ref{ex3:ecpasst}).

A structure that Lexical Insertion does allow a Existentially closed, passivized $\theta$-role to occupy, however, is~(\ref{ex3:ecpasst2}).
\begin{exe}
	\raggedright
	\ex \label{ex3:ecpasst2} 
	\small
	\begin{forest}
	[vP [??\\\ibar{v}
		[$\lambda P$ $\lambda \epsilon$ $\exists x$  \textsc{agent}($x$)($\epsilon$) $\wedge$ $P(\epsilon)$\\ \ibar{v} [$\lambda P$ $\lambda \epsilon$ $\exists x$  \textsc{agent}($x$)($\epsilon$) $\wedge$ $P(\epsilon)$\\ \obar{v} [\textsc{pass}(agent)]]]
		[$\lambda \epsilon$ \textsc{patient}(\denotes{it})($\epsilon$) $\wedge$ \textsc{melt}($\epsilon$)\\ VP [\ibar{V}
			[\obar{V} [melt] ]
			[DP [it, roof] ]
		]]
	]]
	\end{forest}
\end{exe}
\marginnote[0pt]{It may be useful to return to the discussion of (\ref{ex3:LI}) on page~\pageref{ex3:LI}.} Recall that the c-selection requirements of a morpheme can be satisfied if the \obar{X} position it is inserted in has no sister. This happens, however, only when Existential Closure occurs, and that is required in (\ref{ex3:ecpasst2}) as well for otherwise the highest \ibar{V} will get no denotation. Hence, applying Existential Closure to a passivized external $\theta$-role produces VPs like~(\ref{ex3:imppassf}).
\begin{exe}
	\raggedright
	\ex \label{ex3:imppassf}
	\small
	\begin{forest}
	[$\lambda \epsilon$ \textsc{agent}($x$)($\epsilon$) $\wedge$ \textsc{patient}(\denotes{it})($\epsilon$) $\wedge$ \textsc{melt}($\epsilon$) \\vP[$\lambda \epsilon$ \textsc{agent}($x$)($\epsilon$) $\wedge$ \textsc{patient}(\denotes{it})($\epsilon$) $\wedge$ \textsc{melt}($\epsilon$)\\ \ibar{v}
		[$\lambda P$ $\lambda \epsilon$ $\exists x$ \textsc{agent}($x$)($\epsilon$) $\wedge$ $P(\epsilon)$\\ \ibar{v} [$\lambda P$ $\lambda \epsilon$ $\exists x$  \textsc{agent}($x$)($\epsilon$) $\wedge$ $P(\epsilon)$\\ \obar{v} [\textsc{pass}(agent)]]]
		[$\lambda \epsilon$ \textsc{patient}(\denotes{it})($\epsilon$) $\wedge$ \textsc{melt}($\epsilon$)\\ VP [\ibar{V}
			[\obar{V} [melt] ]
			[DP [it, roof]]
		] ]
	]]
	\end{forest}	
\end{exe}
This is the right outcome.

Note that in (\ref{ex3:imppassf}) \obar{v} does not c-command \emph{it}. This, I suggest, is why there is no accusative Case available to this DP in the Passive. As a consequence, A Movement must apply, causing the object DP to relocate into a Case marked position. 

We have so far only looked at English Passive sentences in which the external argument is not expressed, but it is also possible to express the external argument. This is shown in sentences like~(\ref{ex3:by1}).\marginnote[-30pt]{When a Passive construction includes an expression that contains the external $\theta$-role bearer, as the \emph{by}-phrase does in English, it is called a ``long passive.'' The ones with no expression of the external $\theta$-role bearer are called ``short passives.'' Not all languages have both.}
\begin{exe}
	\raggedright
	\ex \label{ex3:by1} 
	\begin{xlist}
		\ex The ice was melted by Janina.
		\ex We were bothered by the news.
		\ex It was seen by the boy.
		\ex This is entailed by that.
	\end{xlist}
\end{exe}
As this series indicates, the $\theta$-role that the active versions of these verbs sanction on the subject argument are borne by the DP inside the \emph{by}-phrase. This correspondence is completely regular in English: the external $\theta$-role of an active sentence is the same that is assigned to the DP in a \emph{by}-phrase in the corresponding Passive sentence. Chris Collins\cite{Collins:2005a} has suggested that this is because the external $\theta$-role, \obar{v}, is pronounced in the Passive as \emph{by}. Let's adopt this idea. In these sentences, then, the external $\theta$-role does not undergo Existential Closure. Instead it combines with the DP argument that it c-selects. This is illustrated by~(\ref{ex3:by2}).
\begin{exe}
	\raggedright
	\ex \label{ex3:by2} 
	\small
	\begin{forest}
	[$\lambda P$ $\lambda \epsilon$ \textsc{agent}(\denotes{Janina})($\epsilon$)\\ vP [$\lambda P$ $\lambda \epsilon$ \textsc{agent}(\denotes{Janina})($\epsilon$)\\ \ibar{v}
		[$\lambda x$ $\lambda P$ $\lambda \epsilon$ \textsc{agent}($x$)($\epsilon$)\\ \obar{v} [\textsc{pass}(agent)] ]
		[DP [Janina, roof] ]
	]]
	\end{forest}
\end{exe}
In this scenario, \obar{v} cannot embed the VP it combines with---here, the \obar{v} combines with the DP it c-selects. Another way of parsing a sentence with (\ref{ex3:by2}) will be have to be found. One parse that works arises if the vP is put inside the projection of the verb; (\ref{ex3:by3}) illustrates.
\begin{exe}
	\raggedright
	\ex \label{ex3:by3} 
	\hspace*{-40pt}
	\small
	\begin{forest}
	[$\lambda \epsilon$ \textsc{agent}\denotes{Janina}($\epsilon$) $\wedge$ \textsc{patient}(\denotes{it})($\epsilon$) $\wedge$ \textsc{melt}($\epsilon$)\\ VP 
	[$\lambda \epsilon$ \textsc{agent}\denotes{Janina}($\epsilon$) $\wedge$ \textsc{patient}(\denotes{it})($\epsilon$) $\wedge$ \textsc{melt}($\epsilon$)\\ \ibar{V}
		[$\lambda \epsilon$ \textsc{patient}(\denotes{it})($\epsilon$) $\wedge$ \textsc{melt}($\epsilon$)\\ \ibar{V} 
			[$\lambda \epsilon$ \textsc{patient}(\denotes{it})($\epsilon$) $\wedge$ \textsc{melt}($\epsilon$)\\ \ibar{V} 
				[$\lambda x$ $\lambda \epsilon$ \textsc{patient}($x$)($\epsilon$) $\wedge$ \textsc{melt}($\epsilon$)\\ \obar{V} [melt] ]
				[DP [it, roof]]
			]
		]
		[$\lambda P$ $\lambda \epsilon$ \textsc{agent}(\denotes{Janina})($\epsilon$) $\wedge$ $P(\epsilon)$\\ vP [$\lambda P$ $\lambda \epsilon$ \textsc{agent}(\denotes{Janina})($\epsilon$) $\wedge$ $P(\epsilon)$=T \ibar{v}
			[$\lambda x$ $\lambda P$ $\lambda \epsilon$ \textsc{agent}($x$)($\epsilon$) $\wedge$ $P(\epsilon)$=T \\ \obar{v} [\textsc{pass}(agent)] ]
			[DP [Janina, roof] ]
		]]
	]]
	\end{forest}
\end{exe}
In (\ref{ex3:by3}), \obar{v} c-commands \emph{Janina}, and it is to this DP that \obar{v} assigns accusative Case. Here too, then, the DP \emph{it} does not get accusative Case in its D-structure position and it must move to satisfy the Case filter.

This sketch of how the Passive works captures how the external $\theta$-role survives in the passive construction, and meets the desiderata of a theory of Passive laid out in (\ref{ex3:passivedesi}). It leaves unexplained why the Passive \obar{v} is pronounced as \emph{by} when it combines with a DP, and not apparently pronounced at all otherwise. I have nothing insightful to offer for this problem.

With the adoption of a syntax that puts external $\theta$-roles outside the VP their meaning is part of, we capture a variety of facts. With suitable adjustments to the denotations of verbs and $\theta$-roles, we can improve the Projection Principle. It is no longer a strange disjunction of conditions -- one for external $\theta$-roles and another for internal $\theta$-roles. Instead, it is a simple condition that applies to all $\theta$-roles alike. We also have a way of modeling Burzio's Generalization: external $\theta$-roles assign accusative Case, not verbs. And it is possible to design an account of Passive that captures some of its puzzles. 

But putting the external $\theta$-role outside of the VP it is construed with raises problems as well. As noted above, it is inconsistent with our conditions on accusative Case assignment. And it raises the salient puzzle of how to connect the external $\theta$-role with the verb that heads its complement. Which external $\theta$-role is found seems to depend on the selection of the verb. Indeed, whether an external $\theta$-role is found also seems to depend on the choice of verb. By syntactically separating the verb from its external $\theta$-role, we can no longer explain this dependence by assuming that the external $\theta$-role is part of the verb that is inserted in \obar{V} position. If this is the correct direction, we need to find a way for the external $\theta$-role in \obar{v} position to communicate \underline{morphologically} with the verb in \obar{V} position.

We have already encountered a problem of a similar type. When a finite IP is headed by present or past tense, as in~(\ref{ex3:last}), this information must communicate morphologically with the verb inserted in the neighboring \obar{V} position.
\begin{exe}
	\raggedright
	\ex \label{ex3:last} 
	\small
	\begin{forest}
	[IP	
		[DP [she, roof] ]
		[\ibar{I} 
			[\obar{I} [pres] ]
			[VP [\ibar{V} 
				[\obar{V} [be] ]
				[AP [happy, roof] ]
			] ]
		]
	]
	\end{forest}
\end{exe}

We need a theory of the morphology/syntax interface that achieves those goals. This is the job of the next chapter.